% =============================================================================
% KlimaGG-Web — klimagg_snapshot_template.tex
% -----------------------------------------------------------------------------
% Version:        v1.0.14
% Stand:          2025-12-13
% Zweck:          LaTeX-Template für Snapshot-/Release-Export (PDF) aus KlimaGG-Web
% Verwendet von:  export_snapshot.py
% Bearbeitungsstand:
%   - Stabiler Basis-Satz (KOMA-Script), sauberer TOC, Kopfzeilen, neutrale Links
%   - Enthält Transparenz-/Datenschutz-Hinweis (keine personenbezogenen Daten)
%
% Platzhalter (werden von export_snapshot.py ersetzt):
%   2025-12-30         = Datum/Zeitpunkt des Snapshots (z.B. 2025-12-08)
%   Artikel: 71 \quad Kommentare: 0 \quad Votes: 0 \quad Nutzer: 1 = Kennzahlenzeile(n) (z.B. Votes/Kommentare/Artikelstände)
%   %%ARTICLES%%     = vollständige Artikel-Abschnitte (inkl. Formatierungen)
%
% WICHTIG:
%   - Die Platzhalter-Namen (2025-12-30 etc.) nur im Dokumentkörper verwenden.
%   - Keine Textausgabe vor \documentclass (LaTeX würde sonst fehlschlagen).
% =============================================================================

\documentclass[12pt,a4paper]{scrartcl}

% --- Sprach- und Schriftgrundlagen ---
\usepackage[ngerman]{babel}
\usepackage[utf8]{inputenc}   % UTF-8 Eingabe (für breite Kompatibilität)
\usepackage[T1]{fontenc}      % saubere Trennung/ Akzente
\usepackage{lmodern}          % Latin Modern als Standardschrift
\usepackage{microtype}        % bessere Silbentrennung, Laufweite

% --- Layout & Seitenränder ---
\usepackage[a4paper,margin=3cm]{geometry}

% --- Zitate, Listen, Links ---
\usepackage{csquotes}
\usepackage{enumitem}
\setlist[itemize]{leftmargin=2em}
\setlist[enumerate]{leftmargin=2em}
\usepackage[
  colorlinks=false,
  pdfborder={0 0 0},
  pdfauthor={Projekt KlimaGG},
  pdftitle={Klima-Generationen-Gesetz – Stand 2025-12-30}
]{hyperref}

% --- Kopf- und Fußzeilen (KOMA) ---
\usepackage[headsepline,plainheadsepline]{scrlayer-scrpage}
\clearpairofpagestyles
\ihead{Klima-Generationen-Gesetz}
\ohead{Stand: 2025-12-30}
\cfoot{\pagemark}
\pagestyle{scrheadings}

% --- Absatz- und Überschriftenstil ---
\setlength{\parskip}{0.75\baselineskip}
\setlength{\parindent}{0pt}

\setkomafont{section}{\bfseries\large}
\setkomafont{subsection}{\bfseries\normalsize}
\setkomafont{sectionentry}{\normalfont}
\setkomafont{pagenumber}{\normalfont}

% --- Metadaten / Titelseite ---
\title{Klima-Generationen-Gesetz\\\large Stand: 2025-12-30}
\author{Projekt KlimaGG}
\date{}  % Datum im Titel, nicht hier

\begin{document}
\maketitle

\thispagestyle{scrheadings}

\tableofcontents
\newpage

\section*{Hinweis}
Diese Fassung wurde aus der KlimaGG-Web-Datenbank erzeugt. Sie ist
für Transparenz und Diskussion bestimmt, enthält aber \emph{keine}
personenbezogenen Daten. Kommentare und Votes sind nur anonymisiert bzw.
in aggregierter Form berücksichtigt.

Für die rechtlich verbindliche Fassung des Gesetzestextes sind allein
die jeweils maßgeblichen amtlichen Veröffentlichungen entscheidend.

\bigskip
\noindent
\textbf{Kennzahlen zum Erzeugungszeitpunkt:}\\
Artikel: 71 \quad Kommentare: 0 \quad Votes: 0 \quad Nutzer: 1

\bigskip

% --------------------------------------------------------------------
% Gesetzestext (Artikel)
% --------------------------------------------------------------------


%% ----------------------------------------
\section[Grundlegendes]{Grundlegendes}
\noindent\textit{ID: 2 -- grundlegendes (meta)}

\begin{quote}
Wir Deutschen können Regeln. Das ist keine Marotte, sondern ein Kulturpfad: von der Hanse mit ihren Kaufmannsordnungen bis zum Frankfurter Kranz, der nur gelingt, wenn man das Rezept ernst nimmt. Das Zusammenspiel aus geschriebenen und gelebten Regeln hat unsere Wirtschaft groß gemacht. Made in Germany war nie nur ein Etikett, sondern ein Versprechen: Qualität, die aus Standards entsteht --- aus Daten und Prüfungen, aus Stichproben und Verantwortung, aus DIN, TÜV und dem schlichten Satz „Wir stehen dafür ein“. Genau dafür bezahlt der Markt einen Aufpreis. Nicht für Pathos, sondern für Verlässlichkeit. Ein Motor hält, weil er nach Norm gebaut ist. Ein Chemieprozess liefert, weil Messwerte und Grenzwerte nicht verhandelbar sind. Ein Auto verdient Vertrauen, wenn Rückrufe transparent geregelt sind. Qualität ist messbar --- und darum etwas wert. Diesen Grundsatz übertragen wir auf das größte Modernisierungsprojekt unserer Zeit: den Schutz des Klimas und die Erneuerung unserer Infrastruktur. Das Klima-Generationen-Gesetz ist kein Katalog von Verboten, sondern ein Klima-Standard. Es macht Klimaschutz marktfähig, indem es drei einfache Dinge verbindlich regelt.

 Der Klima-Standard:
- Ein Preisschild: Ein verlässlicher Treibhausgas-Preis als Leitplanke, damit Emissionen überall den gleichen, planbaren Gegenwert haben --- heute, morgen und über Legislaturperioden hinaus.

- Einen Leistungsstandard: Wie Klimaleistung gemessen und angerechnet wird. Damit entsteht Qualität statt Greenwashing --- und der Markt zahlt dafür ein Premium.

- Eine Mindestgarantie: Der Staat setzt den Rahmen für Daten, Abwicklung und Haftung. Das sichert Planbarkeit damit Unternehmen und Kommunen investieren können.
 
  Mit diesen Standards wird Klimaschutz von der Kostenstelle zur Wertquelle. Ein Produkt ist künftig nicht nur „sauber“, sondern nachweisbar wirksam --- und wird dafür am Markt vergütet. Das erhöht die Exportchancen deutscher Technik und reduziert die Finanzierungskosten, weil Regeln Risiko ersetzen. Deutschland hat dafür besondere Gründe. Als Mittelland ohne Rohstoffreichtum entsteht unser Vorsprung nicht aus Öl oder Erz, sondern aus Standards, Präzision und Effizienz. Wer keine Rohstoffe verkauft, muss Technologie- und Qualität für sich sprechen lassen. Der Klima-Standard schafft dafür Schnittstellen, Datenregeln und Prüfroutinen, die aus Technologieabhängigkeit Technologieoffenheit machen. So bleiben Märkte offen und Innovationen anschlussfähig --- und Mittelstand wie Kommunen können die Transformation nach Zweckmäßigkeit gestalten. So entsteht ein neues Wirtschaftswunder. Nicht aus Zufall, sondern aus Ordnung. Die Voraussetzungen dafür ist der Generationen-Pakt.

 Der Generationen-Pakt:
- Kreditvolumen: Mit klaren Tilgungswegen für den riesigen Investitionsbedarf in Netze, Schienen, Wärme und Speicher. Das sind volle Auftragsbücher für Handwerk und Maschinenbau.

- Sozialausgleich: damit alle anpacken können. 

- Forschung und Innovation: Für neue Technologie.

- Vorsorge: Die Städte und Regionen gegen Hitze, Hochwasser und Dürre resilient macht.

  Die Kombination aus Klima-Standard und Generationen-Pakt, das bedeutet: Sicherheit für die älteren Jahrgänge, Chancen für die jüngeren --- und Regeln, auf die sich alle verlassen können. So kann privates Kapital, das weltweit nach verlässlichen Renditen sucht, profitabel eingesetzt werden. Dieses Gesetz knüpft an das an, was Deutschland stark gemacht hat: Standards, die Vertrauen schaffen. Es übersetzt dieses Erfolgsrezept ins 21. Jahrhundert. Nicht „mehr Staat“, sondern besserer Staat: ein Schiedsrichter, der das Spielfeld markiert, die Uhr stellt und die Tabelle führt --- und dann die Mannschaften spielen lässt. So wird Klimaschutz kein Kulturkampf, sondern ein Qualitätsstandard, der investierbar, exportierbar und für alle sichtbar ist. Wenn wir das ernst nehmen, entsteht wieder das, wofür dieses Land bekannt ist: solide Arbeit, die hält. Und ein Markt, der das anerkennt --- mit Aufträgen, mit Jobs, mit einem neuen Vertrauen in das kleine Wort auf großen Produkten:

 Made in Germany. Jetzt ergänzt um eine neue Zeile: Klimaleistung - geprüft, angerechnet, fair vergütet/bezahlt.

Ein kurzer Blick auf den Status quo:
 Lisa, Pflegekraft in Gelsenkirchen. Schichtwechsel um 6:00 Uhr. Lisa fährt mit dem Auto -- teurer geworden, seit der Treibhausgas-Preis stieg. Aber der Bus fällt wegen Baustelle an der Brücke wieder aus. In ihrer Mietwohnung hat der Vermieter „energetisch modernisiert“; die Heizung ist effizienter, aber die Warmmiete ist kaum gesunken, eher gestiegen. Sie spürt vor allem Mehrkosten. In der Klinik klagt man über Zugluft und kaputte Fenster -- Sanierung „läuft“. Lisa versteht, warum Klimaschutz nötig ist, aber die Regeln fühlen sich zufällig an: Preise rauf, Netze alt, und wer wenig verdient, zahlt relativ mehr. Soll es wirklich so weitergehen?

Schnellstart: Klimaschutz wird im KlimaGG mit einem verbesserten Treibhausgas-Preis (Art. 7) und einer Vergütung für nachweisbare Klimaleistungen, dem Carbon Reward (Art. 5), umgesetzt. Das lenkt Investitionen und ermöglicht neue Geschäftsmodelle. Soziale Fairness wird durch die Klimadividende (Art. 8) verankert. Damit alle anpacken können, stehen ein Transformationskredit (Art. 3) sowie Beratungsangebote und digitale Verwaltungsführung bereit: die Transformationsagentur (Art. 10). Das Gesetz ist technologieoffen, marktwirtschaftlich, harmonisiert mit nationalem und europäischem Recht und legt den Grundstein für internationale Kooperation zum wirksamen Klimaschutz. Dazu verankert das Gesetz staatliche Verantwortung (Art. 11 ff.) im Bereich Vorsorge und Katastrophenschutz. Spring im Inhaltsverzeichnis direkt zu den Artikeln. Mit den Schaltern oben blendest du Textebenen ein und aus. Begriffe sind im FAQ erläutert.
\end{quote}


%% ----------------------------------------
\section[Präambel]{Präambel}
\noindent\textit{ID: 3 -- praeambel (meta)}

\begin{quote}
Was passiert, wenn ein öffentlicher Transformationskredit, eine verlässliche Vergütung für eingesparte Emissionen und eine lokale Rückerstattung pro Kopf Wirklichkeit werden? Drei kurze Szenen.	
    

Max, Autohändler aus Böblingen. Max liebt Status. Penthouse, Design, Partys, laute Musik. Klima? Für ihn kein Thema. Und doch dreht sich um ihn herum vieles -- leise. Der Hauseigentümer nutzt den Transformationskredit und modernisiert: bessere Hülle, leise Wärmepumpe, smarte Steuerung. Der Strommix wird Jahr für Jahr sauberer, weil Carbon Rewards neue Wind- und Solaranlagen bankfähig machen -- ohne Zuschlagsroulette, mit verlässlichem Mindestpreis. Im Hintergrund greift noch etwas: Divestments. Alte Kohle- und Gaskapazitäten gehen planbar vom Netz, rechtssicher. Das spürt Max nicht direkt -- aber sein Fußabdruck sinkt, weil sein Zuhause, seine Marken und seine Lieferketten weniger Emissionen verursachen. Er behält alle Freiheiten. Das System erledigt den Rest. Auch wer nichts ändert, rutscht mit - und die Freiheit bleibt. CO\$\_2\$e vorher: 9 t/Jahr CO\$\_2\$e nachher: 5--6 t/Jahr (Flüge außen vor)

Lukas, Schweinebauer aus Buxtehude. Schweine, Mais-Monokultur, dünne Margen. Der Hof stand kurz vor der Aufgabe. Dann kam die Transformationsagentur auf den Hof -- nicht mit einem Stapel Formulare, sondern mit Gummistiefeln und einer Checkliste. Gemeinsam wird ein Plan geschnürt: PV auf die Ställe, ein kleiner Speicher fürs Netz, bodenschonende Bewirtschaftung, weniger Mineraldünger und Kompostierung von Reststoffen. Entlang der Gräben werden Hecken gesetzt; zwei neue Reihen Bäume spenden Schatten für die Tiere. Und weil der Bach regelmäßig über die Ufer tritt, schlägt die Kommune vor: „Wir finanzieren Klimaschutz auf Ihrer Fläche“. Der Landwirt duldet die Biber, baut kleine Durchlässe und akzeptiert Überflutungsflächen -- als natürlichen Katastrophenschutz. Die Klimaleistungen laufen: weniger Emissionen aus Boden und Stall, dauerhafter Kohlenstoff im Boden. Der Transformationskredit deckt die Investitionen; die Agentur bleibt ansprechbar, wenn es hakt. Ein Jahr später staunen alle: Der Boden hält Wasser wie ein Schwamm, die Erträge sind stabil, die Tiere haben es kühler -- und der Hof hat neben den klassischen Erlösen einen zweiten, planbaren Geldstrom. Beratung, Finanzierung, Carbon Reward und kommunaler Katastrophenschutz greifen ineinander. Der Hof wird widerstandsfähig -- ökologisch und wirtschaftlich. CO\$\_2\$e vorher: 1 700 t/Jahr CO\$\_2\$e nachher: 900 t/Jahr - und dazu landen 300 t/Jahr als Kohlenstoff im Boden

Familie Kunz aus Magdeburg. Die Kunzes wohnen am Stadtrand. Wärmepumpe und Dämmung sind drin, die PV auf dem Dach brummt seit Jahren. Urlaub? Mit der Bahn. Seit kurzem fährt ein kleines E-Auto. Neu ist: Ihre Klimaleistung wird sichtbar -- und bezahlt. Jede eingesparte Tonne bekommt einen Gutschrift; ein Teil tilgt den alten Sanierungskredit, der Rest ist frei verfügbar. Die Familie beschließt: Gemeinwohl zuerst. Sie nutzen ihre jährliche Dividende und legen nochmal etwas drauf. Damit finanzieren sie kommunalen Klimaschutz: eine Moorfläche am Stadtrand wird wiedervernässt, die Schule erhält ein Holz-Modulgebäude, und die Stadtbibliothek testet ein offenes Forschungsprojekt für „Strom sparen ohne Komfortverlust“. Die Kunzess verfolgen alles in der App, spenden ihre Claim-Rechte an die Gemeinde und bekommen dafür ein „Dankeschön“. Reicht ihnen völlig. Ihr Alltag? Wie vorher -- nur ruhiger im Kopf: Sie zahlen nicht mehr drauf, sondern gestalten mit. Ihre harte Vorarbeit für Deutschlands Zukunft wird belohnt. Die Familie senkt Kosten, unterstützt die Nachbarschaft und beschleunigt lokale Forschung -- ganz ohne neue Bürokratie. CO\$\_2\$e vorher: 2,5 t/Jahr  CO\$\_2\$e nachher: 0,8 t/Jahr (Freiwillige Beiträge zu Gemeindereduktionen nicht gezählt)

    Worum es geht:
- Planbarkeit \& Eigentumsschutz: Verlässliche Regeln, die Innovation ermöglichen, Vermögen nicht entwerten und Haushalten wie Unternehmen Sicherheit geben.

- Bezahlbare Energie \& starker Mittelstand: Ein wirksames, kalkulierbares und sozial ausgestaltetes CO\$\_2\$e-Preissignal entlasten Budgets und stärken Wettbewerbsfähigkeit.

- Strukturwandel als Chance: Neue Einkommensquellen u. a. in Landwirtschaft und Forstwirtschaft durch messbare Klimaleistungen.

- Sozialer Ausgleich: Mehrbelastungen werden pro Kopf zurückgegeben; besonders betroffene Regionen und Berufsgruppen werden gezielt unterstützt.

- Weniger Bürokratie, klare Zuständigkeiten: Eine zuständige Behörde, transparente Daten, digitale Verfahren. Altregelungen werden geordnet abgelöst.

- Europa fest im Blick: EU-Recht, fairer Wettbewerb und offene Märkte sind Leitlinie; nationale Regeln ergänzen europäische Mechanismen sinnvoll.

    Kurz: Das KlimaGG ist ein verlässlicher Ordnungsrahmen, der Umwelt schützt, Freiheit sichert und Wohlstand erneuert.
\end{quote}

- Dieses Gesetz dient dem dauerhaften Schutz der natürlichen Lebensgrundlagen und der Bevölkerung und stärkt Grund- und Freiheitsrechte durch wirksame Vorsorge gegen Klimaschäden.

- Leitprinzipien sind Vorsorge und Verursacherprinzip, Technologieoffenheit, Sozial- und Generationengerechtigkeit, Transparenz, Justiziabilität, Subsidiarität und Wettbewerbsfähigkeit.

- Die Ausgestaltung erfolgt im Einklang mit Unions- und Völkerrecht; kollidierende Vorschriften werden durch Harmonisierungsvorschriften dieses Gesetzes geordnet angepasst.

- Öffentliche Rechenschaft, Open-Data-Veröffentlichung und der Zugang zu gerichtlichem Rechtsschutz sichern die Durchsetzbarkeit der Ziele.

{\footnotesize\begin{quote}
Hover-Hilfe grade für den Einstieg und sperrige Begriffe.
\end{quote}}



%% ----------------------------------------
\section[Artikel 1: Verfassungsänderung für effektiven Klimaschutz]{Artikel 1: Verfassungsänderung für effektiven Klimaschutz}
\noindent\textit{ID: 4 -- artikel1 (gesetz)}

\begin{quote}
Dr. Stern, Richter am Bundesverfassungsgericht. Morgens fährt er mit dem Rad durch den Schlossgarten zum Dienst. Nebel hängt flach über den Wiesen, hinten zeichnen sich die Berge ab. „Früher kamen die Klimasachen als Sturm über uns“, denkt er, „Eilanträge, widersprüchliche Gutachten, jeder wollte sofort Recht bekommen.“ Es fehlte ein Halteseil. In der Mittagspause sitzt Dr. Stern mit einem Kaffee auf den Stufen vor dem Gericht. Er erinnert sich an den ersten Winter nach der Verfassungsänderung. Die Akten wurden nicht dünner -- aber leichter. Plötzlich hatten Entscheidungen eine gemeinsame Richtung: Was schützt, was trägt, was ist den Menschen zumutbar? „Wir mussten nicht mehr das ersetzen, was die Politik offen ließ“, sagt er leise, „wir konnten endlich prüfen statt reparieren.“

 Der erste große Fall nach der Grundgesetzänderung. Die Kläger rügten die Bundesregierung. Plötzlich aber mit klarem Streitgegenstand: eine Abweichung vom Budgetpfad. Mit Prüfmaßstab, Darlegungslast und mit Rechtsfolgen, Kompetenzzuweisungen und Finanzierungsgrundlagen. Keine allgemeine Programmkritik der Kläger mehr, keine bloßen Absichtserklärungen aus den Ministerien. In Karlsruhe konnte man tun, was Karlsruhe soll: Klarheit schaffen, Recht sprechen.

 Das neue Gesetz hatte dem Gericht nicht nur Ziele geliefert, sondern auch Regeln, wie der Staat mit Schulden, Investitionen und der Verantwortung der Länder umgehen muss --- sichtbar gemacht in transparenter Haushaltsführung.

 Abends, zu Hause, liegt ein Brief auf dem Küchentisch. Eine Gemeinde bedankt sich für den schnellen Beschluss. Daneben eine Nachricht seines Neffen: Ausbildungsplatz in einer Firma, die Wärmepumpen einbaut. Dr. Stern lächelt. Es war nie die große Geste, die ihn überzeugt hat. Es war das, was danach nicht mehr passierte: das Gerede, die Vertröstungen, die Lücken. „Diese Änderung hat uns einen Kompass gegeben“, denkt er, „und eine Uhr.“ Und auf einmal blieb genug Zeit für das Wesentliche: entscheiden -- und weiterarbeiten.

Kurz: Die Verfassungsebene verankert die Zielrichtung, die Verbindlichkeit, die Finanzierungsfähigkeit sowie Daten-/Transparenzgrundsätze und effektiven Rechtsschutz. Sie schafft damit einen dauerhaften, justiziablen Rahmen, ohne Detailsteuerung zu versteinern. Mit folgenden Änderungen im Grundgesetz:
- Staatsziel: Klimaschutz wird als Staatsziel mit Bindungs- und Nicht-Rückschrittsprinzip festgeschrieben.

- Budgetbindung: Die Staatlichkeit verpflichtet sich an ein kumuliertes Treibhausgas-Budget und an jährliche Zielpfade.

- Budgetbindung: Die Staatlichkeit verpflichtet sich verfassungsrechtlich, ihr Handeln an einem kumulierten Treibhausgas-Budget und an jährlichen Zielpfaden auszurichten und die Erreichung von Klimaneutralität mit anschließender Netto-Negativ-Phase zu gewährleisten; die einfachgesetzliche Konkretisierung dieser Ziele erfolgt im KlimaGG, insbesondere in Artikel 2a (Klimaziele und Emissionsbudget).

- Bundeskompetenz \& Kohärenz: Der Bund erhält eine ausdrückliche Gesetzgebungskompetenz für Klimaschutz einschließlich Daten- und Registerstandards.

- Finanzverfassung: Spezifische CO\$\_2\$-Steuerkompetenz des Bundes; Transformationskredite werden zielgebunden unter Schutzgeländern der Schuldenregeln verfassungsfest ermöglicht.

- Daten/Transparenz: Verfassungsgrundsätze für offene Klima-Daten, Monitoring und Beteiligung.

- Daten/Transparenz: Verfassungsgrundsätze für offene Klima-Daten, Monitoring, Beteiligung und digitale Register; die einfachgesetzliche Konkretisierung erfolgt insbesondere im KlimaGG durch Artikel 2g (Transparenz, Beteiligung und Ampel-Mechanismus) sowie Artikel 17.14 (Berichts- und Dashboardsystem).

- Backstop \& Notfall: Vorsorgemechanismus für Zielabweichungen und eng gefasste Notfallklausel mit Sunset.

- Backstop \& Notfall: Vorsorgemechanismus für Zielabweichungen und eine eng gefasste Notfallklausel mit Sunset-Klausel, die als äußerer Schutzring oberhalb der einfachgesetzlichen Ampel- und Nachsteuerungsmechanismen des KlimaGG wirkt (insbesondere Artikel 2g und Artikel 17.14), ohne diese zu ersetzen.

- Rechtsschutz: Garantie effektiven und beschleunigten Rechtsschutzes; die Ausgestaltung erfolgt einfachgesetzlich.

Warum auf Verfassungsebene? Damit zentrale Leitplanken stabil, justiziabel und über Legislaturen hinweg gesichert sind. Die Feinsteuerung bleibt einfachgesetzlich und damit anpassbar.

Was bleibt dem einfachen Gesetz vorbehalten? Konkrete Preis-, Förder- und Registermechanik, MRV-Standards, Programmvolumina, Behördenzuschnitte, Verfahrensfristen, Sanktionskataloge und EU-Beihilfe-Umsetzung -- alles präzise, aber unterhalb des Grundgesetzes.

Was bleibt dem einfachen Gesetz vorbehalten? Konkretisierung, Instrumentenwahl und operative Steuerung bleiben einfachgesetzlich. Das Klima-Generationen-Gesetz (KlimaGG) entfaltet diese Verfassungsleitplanken, indem es in Artikel 2a bis 2g die Ziele, Pfade, digitalen Grundlagen und Governance-Mechanismen einfachgesetzlich ausgestaltet und in den Artikeln 3 ff. konkrete Instrumente regelt.

Hinweis zur Einordnung: Die folgende Verfassungsänderung verankert ausschließlich solche Leitlinien, die unmittelbare, dauerhafte grundrechtliche oder finanzverfassungsrechtliche Wirkung entfalten müssen. Operative und technische Regelungen (insbesondere MRV-Templates, Plattform-Spezifikationen, Ampel-Parameter, Trust-Provider-Standards, unterstützende Assistenzmaßnahmen) verbleiben im einfachen Gesetz und in den zugehörigen Verordnungen, um Anpassungsfähigkeit und technologische Entwicklung nicht verfassungsrechtlich zu fixieren.

Hinweis zur Einordnung: Die folgende Verfassungsänderung skizziert Leitplanken und Mindestinhalte, ohne die einfache Gesetzgebung zu versteinern. Im Klima-Generationen-Gesetz (KlimaGG) beginnt die Ausbuchstabierung dieser Leitplanken mit Artikel 2 (Zielsetzung und Leitplanken) und Artikel 2a bis 2g (Ziele, Budget, digitale Grundlagen, Governance) und führt über die Fachartikel 3 ff. zu konkreten Instrumenten und Verfahren.
\end{quote}

Juristischen Volltext einblenden (Die einzelnen Änderungen auf Artikelebene)

- Staatsziel \& Nicht-Rückschritt -- Art. 20a GG (neu gefasst):
- „Der Staat schützt auch in Verantwortung für die künftigen Generationen die natürlichen Lebensgrundlagen und das Klima. Gesetzgebung, vollziehende Gewalt und Rechtsprechung sind verpflichtet, die Einhaltung völker- und unionsrechtlicher Klimaziele sowie eines mit diesen Zielen vereinbaren nationalen Treibhausgasbudgets sicherzustellen und die nationalen Klimaziele, einschließlich der Erreichung der Treibhausgas-Neutralität und anschließender netto-negativer Emissionen, zu erreichen.“

- „Eine Absenkung eines erreichten Schutzniveaus ist nur zulässig durch Maßnahmen gleicher oder höherer Zielwirksamkeit, unter Wahrung der Verhältnismäßigkeit und des gesetzgeberischen Einschätzungsspielraums sowie soweit nicht überwiegende verfassungsrechtliche Güter zwingend entgegenstehen.“

- Auftrag und Mindestinhalte einer Bundesgesetzgebung -- Art. 20c GG (neu):
- „Der Bund regelt durch Bundesgesetz die grundlegenden Mechanismen des Klimaschutzes. Dieses Gesetz stellt mindestens sicher: (1) Zielpfade und Emissionsbudgets, (2) die Mechanik der Emissionsbepreisung einschließlich eines Mindestpreis-Korridors, (3) die Pro-Kopf-Rückverteilung (Klimadividende), (4) ein marktbasiertes System zusätzlicher Klimaleistungsvergütung (Carbon Rewards) mit Register, (5) Risikovorsorge/Backstop-Mechanismen, (6) Grundsätze kommunaler Klima- und Resilienzmittel, (7) Transparenz-, Monitoring- und Nachsteuerungsmechanismen auf Basis digitaler Daten- und Registerinfrastrukturen. Die Regelungen können in einem oder mehreren Bundesgesetzen erfolgen; die Bezeichnung ist unerheblich.“

- „Abweichungen im Sinne einer Reduktion oder Aufhebung der in Absatz 1 genannten Mindestinhalte bedürfen eines Bundesgesetzes, das die Voraussetzungen des Artikels 79 Absatz 2 erfüllt; der Zustimmung des Bundesrates bedarf es.“

- Gesetzgebungskompetenz -- Art. 74 Abs. 1 GG (ergänzt) i. V. m. Art. 72 Abs. 2 GG:
- „Nach Nummer 32 wird eingefügt: ‚32. Klimaschutz einschließlich Emissionsbepreisung, Emissionshandel, Kohlenstoff-Entnahme und -Speicherung, Daten- und Registerinfrastrukturen sowie die Verteilung klimaschutzbedingter Einnahmen an die Bevölkerung.‘“

- „Artikel 72 Absatz 2 gilt; im Bereich des Klimaschutzes ist regelmäßig die Wahrung der Rechts- und Wirtschaftseinheit oder gleichwertiger Lebensverhältnisse berührt.“

- Steuerkompetenz \& Verwendung -- Art. 105 Abs. 2 GG (neu gefasst) \& Art. 106 GG (ergänzt):
- „Artikel 105 Absatz 2 Satz 1 (Ergänzung): Der Bund besitzt die konkurrierende Gesetzgebung für Steuern auf die Emission von Treibhausgasen.“

- „Artikel 106 Absatz 7 (neu): Das Aufkommen aus Steuern auf die Emission von Treibhausgasen steht dem Bund zu. Es ist überwiegend für Maßnahmen des Klimaschutzes und für die Auszahlung einer Pro-Kopf-Rückverteilung (Klimadividende) zu verwenden; Näheres regelt ein Bundesgesetz.“

- Klima-Investitionsfenster in der Schuldenbremse -- Art. 109 und 115 GG (ergänzt):
- „Artikel 109 Absatz 3 (Ergänzung): Abweichend von Satz 1--5 sind kreditfinanzierte Ausgaben zulässig, soweit sie investiven Klimaschutzmaßnahmen dienen, die nach Maßgabe eines Bundesgesetzes MRV-basiert (Messung, Berichterstattung und Verifizierung) eine dauerhafte Emissionsminderung oder eine nachhaltige Resilienzerhöhung bewirken (Klima-Investitionsfenster). Voraussetzung sind: (1) Zielkompatibilität nach Gesetz, (2) Transparenz- und Berichtspflichten, (3) Evaluations- und Korrekturmechanismen, (4) Begrenzung durch mehrjährige Ausgabenpfade und (5) Vereinbarkeit mit Unionsrecht.“

- „Artikel 115 Absatz 2 (Ergänzung): Für das Klima-Investitionsfenster legt ein Bundesgesetz insbesondere die Definition der Investivität, den MRV-Nachweis, Tilgungsregeln, eine Rücklage und die ESA-Konformität fest. Die Kreditaufnahme ist auf Maßnahmen mit dauerhaftem Emissionsminderungseffekt oder resilienter Infrastruktur auszurichten.“

- Klimasenat beim Bundesverfassungsgericht -- Art. 94 GG (ergänzt):
- „Beim Bundesverfassungsgericht wird ein zusätzlicher Senat (Klimasenat) gebildet. Er entscheidet über abstrakte Normenkontrollverfahren, konkrete Normenkontrollen und Organstreitigkeiten mit maßgeblichem Bezug zu Artikel 20a. Die Zuweisung weiterer Verfahren, insbesondere Verfassungsbeschwerden, und das Nähere regelt ein Bundesgesetz (BVerfGG).“

- Hinweis: Alternativ kann ein eigenständiges Gericht durch einen neuen Artikel 92a GG geschaffen werden; hierfür bedarf es einer institutionellen Neuregelung im Bundesverfassungsgerichts-Gesetz.

- Ewigkeitsgarantie -- Art. 79 Abs. 3 GG (Klarstellung, unverändert im Wortlaut):
- „Die Ewigkeitsgarantie bleibt unberührt; die vorstehenden Änderungen sind mit den Artikeln 1 und 20 in Einklang auszugestalten.“

{\footnotesize\begin{quote}
- Systematik: Verfassungsrang für Bindung/Nicht-Rückschritt, Kompetenz, Finanzierungsspielraum und Rechtsschutz. Inhalte/Mechanik folgen einfachgesetzlich; Detailfragen verbleiben im Ausführungsrecht.

- Schuldenbremse: Das Klima-Investitionsfenster schafft klar geregelte Ausnahmen mit Pfaden, Tilgung und Transparenz statt pauschaler Aufweichung.

- Gericht: Der Klimasenat beim BVerfG fokussiert Grundsatz- und Normenkontrollfragen; Massenverfahren werden über das BVerfGG gesteuert.

- Änderungsschutz: Die Mindestinhalte der Bundesgesetzgebung erhalten einen qualifizierten Schutz: Abweichungen nur mit 2/3-Mehrheit (entspricht GG-Änderung) und Bundesratszustimmung.

- Bundeskompetenz (eng, funktional begrenzt): Der Bund erhält eine ausdrücklich auf das Ziel der klimaneutralen Transformation bezogene Gesetzgebungskompetenz, die auch die Einrichtung von nationalen Register- und Datenstandards umfasst, soweit dies zur Wahrung der Rechts- und Wirtschaftseinheit erforderlich ist; detaillierte Betriebsregelungen blei ben der einfachen Gesetzgebung und den Verordnungen vorbehalten.

- Klima-Investitionsfenster: Kreditfinanzierte Investitionen mit dauerhaftem Emissionsminderungseffekt können verfassungsrechtlich gerechtfertigt werden, soweit sie an Mehrjahres-Tilgungspläne, transparente Rücklagenbildung und parlamentarische Kontrollmechanismen gebunden sind.
\end{quote}}



%% ----------------------------------------
\section[Artikel 2: Zielsetzung und Leitplanken]{Artikel 2: Zielsetzung und Leitplanken}
\noindent\textit{ID: 5 -- artikel2 (gesetz)}

\begin{quote}
Worum es in diesem Gesetz am Ende geht, lässt sich in drei Fragen zusammenfassen: Was ist das klimabezogene Ziel Deutschlands, womit wird gemessen, ob wir auf Kurs sind, und was passiert, wenn wir vom Kurs abweichen? Artikel 2 beantwortet diese drei Fragen und macht daraus eine verbindliche Zielsetzung für alle folgenden Artikel.

Das KlimaGG soll sicherstellen, dass Deutschland sein wissenschaftlich abgeleitetes Treibhausgasbudget einhält, spätestens klimaneutral und danach netto-negativ wird, Klimafolgen vorsorgend begrenzt und zugleich eine sozial verträgliche, wirtschaftlich tragfähige Transformation gelingt. Die Ziele gelten nicht abstrakt, sondern werden in messbare Budgets, Pfade und Indikatoren übersetzt, die sich mit den Mitteln moderner Daten- und Registerinfrastruktur überwachen und nachsteuern lassen.

Die Unterartikel 2a bis 2g konkretisieren diese Zielsetzung: Art. 2a legt Endziel, Zwischenziele, Budgetlogik und Versionierung fest; Art. 2b bis 2f beschreiben die dafür notwendige digitale Infrastruktur, die MRV-Integrität, den beihilferechtlich und haushalterisch abgesicherten Finanzierungsrahmen sowie die Transformationspfade für Sektoren; Art. 2g verankert Transparenz, Beteiligung und ein Ampel-Frühwarnsystem. Gemeinsam bilden sie den roten Faden, an dem sich alle Instrumente des KlimaGG auszurichten haben.
\end{quote}

- Gegenstand. Dieses Kapitel legt die Zielsetzungen und Leitprinzipien des KlimaGG fest. Es bestimmt, dass das Gesetz
- die Einhaltung eines mit Art. 1 vereinbaren Treibhausgasbudgets,

- die Erreichung von Klimaneutralität und den Übergang zu Netto-Negativ-Emissionen,

- vorsorgende Klimaanpassung und Resilienz,

- soziale Fairness und Generationengerechtigkeit sowie

- makroökonomische Stabilität und geordnete Haushaltsführung

    gleichzeitig verfolgt und hierzu ein einheitliches, datenbasiertes Steuerungssystem etabliert.

- Dreifache Zielbestimmung. Das KlimaGG
- legt in Art. 2a Endziel, Zwischenziele und Emissionsbudget Deutschlands fest,

- ordnet den Fachartikeln 3 bis 14 und Art. 17 die Rolle zu, diese Ziele technologieoffen, regelgebunden und überprüfbar umzusetzen, und

- schafft mit Art. 2g und Art. 17.14 ein verbindliches System für Monitoring, Berichterstattung, Ampel-Indikatoren und Nachsteuerung.

    Diese dreifache Zielbestimmung bildet den Rahmen für alle weiteren Normen dieses Gesetzes.

- Verhältnis zu Art. 1 und Art. 17. Die in Art. 2a bis 2g geregelten Ziele und Leitprinzipien konkretisieren die in Art. 1 verfassungsrechtlich verankerte Zielrichtung (Staatsziel Klimaschutz, Budgetbindung, Daten- und Transparenzgrundsätze) auf einfachgesetzlicher Ebene. Die Fachartikel des KlimaGG, insbesondere Art. 3 bis 14 und Art. 17, sind so auszugestalten, dass sie diese Zielsetzung unterstützen und mit den Harmonisierungsvorschriften des Art. 17 kohärent zusammenspielen.

- Auslegungs- und Vorrangregel. Bei der Auslegung und Anwendung des KlimaGG, seiner Verordnungen und Verwaltungsakte sind die Zielsetzungen dieses Artikels vorrangig zu berücksichtigen. Insbesondere sind
- Digitalisierung und Verwaltungsmodernisierung (Art. 2b),

- MRV-Integrität und Stufung (Art. 2c),

- beihilferechtliche Einbettung (Art. 2d) und

- haushalts- und finanzverfassungsrechtliche Leitplanken (Art. 2e)

    als Ermöglichungsbedingungen der Klimaziele nach Art. 2a zu verstehen und nicht als nachträgliche Vetoinstrumente, die ambitionierte Zielpfade ohne sachliche Rechtfertigung aushebeln dürfen. Konflikte zwischen Vorschriften sind, soweit rechtlich möglich, zielkonform im Sinne dieses Artikels aufzulösen.

- Europäische Einbettung. Die Zielsetzung des KlimaGG ist im Lichte des europäischen Klimarechts, insbesondere des europäischen Klimagesetzes, des Emissionshandelssystems, der beihilferechtlichen Vorgaben und der einschlägigen Digital- und Datenrechtsakte (einschließlich eIDAS sowie relevanter Daten- und Produktregime), auszulegen. Die in Art. 2a bis 2g vorgesehenen Mechanismen sollen mit europäischen Instrumenten kompatibel sein und können als Modell für deren Weiterentwicklung dienen.

- Verordnungsermächtigung. Die Bundesregierung wird ermächtigt, durch Rechtsverordnung mit Zustimmung des Bundestages nähere Bestimmungen zu den in Art. 2a bis 2g genannten Zielen, Indikatoren und Verfahren zu erlassen, sie regelmäßig zu überprüfen und an den Stand von Wissenschaft, Technik und europäischem Recht anzupassen, soweit der Gesetzgeber nicht binnen der dort genannten Fristen anders entscheidet. Dabei sind die in diesem Artikel niedergelegten Zielsetzungen und die Ambitionsneutralität zu wahren.

{\footnotesize\begin{quote}
Ziel: Artikel 2 macht die Zielrichtung des KlimaGG transparent, verbindet sie mit den verfassungsrechtlichen Leitplanken aus Artikel 1 und erklärt ausdrücklich, dass Digitalisierung, Beihilferecht, Haushaltsregeln und sektorale Detailnormen dem Erreichen der Klimaziele dienen. Artikel 2a bis 2g werden damit zur gemeinsamen Auslegungsgrundlage für alle weiteren Artikel und Verordnungen.
\end{quote}}



%% ----------------------------------------
\section[Artikel 2a: Klimaziele, Emissionsbudget und Versionierung]{Artikel 2a: Klimaziele, Emissionsbudget und Versionierung}
\noindent\textit{ID: 6 -- artikel2a (meta)}

\begin{quote}
Was ist das Ziel und womit messen wir es? Artikel 2a übersetzt die verfassungsrechtliche Zielrichtung aus Artikel 1 in konkrete, messbare Größen: ein Endziel (Klimaneutralität mit anschließender Netto-Negativ-Phase), verbindliche Zwischenziele, ein kumuliertes Emissionsbudget und jährliche Obergrenzen. Diese Größen sind die gemeinsame Referenz für alle Fachartikel des Gesetzes.

Die Zahlen selbst werden nicht in Stein gemeißelt, sondern mit einem klaren Verfahren versioniert: Anpassungen an neue wissenschaftliche Erkenntnisse dürfen nur in Richtung höherer Ambition erfolgen oder müssen zumindest ambitionneutral sein. Jede Veränderung wird offen dokumentiert und maschinenlesbar veröffentlicht.
\end{quote}

- Endziel. Deutschland erreicht spätestens bis zum Jahr 2045 Treibhausgas-Neutralität im Sinne eines ausgeglichenen Netto-Saldo aller vom Menschen verursachten Emissionen und dauerhaften Entnahmen aus der Atmosphäre. In der Folge strebt Deutschland eine Netto-Negativ-Position an, bei der dauerhaft mehr Treibhausgase entnommen als emittiert werden.

- Zwischenziele.
- Für das Jahr 2030 wird eine Reduktion der Netto-Treibhausgasemissionen um mindestens 70 \% gegenüber 1990 angestrebt.

- Weitere Zwischenziele (insbesondere für die Jahre 2035 und 2040) werden durch Rechtsverordnung nach Satz 3 präzisiert; die Verordnung legt Emissionsmengen oder prozentuale Reduktionsziele fest, die mit dem Endziel nach Nummer 1 vereinbar sind.

- Die Klimakommission kann ergänzend zu CO\$\_2\$e/GWP100 gas- oder sektorenspezifische Indikatoren (insbesondere für kurzlebige Klimagase wie Methan) zur Steuerung und Bewertung entwickeln und empfehlen.

- Die Bundesregierung wird ermächtigt, durch Rechtsverordnung (VO 2a) die genauen Zielwerte und Zeitpunkte festzulegen, zu aktualisieren und in einer amtlichen, maschinenlesbaren Tabelle zu führen.

- Nationales Emissionsbudget und jährliche Obergrenzen.
- Für den Zeitraum bis zum Erreichen der Treibhausgas-Neutralität wird ein kumuliertes nationales Emissionsbudget festgelegt, das mit den Verpflichtungen der Bundesrepublik Deutschland im Rahmen des europäischen und internationalen Klimarechts vereinbar ist.

- Die Klimakommission hat das deutsche Budget aus IPCC-AR6-Carbon-Budgets und anerkannten Fair-Share-Ansätzen herzuleiten, regelmäßig zu aktualisieren und transparent zu begründen.

- Das Budget wird in jährliche Emissionsobergrenzen übersetzt, die in VO 2a tabellarisch, maschinenlesbar und sektorenübergreifend festgelegt werden.

- Änderungen des Budgets oder der jährlichen Obergrenzen sind nur zulässig, wenn sie mit dem 1,5-Grad-Ziel vereinbar sind und die Ambitionsneutralität gewahrt bleibt.

- Verteilung auf Sektoren und Pfade. Die Zuordnung der nationalen Zielgrößen zu Sektoren, Instrumenten und Pfaden erfolgt durch Fachrecht und Verordnungen, insbesondere durch VO 2a und die Artikel 3 bis 14. Dabei sind Technologieoffenheit, Effizienz, soziale Fairness und die Vermeidung von Carbon-Leakage zu beachten.

- Versionierung und Ambitionsneutralität.
- Alle Zielgrößen nach diesem Artikel sind mit einer eindeutigen Versionskennung zu versehen und in einem öffentlich zugänglichen, maschinenlesbaren Register zu dokumentieren.

- Anpassungen von Zielwerten und Budgets dürfen das Ambitionsniveau nicht absenken, es sei denn, zwingende völker- oder unionsrechtliche Vorgaben machen eine Anpassung erforderlich; in diesem Fall ist die Abweichung zu begründen und parlamentarisch zu beraten.

- Jede Änderung ist mit einem Änderungsjournal zu veröffentlichen, das Zeitpunkt, Begründung und veranlassende Datenbasis ausweist.

- MRV-Anbindung. Alle Zielgrößen nach diesem Artikel sind mit klaren Kriterien für Monitoring, Berichterstattung und Verifizierung (MRV) zu hinterlegen. VO 2a legt fest, welche Datenquellen, Methoden und Aggregationsstufen für die Zielverfolgung verbindlich sind und wie sie mit den MRV-Verordnungen anderer Kapitel (insbesondere LMRV, CMRV, KMRV) verzahnt werden.

{\footnotesize\begin{quote}
Artikel 2a übersetzt das Staatsziel aus Artikel 1 in konkrete Kennzahlen und sorgt mit Versionierung und Ambitionsneutralität dafür, dass Anpassungen an neue Erkenntnisse möglich sind, ohne das Grundniveau der Ambition zu unterlaufen.
\end{quote}}



%% ----------------------------------------
\section[Artikel 2b: Digitale Infrastruktur, Register und Once-Only-Prinzip]{Artikel 2b: Digitale Infrastruktur, Register und Once-Only-Prinzip}
\noindent\textit{ID: 7 -- artikel2b (meta)}

\begin{quote}
Ohne digitale Grundlagen bleibt jedes Klimagesetz Stückwerk. Artikel 2b legt fest, dass das KlimaGG auf einer durchgängig digitalen Infrastruktur beruht: mit elektronischer Aktenführung, interoperablen Registern, sicheren Identitäten und offenen Schnittstellen.

Ziel ist, dass Daten nur einmal erhoben und dann vielfach genutzt werden (Once-Only-Prinzip), dass Entscheidungen nachvollziehbar und prüfbar sind und dass Verwaltung, Wirtschaft und Bürgerinnen und Bürger auf verlässliche digitale Werkzeuge zugreifen können.
\end{quote}

- Digitale Grundprinzipien. Das KlimaGG folgt den Grundsätzen
- digital-by-design (digitale Ausgestaltung von Verfahren von Anfang an),

- once-only (Daten werden nach Möglichkeit nur einmal erhoben und mehrfach genutzt),

- security-by-design und privacy-by-design sowie

- interoperablen, offenen Standards für Schnittstellen und Datenmodelle.

- Elektronische Akte und Register.
- Vorgänge und Entscheidungen im Anwendungsbereich des KlimaGG sind grundsätzlich in elektronischen Akten zu führen.

- Zentrale oder verteilte Register nach Art. 17 erfassen wesentliche klimabezogene Rechte, Pflichten, Zertifikate, Förderungen und relevanten Datenzustände in nachvollziehbarer Form.

- Die Register sind so auszugestalten, dass sie sowohl den Anforderungen des Datenschutzes als auch denen der Nachvollziehbarkeit und Beweisführung genügen.

- Identitäten, Signaturen und Wallets. Die Identifizierung von Personen und Organisationen sowie die Signatur klima-relevanter Vorgänge erfolgt vorrangig mit Mitteln, die mit eIDAS und dem europäischen digitalen Identitätsrahmen (einschließlich Wallet-Lösungen) kompatibel sind. Nähere Bestimmungen werden in Art. 17 und den darauf beruhenden Verordnungen getroffen.

- Once-Only-Pflichten. Behörden und andere öffentliche Stellen, die im Rahmen des KlimaGG handeln, sind verpflichtet, das Once-Only-Prinzip umzusetzen, soweit dem keine zwingenden rechtlichen oder technischen Gründe entgegenstehen. VO 2b legt Übergangsfristen, Prioritäten und technische Mindestanforderungen fest.

- Schnittstellen und Datenräume. Die für das KlimaGG maßgeblichen Register und Fachverfahren sind über standardisierte Schnittstellen in Datenräumen zu verknüpfen, die sichere, kontrollierte Datennutzung sektorenübergreifend ermöglichen.

{\footnotesize\begin{quote}
Artikel 2b legt die digitale Basis, auf der die späteren Fachregeln laufen. Er ersetzt kein Fachrecht, sorgt aber dafür, dass Klimapolitik nicht an analogen Verfahrensgrenzen scheitert.
\end{quote}}



%% ----------------------------------------
\section[Artikel 2c: MRV-Integrität und Stufung]{Artikel 2c: MRV-Integrität und Stufung}
\noindent\textit{ID: 8 -- artikel2c (meta)}

\begin{quote}
Ohne verlässliche Daten gibt es keine glaubwürdigen Carbon Rewards. Artikel 2c legt die Grundsätze für Monitoring, Berichterstattung und Verifizierung (MRV) fest. Er verbindet hohe Integritätsanforderungen mit pragmatischer Staffelung: nicht jedes Kleinprojekt braucht eine Vollprüfung, aber große Volumina müssen besonders streng geprüft werden.

Die Details werden in sektorspezifischen MRV-Verordnungen geregelt (LMRV, CMRV, KMRV, AMRV usw.). Artikel 2c sorgt dafür, dass diese Verordnungen einem gemeinsamen Integritätsrahmen folgen.
\end{quote}

- Integritätsprinzip. MRV-Regelungen im Anwendungsbereich des KlimaGG müssen gewährleisten, dass
- Emissionen, Einsparungen und Entnahmen nachvollziehbar und prüfbar sind,

- Doppelzählungen und Scheineinsparungen vermieden werden und

- Reversals (Rückgänge von Einsparungen oder Entnahmen) erfasst und adressiert werden.

- Stufung nach Volumen und Risiko.
- MRV-Anforderungen werden nach Volumen, Komplexität und Risikoprofil der Maßnahmen gestuft. Kleinvolumige, standardisierte Maßnahmen können vereinfachten Pfaden folgen, während großvolumige oder risikoreiche Vorhaben strengere Prüfungen erfordern.

- VO 2c und die sektorspezifischen MRV-Verordnungen definieren Mindestanforderungen und Stufenmodelle (z. B. AL-1 bis AL-3), einschließlich Datenanforderungen, Prüfintervallen und Assurance-Leveln.

- Cluster- und Methodenräte. Für komplexe Sektoren (insbesondere Industriecluster, Chemie, Netzinfrastruktur) können anerkannte Methodenräte und Cluster-MRV-Strukturen geschaffen werden, die MRV-Methoden entwickeln, erproben und standardisieren. Kriterien für Anerkennung, Aufgaben und Aufsicht werden in VO 2c festgelegt.

- Offenheit und Reproduzierbarkeit. Soweit möglich, sind MRV-Methoden, Datenmodelle und Standardpfade öffentlich zugänglich zu machen, um Nachvollziehbarkeit und Reproduzierbarkeit zu fördern und Missbrauch vorzubeugen.

- Verzahnung mit Budget und Ampel. MRV-Daten sind so aufzubereiten, dass sie unmittelbar in die Budgetüberwachung nach Art. 2a und die Ampelindikatoren nach Art. 2g einfließen können.

{\footnotesize\begin{quote}
Artikel 2c stellt sicher, dass Carbon Rewards und andere Instrumente auf einem belastbaren Datenfundament stehen, ohne kleine Akteure durch überzogene Anforderungen auszuschließen.
\end{quote}}



%% ----------------------------------------
\section[Artikel 2d: Beihilferechtliche Einbettung und EU-Governance]{Artikel 2d: Beihilferechtliche Einbettung und EU-Governance}
\noindent\textit{ID: 9 -- artikel2d (meta)}

\begin{quote}
Carbon Rewards und andere Instrumente müssen europarechtsfest sein. Artikel 2d legt Grundsätze fest, nach denen das KlimaGG mit dem europäischen Beihilferecht und anderen EU-Vorgaben in Einklang gebracht wird.

Ziel ist, Doppelstrukturen und Blockaden zu vermeiden: Wo möglich, sollen Standardkorridore unter Nutzung der allgemeinen Gruppenfreistellungsverordnung (GBER) und anderer EU-Instrumente geschaffen werden, statt jedes Projekt einzeln notifizieren zu müssen.
\end{quote}

- Beihilfekonformität. Instrumente des KlimaGG, die staatliche Beihilfen im Sinne des Unionsrechts darstellen, sind so zu konzipieren, dass sie mit Art. 107 ff. AEUV vereinbar sind.

- Standardkorridore.
- Für typische Maßnahmen sollen, soweit möglich, standardisierte Beihilfekorridore geschaffen werden, die unter bestehende Freistellungstatbestände fallen (insbesondere GBER).

- VO 2d legt Parameterbereiche (z. B. Förderhöhen, Intensitäten, Kumulationsregeln) fest, die ohne Einzelnotifizierung genutzt werden können.

- EU-Gate. Für Instrumente und Programme, die eine individuelle Notifizierung erfordern, richtet die Bundesregierung ein koordiniertes „EU-Gate“ ein, das Verfahren bündelt, beschleunigt und transparent macht. Die Rolle des EU-Gates wird in VO 2d näher bestimmt.

- Transparenz und Berichterstattung. Beihilferelevante Maßnahmen nach diesem Gesetz sind in einem zentralen, öffentlich einsehbaren Verzeichnis zu dokumentieren, das den EU-Transparenzanforderungen genügt.

- Zusammenwirken mit EU-Instrumenten. Das KlimaGG ist so auszulegen, dass es mit ETS, CBAM, dem europäischen Klimagesetz und sonstigen relevanten EU-Instrumenten kohärent zusammenspielt und Doppelzählungen vermeidet.

{\footnotesize\begin{quote}
Artikel 2d sorgt dafür, dass die beabsichtigten Anreize und Finanzierungsströme nicht an Brüssel scheitern, sondern in geordneten, transparanten Bahnen verlaufen.
\end{quote}}



%% ----------------------------------------
\section[Artikel 2e: Haushalts- und Finanzrahmen für Klimainstrumente]{Artikel 2e: Haushalts- und Finanzrahmen für Klimainstrumente}
\noindent\textit{ID: 10 -- artikel2e (meta)}

\begin{quote}
Klimapolitik braucht einen klaren Finanzrahmen. Artikel 2e legt fest, wie Verpflichtungen aus Carbon Rewards und anderen Instrumenten haushaltsrechtlich einzuordnen sind und wie sie mit der Finanzverfassung (einschließlich Schuldenregeln) zusammenspielen.

Ziel ist, Klimarisiken finanziell sichtbar, planbar und steuerbar zu machen, statt sie zu verdrängen. Carbon Rewards werden als konditionale Verpflichtungen mit Caps, Puffern und Tilgungsregeln behandelt, nicht als unkontrollierbare Ausgabenversprechen.
\end{quote}

- Einordnung als Eventualverbindlichkeit. Verbindlichkeiten aus Carbon-Reward-Versprechen gelten haushaltsrechtlich als konditionale Eventualverbindlichkeiten, die an klar definierte Klimaleistungen geknüpft sind. Nähere Bestimmungen enthalten insbesondere Art. 17.2 und Art. 17.4.

- Caps und Floors.
- Für die Gesamtverpflichtungen können Ober- und Untergrenzen (Caps und Floors) festgelegt werden, die Sicherheit für Haushaltsplanung und Investitionen schaffen.

- VO 2e und Art. 17.2 legen Parameter, Anpassungsregeln und Transparenzanforderungen fest.

- Rücklagen und Tilgung. Es sind angemessene Rücklagen- und Tilgungsmechanismen zu schaffen, die gewährleisten, dass aus Carbon-Reward-Verpflichtungen keine unkontrollierbaren Belastungen zukünftiger Haushalte entstehen.

- Vereinbarkeit mit der Schuldenverfassung. Die Ausgestaltung der Finanzinstrumente nach diesem Gesetz hat den verfassungsrechtlichen Rahmen der Schuldenregeln zu beachten. Soweit das Grundgesetz besondere Klimainvestitionsfenster vorsieht (Art. 1), sind diese in den einfachgesetzlichen Regelungen zu konkretisieren.

- Transparenzberichte. Die Bundesregierung berichtet regelmäßig über Stand, Risiken und Puffer der klimabezogenen Finanzverpflichtungen, insbesondere im Rahmen der Berichterstattung nach Art. 17.14.

{\footnotesize\begin{quote}
Artikel 2e macht klar, dass Klimaverpflichtungen nicht hinter dem Haushalt versteckt werden, sondern sichtbar und steuerbar sind.
\end{quote}}



%% ----------------------------------------
\section[Artikel 2f: Sektorale Transformationspfade und fossile Übergänge]{Artikel 2f: Sektorale Transformationspfade und fossile Übergänge}
\noindent\textit{ID: 11 -- artikel2f (meta)}

\begin{quote}
Der Weg in die Klimaneutralität verläuft durch konkrete Sektoren. Artikel 2f beschreibt Grundsätze für Transformationspfade in Energie, Industrie, Verkehr, Gebäude, Landwirtschaft und Landnutzung.

Fossile Geschäftsmodelle werden nicht abrupt abgeschaltet, sondern in geordnete Übergänge mit klaren Endpunkten überführt. Dabei sind soziale Abfederung, Standortfragen und Versorgungssicherheit mitzudenken.
\end{quote}

- Sektorale Pfade. Für die wesentlichen Emissionssektoren (insbesondere Strom- und Wärmeerzeugung, Industrie, Verkehr, Gebäude, Landwirtschaft, Landnutzung und Abfall) sind sektorale Transformationspfade zu entwickeln, die mit dem nationalen Budget und den Zielen nach Art. 2a vereinbar sind.

- Fossile Übergänge.
- Fossile Infrastrukturen und Geschäftsmodelle sind auf einen auslaufenden Betriebspfad auszurichten, der Investitionssicherheit mit Klimaverträglichkeit verbindet.

- Unterstützende Instrumente (z. B. Transformationskredite, Stilllegungsprämien, Qualifizierungsprogramme) sollen den Übergang flankieren und soziale Härten abfedern.

- Cluster-Ansätze. Industrie- und Energiecluster können eigene Transformationspfade entwickeln, die mit MRV- und Finanzierungslogiken (Art. 2c, 2d, 2e) verzahnt sind.

- Just Transition. Bei der Ausgestaltung der sektoralen Pfade sind die Prinzipien einer gerechten Transformation (Just Transition) zu berücksichtigen, insbesondere mit Blick auf Beschäftigte, Regionen und einkommensschwache Haushalte.

- Interaktion mit anderen Rechtsbereichen. Sektorale Pfade nach diesem Artikel sind mit bestehenden Fachgesetzen (z. B. im Energie-, Verkehrs-, Bau- und Landwirtschaftsrecht) zu koordinieren.

{\footnotesize\begin{quote}
Artikel 2f schafft den Rahmen, in dem spätere Fachartikel und Verordnungen konkrete Sektorstrategien entwickeln, ohne das Gesamtziel aus den Augen zu verlieren.
\end{quote}}



%% ----------------------------------------
\section[Artikel 2g: Transparenz, Beteiligung und Ampel-Mechanismus]{Artikel 2g: Transparenz, Beteiligung und Ampel-Mechanismus}
\noindent\textit{ID: 12 -- artikel2g (meta)}

\begin{quote}
Was passiert, wenn wir vom Kurs abweichen? Artikel 2g beantwortet diese Frage: Er legt fest, wie Fortschritt und Rückstände transparent gemacht werden, wie Beteiligung organisiert ist und wie das Ampel-System funktioniert, das politisches Handeln auslöst, statt es zu blockieren.

Die Ampel ist ein Frühwarn- und Priorisierungsinstrument, kein automatischer „Selbststopp“. Sie zeigt, wo Handlungsbedarf besteht, und verknüpft dies mit klaren Prozessen zur Nachsteuerung nach Art. 17.14.
\end{quote}

- Transparenz.
- Zentrale Kennzahlen zu Emissionen, Budgets, Pfaden und Finanzverpflichtungen sind in einem öffentlichen Dashboard aufzubereiten, das für Fachleute und interessierte Bürgerinnen und Bürger verständlich ist.

- Die zugrunde liegenden Daten sind, soweit rechtlich zulässig, als offene, maschinenlesbare Daten bereitzustellen.

- Beteiligungsformate. Zur Ausgestaltung von Methoden, Verordnungen und Prioritäten sind geeignete Beteiligungsformate vorzusehen, einschließlich Konsultationen, Fachdialogen und Pilotprogrammen mit begleitender Evaluation.

- Ampel-System.
- Für zentrale Zielgrößen (Emissionen, Budgetausschöpfung, MRV-Integrität, Finanzpuffer u. a.) werden Indikatoren definiert, die in einem Ampelschema (Grün/Gelb/Rot) dargestellt werden.

- Die Ampelindikatoren und die zugehörigen Schwellenwerte und Reaktionspflichten werden in Art. 17.14 und den dazugehörigen Verordnungen näher geregelt.

- Ein roter oder gelber Status löst grundsätzlich eine Pflicht zur Prüfung und, soweit erforderlich, zur Nachsteuerung aus, ersetzt aber nicht die politische Entscheidung über konkrete Maßnahmen.

- Verzahnung mit Backstop. Die einfachgesetzlichen Ampel- und Nachsteuerungsmechanismen lassen den verfassungsrechtlichen Backstop nach Art. 1 unberührt. Sie dienen dazu, Zielabweichungen früh zu erkennen und zu beheben, bevor Notfallmechanismen greifen müssen.

{\footnotesize\begin{quote}
Artikel 2g macht sichtbar, ob das KlimaGG wirkt, und verknüpft diese Sichtbarkeit mit klaren Prozessen, ohne Automatismen zu schaffen, die demokratische Entscheidungsräume ersetzen.
\end{quote}}



%% ----------------------------------------
\section[VO 2a: Verordnung zur Zielbindung, Pfadfortschreibung und Backstop-Auslösung]{VO 2a: Verordnung zur Zielbindung, Pfadfortschreibung und Backstop-Auslösung}
\noindent\textit{ID: 13 -- vo2a (meta)}

\begin{quote}
Diese Verordnung konkretisiert Artikel 2: Methodik, Budget \& jährliche Obergrenzen, Zwischenziele, Trigger/Fristen, Backstop-Zwischenmaßnahmen, Transparenz und EU-Kohärenz. So werden die Ziele und Pfade des KlimaGG regelbasiert an neue wissenschaftliche Erkenntnisse, insb. IPCC-Berichte, angepasst werden -- transparent, mit klaren Fristen und demokratischer Kontrolle. Das schafft Planungssicherheit und verhindert Ad-hoc-Politik: Aktualisierungen folgen einem festen Verfahren, nicht Stimmungen.
\end{quote}

{\footnotesize\begin{quote}
VwV-Details (Prüfschemata, Datenformate, Rollenmodelle) können gesondert als Verwaltungsvorschrift veröffentlicht werden.
\end{quote}}


- Gegenstand, Zweck und Rechtsgrundlage. Diese Verordnung bestimmt Verfahren und Parameter zur Umsetzung von Artikel 2 KlimaGG. Rechtsgrundlage ist Artikel 2 (Zielsetzung) in Verbindung mit Artikel 10 (Institutionen) und Artikel 17.14 (Monitoring \& Open Data).

- Begriffe und Messmethodik. Bilanzierungsstandard: CO\$\_2\$-Äquivalente nach GWP100; Basisjahr 1990. LULUCF wird separat bilanziert; Anrechenbarkeit von Entnahmen setzt Zusätzlichkeit, Permanenz, robustes Monitoring und einen Kappungswert voraus. Bei methodischer Unsicherheit sind konservative, risikobasierte Annahmen anzuwenden.

- Nationales Emissionsbudget und jährliche Obergrenzen. Das nationale Budget und die jährlichen Obergrenzen werden als maschinenlesbare Datensätze veröffentlicht (Versionierung inkl. Änderungsjournal). Die Wissenschaftliche Klimakommission schlägt Budget und Pfad vor; die Bundesregierung setzt sie um.

- Zwischenziele und Pfadkorrektur. Zwischenziele (2035, 2040) folgen dem festgestellten Budgetpfad; Anpassungen erfolgen ambitionsneutral. Überholte Annahmen werden in den Datensätzen kenntlich gemacht; rückwirkende Neu-Bewertungen werden transparent versioniert.

- Trigger für Überprüfung. Neue IPCC-Berichte, methodische Updates international anerkannter MRV-Standards oder belegte Änderungen relevanter Budgets lösen eine Prüfung aus. Die Vorprüfung beginnt unverzüglich nach Veröffentlichung der maßgeblichen Quelle.

- Verfahren, Fristen und Beteiligung. Vorprüfung binnen 6 Wochen; Entwurf binnen 6 Monaten mit Wirkungs- und Verteilungsanalyse. Öffentliche Konsultation mind. 30 Tage; Auswertung binnen 30 Tagen. Die Verordnung tritt 3 Monate nach Veröffentlichung in Kraft, sofern der Bundestag nicht innerhalb dieser Frist mindestens gleichwirksame Abweichungen beschließt (parlamentarischer Widerspruchsvorbehalt).

- Backstop und interimistische Sicherungsmaßnahmen. Bei absehbarer Verfehlung von Budget oder Pfad wirken automatisch: ein zeitanteiliger Korrekturfaktor im Jahrespfad und eine temporäre Anhebung des CO\$\_2\$-Mindestpreis-Floors (Art. 7), soweit zur Pfadeinhaltung erforderlich. Die Maßnahmen enden mit Inkrafttreten der angepassten Verordnung.

- Transparenzpflicht. Brutto- und Netto-Minderungen werden getrennt ausgewiesen, nach CR-Linie, Sektor und Jahr.

- Transparenz und Open Data. Datensätze, Modelle, Annahmen, Sensitivitäten und ggf. Minderheitsvoten werden maschinenlesbar veröffentlicht (Art. 17.14). Versionierung und Änderungsjournal sind verpflichtend.

- EU-/Völkerrechtskohärenz. Erforderliche Notifizierungen und Kohärenzprüfungen erfolgen nach Artikel 17.1.

- Inkrafttreten und Übergang. Diese Verordnung tritt am ersten Tag des dritten Monats nach Verkündung in Kraft. Bis zum Inkrafttreten gelten die Sicherungsmechanismen nach § 7. [ToDo: Woher kommt das? Referenz anpassen]


%% ----------------------------------------
\section[Artikel 3: Transformationskredit (TFK)]{Artikel 3: Transformationskredit (TFK)}
\noindent\textit{ID: 14 -- artikel3 (gesetz)}

\begin{quote}
Dorfwärme eG, Oberpfalz. In T. wird ein gemeinsame Wärmeversorgung geplant. Der Transformationskredit macht’s für alle möglich. Für Standardmaßnahmen der Wärmewende ist er direkt verfügbar: Zentral kommt ein Wärmenetz, Außenhöfe installieren dezentrale Syteme. Die Dorfwärme eG nimmt für gemeinnützige Beratungsleistung und die Wärmetechnik einen Kredit auf, der eine Zinsstützung und Garantie aus dem kommunalen Pool erhält. Die Haushalte nutzen ihre Standardkredite für eigene Wärmepumpen. Wer seinen Kreditrahmen nicht voll ausschöpft, bereichert den Kommunal-Pool mit einem Anteil - und Familien können für ihre Kinder ins gemeinsame Wohnumfeld investieren. So werden die Keller sauber, die Rechnungen planbar, die Räume warm ohne Öltankgeruch. Der Kredit wird durch Energieeinspaarungen und Klimaleistungen zurückgezahlt - und ganz wichtig: Keiner bleibt im Kalten sitzen. CO\$\_2\$e vorher: 550 t/Jahr CO\$\_2\$e nachher: 200 t/Jahr

  Der Transformationskredit macht Klimainvestitionen für alle planbar - ob Haushalte, Handwerk, Mittelstand oder Kommunen. Standardmäßig zinsfrei (Variante A), zweckgebunden und digital verwaltet senkt er Einstiegshürden und bündelt privates Engagement mit öffentlicher Unterstützung. So entstehen verlässliche Wege zur Senkung von Energiekosten, zur Modernisierung von Gebäuden, Fahrzeugen und Prozessen -- ohne Bürokratie-Marathon. Annahme: \textasciitilde{}10.000€ pro Person. [ToDo: Annahmen sollen in allen Artikeln eine Richtung vorgeben, Verständnis fördern und geben den Geschichten in der jeweiligen Einleitung eine Basis. ToDo: Integration der Detail-Zielsetzung aus Artikel 2 auch hier.]
\end{quote}

- Zweck \& Anspruch. Anspruchsberechtigt sind natürliche und juristische Personen sowie Körperschaften des öffentlichen Rechts für Investitionen zur Emissionsminderung, nachhaltigen Energieversorgung und Effizienzsteigerung einschließlich gemeinschaftlicher Projekte. Querschnittsregeln zu Daten/IDs, Once-Only, Datenschutz und Kontrolle ergeben sich aus Art. 17.1, 17.4, 17.12 und 17.14.

- Verwendungsbereiche. Förderfähig sind insbesondere: a) Gebäude/Wärme (Sanierung, Wärmepumpen, Netze), b) Eigenstrom \& Netzintegration, c) Speicher \& Lastmanagement, d) Mobilität/Flotten, e) Prozesse/Elektrifizierung, f) gemeinwohlorientierte Quartiers-/Kommunalprojekte. Effiziente Kühlung in Kopplung mit Eigenstrom ist zulässig.

- Produktfamilie A/B/C. A (Regelfall für Haushalte/KMU): zinslos, automatische Teiltilgung aus Einsparungen/Rewards. B: Niedrigzins-Annuität mit teilweiser Abtretung. C: Marktmodell mit befristeter staatlicher Portfolio-Garantie/Pooling. Schwellen/Parameter legt VO 3a fest.

- Rückzahlung, Abtretung \& Rang. Tilgung erfolgt vorrangig aus eingesparten Energiekosten und zugewiesenen Carbon-Reward-Zuflüssen (Art. 5). Künftige Ansprüche können zugunsten des TFK abgetreten werden; Priorität und Verrechnung entstehen mit Eintragung im Register nach den Schnittstellen-/Signaturregeln des Gesetzes (Art. 17.12, 17.14). Doppelzählung/Doppelförderung ist ausgeschlossen (Art. 17.1, 17.4).

- Träger \& Finanzierungsrahmen. Der TFK wird als Sondervermögen des Bundes errichtet. Wirtschaftsplan, Tilgungsregeln, Rücklagen, Transparenz und ESA-/Beihilfe-Kohärenz richten sich nach Art. 17.2; Finanzierung erfolgt nicht über Netzentgelte/Umlagen (Art. 17.17).

- Familienregelung. Gesetzliche Vertreter können zweckgebundene TFK-Darlehen für das familiäre Wohnumfeld Minderjähriger aufnehmen. Minderjährige treten keinen persönlichen Schuldbeitritt an; Verbraucherrecht bleibt unberührt.

- Härtefälle \& Missbrauch. Bei Härtefällen sind Stundung/Restrukturierung, in Ausnahmefällen teilweiser Erlass möglich. Zweckwidrige Verwendung führt zu Rückforderung; Sanktions-/Clawback-Regeln nach Art. Art. 17.6 [ToDo: Art. Art. 17.6 richtig referenzieren und ggf. entwerfen] gelten entsprechend.

- Abwicklung \& Aufsicht. Digitale Antragstellung, Vergabe und Monitoring erfolgen über die Transformationsagentur (Art. 10) unter Nutzung der eAkte/Once-Only-Schnittstellen (Art. 17.4). Wertpapier-/Verbriefungsaktivitäten unterliegen den jeweils einschlägigen aufsichtsrechtlichen Vorgaben (Einzelheiten VO 3a).

- Freiwillige Poolzuweisung. Anspruchsberechtigte können ungenutzte Jahresansprüche freiwillig an Bundes-, Landes- oder Kommunal-Pools zuweisen. Poolmittel dürfen ausschließlich für Zinsstützungen, Portfolio-Garantien vorgeprüfter Vorhaben und gemeinnützige Beratungszuschüsse verwendet werden; direkte Darlehen an Private aus Poolmitteln sind ausgeschlossen. Näheres regelt VO 3a in Kohärenz zu Art. 17.6 (KH) und 17.14 (AR).

- Verordnungsermächtigung. Die Bundesregierung wird ermächtigt, durch Rechtsverordnung (VO 3a) Parameter/Prozesse festzulegen, insbesondere Produktparameter (A/B/C), Schwellen, Register-/Verrechnungslogik, Refinanzierungspools, Poolverfahren, Beratungszuschüsse, Risikoreporting und Pilotierung, jeweils kohärent zu Art. 17.* [ToDo: Verweis korrigieren].


%% ----------------------------------------
\section[VO 3a: Durchführung Transformationskredit]{VO 3a: Durchführung Transformationskredit}
\noindent\textit{ID: 15 -- vo3a (meta)}

\begin{quote}
Anmerkung: Die Höhe des Kredites (A) sollte Haushalten die nötigen Investitionen ermöglichen, d.h. \textasciitilde{}10.000€. Hier gibts auch ne gute Geschichte zu!

  
  Diese Verordnung konkretisiert Art. 3: Produktparameter, technische Register-/Prozessregeln, Refinanzierung sowie die Poollösung. Querschnittsvorgaben ergeben sich aus Art. 17 (Haushalt/ESA/Beihilfe, Durchsetzung, ID/eIDAS, Monitoring/Open Data, Kommunalhaushalt, Energiepreise/Netzentgelte).
\end{quote}

- Produktparameter A/B/C. Festlegung von Referenzzins und Aufschlägen (B/C), Ziel-DSCR, Max-Laufzeiten, sicherheitenfreien Schwellen, automatischem Fallback A→B bei Unterdeckung, Frühtilgungsregeln sowie einem standardisierten einmaligen Optionswechsel (Bewertungsformel).

- Register- \& Verrechnungslogik. Technische Abtretung und Prioritätsrang der Zahlungsansprüche (Energieeinsparungen, Reward-Zuflüsse) inklusive Idempotenz/Conflict-Resolution. Statusändernde Ereignisse sind nach ID-VO (Art. 17.12) qualifiziert zu signieren/siegeln und mit qualifizierten Zeitstempeln zu versehen; Datenflüsse werden once-only an das Ampel-Reporting (Art. 17.14) gespiegelt.

- Refinanzierung \& Verbriefung. Bildung von Refinanzierungspools und zulässigen Tranchen; Mindestreserven, Stresstests, Anlegerberichte. Ausgabe elektronischer Wertpapiere nach eWpG und -- soweit einschlägig -- Prospekt-VO/Verbriefungs-VO [ToDo: passen diese Verweise?]. Aufsichtsrechtliche Einordnung (u. a. BaFin) richtet sich nach den einschlägigen Normen.

- Schwellen \& Standardanwendung. Betrags-/Projekt-/Größenklassen, ab denen B Standard wird und C zu prüfen ist; für Haushalte/KMU gilt grundsätzlich A. Schnittstellen zu Beihilfe-Schwellen gemäß Art. 17.2.

- Poollösung Bund/Land/Kommune.
- Zweck: Finanzierung ausschließlich von Zinsstützungen, Portfolio-Garantien für vorgeprüfte Vorhaben (z. B. Bürgerenergie, Speicher, Wärmenetze, Effizienz kommunaler Liegenschaften, offene Forschung) sowie gemeinnützigen Beratungszuschüssen.

- Zuweisung: Freiwillige Zuweisung bis zu 50 \% ungenutzter Jahresansprüche; Widerruf bis zur projektbezogenen Bindung.

- Verteilung \& Caps: Jahreskontingente je Ebene; Schlüssel (Einwohner/Bedarf); Projekt-Cap (max. 15 \% des Jahreskontingents); Re-Allokation nicht gebundener Mittel.

- Auswahl \& Umsetzung: Vorprüfung (CO\$\_2\$-Wirkung/€, Reifegrad, sozialer Nutzen, Open-Access-Kriterien); Umsetzung über Vergaberecht (TCO inkl. CO\$\_2\$) oder beihilfekonforme Standardmodule; Register-Sperrlogik gegen Doppelförderung.

- Transparenz \& Buchung: Öffentliche Dashboards (Zuweisungen, gebundene Projekte, CO\$\_2\$-Wirkungen, Mittelabrufe, Restvolumen) in das AR-Dashboard integriert (Art. 17.14); Haushalts-/Kontierung nach KH-Schema (Art. 17.6).

- Abwicklung \& Aufsicht: Technische Abwicklung über Förderbanken; beihilferechtliche Einordnung nach CEEAG/GBER; jährliche Assurance-Prüfung.

- Pilot \& Evaluierung: 12--24 Monate Pilot mit begrenzten Kontingenten; Parameternachsteuerung per VO.

- Gemeinnützige Beratungszuschüsse.
- Gegenstand: Zuschüsse für Erst-/Orientierungsberatung (z. B. Quartiers-, Wärmenetz-, Effizienz-Checks) mit Gemeinwohlbezug; keine projektindividuelle Bauüberwachung.

- Träger \& Neutralität: Gemeinnützige Organisationen, Kommunen, Bürgerenergie/Genossenschaften; bei Doppelfunktion ist eine organisatorische Trennung (Firewall) verpflichtend.

- Sätze: Deckelnde Pauschalen/Höchstbeträge; offene Providerliste; einfache Eignungsprüfung; Veröffentlichung der Sätze.

- Abgrenzung: Laufende Projektsteuerung, individuelle Fachplanung/Bauüberwachung sind nicht zuschussfähig aus Poolmitteln.

- Risikobericht \& Parametertuning. Zweijährlicher Bericht zu Ausfallrisiken und Rückflusslage; VO-basierte Anpassungen (Schwellen, Reservesätze, DSCR-Ziele). Ampel-Trigger und Haushaltsanpassungen bleiben Art. 17.2/17.14 vorbehalten.

- Härtefall-Standard. Einheitliche Kriterien und Verfahren für Stundung/Restrukturierung (Fristen, Nachweise, Entscheidungswege), koordiniert mit Verbraucherschutz- und Sozialregelungen; Veröffentlichung als VwV-Anlage.

- Inkrafttreten \& Übergang. Stufenplan: M0--M6 Produkt A; M7--M12 A/B; M17--M24 C/Refi-Pools/Pilotzonen. Evaluation und ggf. VO-Anpassung im Jahresbericht nach Art. 17.14.


%% ----------------------------------------
\section[Artikel 4: Transformationsförderung]{Artikel 4: Transformationsförderung}
\noindent\textit{ID: 16 -- artikel4 (gesetz)}

\begin{quote}
Brauerei H., Oberfranken. Die Sudkessel liefen bisher mit Gas. Ein Preiswettbewerb „Prozesswärme ohne Fossil“ prämiert die beste Pilotlösung. Ein Mittelständler entwickelt eine Hochtemperatur-Wärmepumpe, die 120 \$\textasciicircum{}\textbackslash\{\}circ\$C schafft. Die Brauerei bewirbt sich als Reallabor und bekommt einen Ansprechpartner bei der Agentur für Genehmigungen und Meilenstein-Förderung: zuerst Machbarkeit, dann Einbau, dann Monitoring. Nach einem Jahr sinkt der Gasverbrauch gegen null, die Technik ist marktreif -- und zehn weitere Brauereien steigen ein. CO\$\_2\$e vorher: 786 t/Jahr CO\$\_2\$e nachher: 424 t/Jahr

Manche Vorhaben brauchen mehr als Kapital: Koordination, Testfelder, Genehmigungswissen. Die Transformationsförderung schließt diese Lücke -- schlank, wettbewerblich und mission-orientiert --, damit gute Ideen rasch in den Markt kommen und Deutschland bei Schlüsseltechnologien vorn bleibt.
\end{quote}

- Gegenstand und Abgrenzung. Gefördert werden marktnahe Innovation, Demonstration (TRL ≥ 6) und First-of-a-Kind-/Skalierungsvorhaben mit signifikanter Treibhausgasminderung, Systemintegration oder Flexibilitätsnutzen. Grundlagenforschung (Lehre/Tenure) ist nicht Gegenstand dieses Artikels; sie wird querschnittlich geregelt (Art. 4b). [ToDo: Wo und wie ist Grundlagenforschung geregelt?]

- Förderlinien.
- (A) Pilot \& Demonstration: Validierung im realen Betrieb, MRV-fähig.

- (B) First-of-a-Kind/Skalierung: Erst- oder frühe Serienanlagen, Ramp-up/Lerneffekte.

- (C) Missions-/Preiswettbewerbe: Zielklar definierte Auslobungen mit Preisgeld/Abnahme-Versprechen.

- (D) Cluster \& Reallabore: Reale Testfelder mit beschleunigten Verfahren (DV-Verweis) und Datenpflichten.

- (E) Standardisierung \& Open Tech: Interoperabilität, MRV-/Daten-APIs, offene Referenzimplementierungen.

- Förderprinzipien. Wettbewerb, Transparenz, Meilenstein- und Ergebnisauszahlung; Kumulierung mit Art. 3 (TFK) und Art. 5 (Carbon Reward) nur ohne Doppelförderung derselben Tonne und unter Anrechnung in der Beihilfeintensität. Once-Only-Datenfluss und Register-Sperrlogik nach Art. 17.1/17.14/Art. 17.11/Art. 17.12. [ToDo: Art. 17.11 und Art. 17.12 korrekt referenzieren und ggf. Artikel schreiben.]

- Koordination \& Verfahren. Die Transformationsagentur (Art. 10) führt wettbewerbliche Verfahren (Calls/Preise), koordiniert Genehmigungen als Single-Point-of-Contact und nutzt die Beschleunigungs- und Fiktionsregeln aus Art. 17.4 (DV). Digitale Akte, Zustellung, Signaturen gem. Art. 17.12 (ID-VO).

- Förderformen. Zuschuss (investiv/operativ), Preisgeld, rückzahlbare Zuschüsse, beteiligungsähnliche Elemente, ggf. zielgebundene Abnahme-/Verfügbarkeitsverträge. Ausgestaltung beihilferechtskonform (CEEAG/GBER); Beihilfe-Gate vor Bewilligung.

- Förderkonditionen. Förderfähigkeit, Kostenarten, Obergrenzen und Intensitäten werden in der VO 4a festgelegt -- differenziert nach Linie, Unternehmensgröße und Klimawirkung. Förderungen sind meilensteinbasiert; Nichterreichen → Anpassung/Abbruch/Clawback (Art. Art. 17.6).

- MRV \& Offenlegung. Klimawirkungen und Systemnutzen werden nach Art. 17.14 (AR-VO) erfasst und maschinenlesbar veröffentlicht (Open Data) -- Schutz berechtigter Betriebsgeheimnisse gewahrt. Projektdaten werden an das Registry gespiegelt (Art. 17.11) zur Kumulierungskontrolle.

- Vergabe \& TCO. Bei öffentlichen Projektteilen gelten die Vergaberegeln mit TCO/CO\$\_2\$-Kriterien (Art. 8a/VG); Standardmodule können per Mustervergabe genutzt werden.

- Sozial-/Arbeits-Schnittstelle. Qualifizierung/Übergänge werden über Art. 17.7 (SA/TKUG) angebunden; Doppelstrukturen werden vermieden.

- Transparenz \& Evaluation. Jährliche Wirkungsberichte (CO\$\_2\$-/€, Skalen-Effekte, Markteintritt, Standardisierungsbeiträge), Audit/Assurance mind. „limited“. Veröffentlichung im Ampel-Reporting (Art. 17.14).

- Rechts- und Sanktionsbezug. Verstöße gegen Förderauflagen, Daten- und Registerpflichten unterfallen Art. Art. 17.6 (Clawback, Bußgelder, Sperren); Rechtsschutz nach Art. 8.

- Verordnungsermächtigung. Die Bundesregierung wird ermächtigt, durch Rechtsverordnung (VO 4a) die Durchführung zu regeln -- u. a. Linienparameter, Intensitäten, Auswahl-/Scoringkriterien, Meilensteinlogik, Open-Data-/Standardisierungsauflagen, Kumulierungs-/Anrechnungsregeln, sowie Verfahren der Preiswettbewerbe und Reallabore.

{\footnotesize\begin{quote}
Schlanke, marktorientierte Förderung: Demonstration \& Skalierung beschleunigen, Doppelstrukturen vermeiden, Daten \& Standards als Hebel.
\end{quote}}



%% ----------------------------------------
\section[VO 4a: Durchführung Transformationsförderung]{VO 4a: Durchführung Transformationsförderung}
\noindent\textit{ID: 17 -- vo4a (meta)}

\begin{quote}
Diese Verordnung konkretisiert Artikel 4: Förderkonditionen, Verfahren, Kumulierungskontrolle und Datenschnittstellen -- im Einklang mit Art. 17 (DV/ID/AR/Harmonisierung).
\end{quote}

{\footnotesize\begin{quote}
Keine Doppelstrukturen: Genehmigung (Art. 17.4), Ident/Signatur (17.12), Monitoring/Open Data (17.14) und Beihilfe-Gate sind angebunden; Grundlagenforschung folgt eigenständig (Art. 4b).
\end{quote}}


- Budgetierung \& Quoten. Jährliche Budgetfestlegung mit Linienquoten (A--E) und Mindestanteil für Missions-/Preiswettbewerbe; überjährige Übertragbarkeit zulässig.

- Förderintensitäten \& Kostenarten.
- Linie A (Pilot/Demo): Fördersatz bis [VO: \%] förderfähiger Kosten (Anlagen/Integration, MRV/Monitoring, Genehmigungs-/Planungskosten anteilig).

- Linie B (FOAK/Skalierung): Fördersatz bis [VO: \%] (CapEx-Lernkurven, Inbetriebnahme, Systemintegration); degressiv mit Serienzahl.

- Linie C (Mission/Preis): Preisstaffeln, „pay-for-performance“; optional Abnahme-/Verfügbarkeitsverträge.

- Linie D (Reallabore): Infrastruktur-/Betriebskostenzuschüsse befristet; DV-Fast-Track, einheitliche Daten-/Sicherheitsprofile.

- Linie E (Standards/Open Tech): Pauschalen für Normungsbeiträge, Referenzimplementierungen, Interop-Tests, offene MRV-Toolchains.

    Förderintensitäten werden nach Unternehmensgröße, Spillover und CO\$\_2\$-/€-Wirkung differenziert; Beihilfeobergrenzen (CEEAG/GBER) sind einzuhalten.

- Auswahl \& Scoring. Mehrstufiges Verfahren (Eligibility → Technik/MRV → Wirtschaft/Skalierung → Systemnutzen/Netzdienlichkeit). Kernkriterien: CO\$\_2\$-Minderung pro €, Reife/Skalierbarkeit, Netz-/Flex-Nutzen, europäische Interoperabilität, Zeit bis Wirkung, Risiko/Management, Open-Science-/Standardisierungsbeiträge.

- Meilensteine \& Auszahlung. Standard-Meilensteine (Konzept, Genehmigung, Inbetriebnahme, Betriebsnachweis, Transfer). Auszahlungen transchenweise gegen QES/QSeal-gesicherte Nachweise (Art. 17.12); Abbruch-/Pivot-Regeln, Clawback bei Nichterreichen/Fehlangaben.

- Kumulierung \& Doppelförderung. Anrechnung von Art. 3 (TFK) und Art. 4 (Carbon Reward) auf Förderäquivalente; Register-Sperrvermerk (Art. 17.11/Art. 17.12) gegen parallele Inanspruchnahme derselben THG-Einheiten; Zulässigkeit von REWARD + Zuschuss, wenn sachlich getrennte Gegenstände (CapEx vs. verifizierte Tonnen) und Beihilfeobergrenzen eingehalten werden.

- Reallabore. Auswahlverfahren, Zonierung, UVP-/Planbeschleunigung gem. Art. 17.4; Sicherheits-/IT-Profile (NIS2/BSI), BCM (RTO/RPO) und Incident-Meldungen wie Art. 17.14; offene Datenspezifikation (JSON/CSV), Versionierung/Journal.

- Missions-/Preiswettbewerbe. Zielbeschreibung, Teilnahmebedingungen, Bewertungsmatrix, Preisstaffeln, Fristen; Nachnutzung unter fairen, transparenten Bedingungen (Open-Access-Daten; IP-Regeln mit Schutz von Geschäftsgeheimnissen).

- Daten \& Schnittstellen. Einlieferung über die Klima-Datenplattform (Art. 17.14): /v1/funding/projects, /v1/funding/milestones, /v1/funding/results; Pflichtfelder (Projekt-ID, Anlagen-ID, LEI/EORI, Serien-/Reward-Refs, CO\$\_2\$-/€-Kennzahlen). Once-Only-Prinzip; QSeal-gesiegelte Snapshots.

- Vergabe \& TCO. Mustervergaben mit TCO inkl. CO\$\_2\$ (Art. 8a/VG), Safe-Harbor-Standards für häufige Module (z. B. Wärmenetze, Speicher, Flex-IT); Open-Book-Prinzip optional.

- Transparenz \& Audit. Veröffentlichung bewilligter Projekte, Mittelabfluss, Kennzahlen; Assurance (mind. „limited“; Stichproben „reasonable“) und jährlicher Wirkungsbericht in das Ampel-Dashboard (Art. 17.14).

- Rechtsfolgen \& Sanktionen. Fehlgebrauch, Falschangaben, Verstoß gegen Daten-/Registerpflichten: Rückforderung, Zinsen, Ausschluss (Fast-Track-Sperre), Veröffentlichung aggregierter Verstöße; Verfahren nach Art. Art. 17.6, Rechtsschutz nach Art. 8.

- Übergang \& Evaluation. Pilotphase 12--24 Monate mit Kontingenten; Evaluation (CO\$\_2\$-/€, Time-to-Market, Standardisierungsoutput, Netz-KPIs). Parametrische Anpassung per VO-Änderung.


%% ----------------------------------------
\section[Artikel 5: Carbon Reward]{Artikel 5: Carbon Reward}
\noindent\textit{ID: 18 -- artikel5 (gesetz)}

\begin{quote}
Sägewerk H., Oberfranken. Um sechs riecht die Halle nach Harz und frischem Schnitt. Der Gabelstapler summt, Sägemehl liegt wie Puder auf den Schuhen. Früher fraß die Trockenkammer Gas im Dauerbrand -- eine gelbe Flamme für Geld und CO\$\_2\$. Heute landen Rinde und Verschnitt im neuen Ofen. Früher machte der Köhler aus Holz schwarzes Gold; später brauchte die chemische Industrie Holz für Essig und Teere -- und Grillkohle. Heute kehrt die moderne Pyrolyse zurück: Keine Flamme, sondern schwarze Glut: Aus den Resten im Sägewerk entsteht heißes Prozessgas für die Trocknung. Der Alltag ändert sich kaum -- Bretter stapeln, Stapel drehen, Kammertüren schließen. Neu ist die Anzeige am Bürorechner: „Klimaleistung vergütet“. Die Zahlung kommt, weil jetzt fossile Wärme ersetzt wird. Planbar, messbar, auditierbar, jeden Monat, so sicher wie die Lohnliste. Der Banker, sonst knauserig, nickt: „Das rechnet sich.“ Was früher Abfall war, wird Teil der Geschichte: Aus einem Teil der Reststoffe wird Biokohle, sauber verpackt auf Paletten. Die Landwirte aus der Umgebung können gar nicht genug davon bekommen, und von der neue Produktkennzeichnung. Denn die fährt mit. Und wer sie später auf dem Feld einsetzt, kann den Klimawert in seinem Namen verbuchen. Dazu kommt aus dem Prozess etwas Pyrolyseöl, das die regionale Industrie abnimmt. Hier im Werk bleibt vor allem die Wärme -- und die Gewissheit, dass jeder Stapel getrockneter Bretter weniger CO\$\_2\$ gekostet hat. Für die Belegschaft bedeutet es Routine mit Rückenwind -- der gleiche Handgriff, nur klüger. Und wenn abends das Tor zugeht, bleibt ein Gefühl, das man sehen kann: Stapel für Stapel entsteht Wert -- fürs Haus nebenan, für den Wald vor der Tür, und für ein Konto, das Klima heißt. CO\$\_2\$e vorher: 1.800 t/Jahr CO\$\_2\$e nachher: 350 t/Jahr

  Carbon Rewards vergüten nachweisbare Klimaleistung: vermiedene Emissionen, dauerhaft entnommenes CO\$\_2\$e und staatlich beschlossene Verzichts-/Stilllegungsentscheidungen. Das System funktioniert einfach: Projekte werden standardisiert gemessen und geprüft, im Register erfasst und erhalten -- je nach Linie -- eine Mindestvergütung pro Tonne (Floor). Kreislaufwirtschaft (u. a. hochwertiges Recycling, Re-Use, Substitution klimarelevanter Produkte) und LULUCF (Landwirtschaft, Forst, Moore, Böden -- insbesondere CO\$\_2\$e-Aufbau, Langzeitspeicherung, Stilllegung) sind explizite und gleichwertige Säulen der Carbon Rewards. So entstehen planbare Einnahmen, die Investitionen sofort finanzierbar machen. Für alle Sektoren gilt: Doppelzählung ist ausgeschlossen, Bestandschutz (Vintage) gibt Sicherheit. Die Linie C bleibt grundsätzlich geschlossen und wird nur für klar benannte öffentliche oder internationale Fälle geöffnet.
\end{quote}

- Zweck, Geltung, Linien. Vergütung überprüfbarer Minderungen/Entnahmen in Energie, Industrie, Gebäuden, Verkehr, Kreislaufwirtschaft und LULUCF. Linien: A Vermeidung/Substitution; B Entnahme/Speicherung (D100/D40 mit Puffer/Versicherung); C programmatische Verzichts-/Stilllegungsentscheidungen öffentlicher Stellen.

- Integritätsprinzipien. Zusätzlichkeit gegenüber Baselines; keine Doppelzählung und keine Doppelförderung derselben Tonne (Register-Sperrlogik). Beihilferechtliche Genehmigungen sind vor Auszahlung einzuholen (EU-Gate). Unions- und Völkerrecht bleiben unberührt.

- Preisleitplanken. Ein Floor-Korridor wird durch Rechtsverordnung festgelegt. Auszahlungen erfolgen mindestens zum Floor; marktliche Top-ups sind zulässig. Linien- und methodenspezifische Zahlungsmodi (z. B. Annuitäten/Puffer bei Linie B) bestimmt die Verordnung. Bei Cap-Ausschöpfung wird ein Reservierungsbescheid erteilt; der für die Serie bekanntgemachte Floor bleibt gewahrt.

- Kommunalbonus (Verweis). Auf Auszahlungen aus Carbon Rewards findet der Kommunalbonus gemäß Art. 17.17 § 2 (KAV-Neu) Anwendung. Einzug, Verteilung und Reporting richten sich nach Art. 17.17 und Art. 17.6 (KH).

- Volumensteuerung. Jahresobergrenzen (Caps) je Linie mit Ampel-Logik gemäß AR-VO; Banking/Borrowing und Quartalsfenster regelt die Verordnung.

- Claims (No-Offset). Keine Nutzung zur Erfüllung öffentlich-rechtlicher Emissionspflichten oder für pauschale „klimaneutral“-Aussagen laufender Emissionen. Zulässig sind produktgebundene Claims (über DPP) und Contribution-Claims -- jeweils nur nach Retirement im Register. Ausgestaltung, Belege und Sanktionen folgen Art. 17.5 (VC/Green-Claims).

- MRV, Register und Signaturen. Baselines, Assurance-Level, Datenformate, Serienstati und Ampel-Trigger folgen AR-VO (Art. 17.14); Identifizierung, Signaturen/Siegel, Zeitstempel und Zustellung folgen ID-VO (Art. 17.12). Schnittstellen zu ETS/ETS2/CBAM werden „once-only“ gespiegelt (Art. 17.14).

- Linie C -- Zugang \& Sicherungen. Linie C ist geschlossen; Zugang nur für programmatisch verhandelte Stilllegungen/Verzichte (Bund/Länder/Kommunen oder mehrheitlich öffentliche Unternehmen) sowie völkerrechtliche Abkommen. Voraussetzungen: Transparenz, angemessene Entschädigung, No-Restart-Covenant, Sozial-/Strukturplan. Für ETS-Anlagen erfolgt die Zertifikatslöschung nach Maßgabe der Spezialregelungen (z. B. KVBG/EN-VO).

- Sanktionen. Falschangaben, Doppelzählung, Reversal-/No-Restart-Verstöße führen zu Clawback, Serien-Sperre, Bußgeldern nach AR-VO/DV; Prüferhaftung bei grober Fahrlässigkeit/Vorsatz.

- Übergang \& Evaluation. Pilotphase 12--24 Monate mit Teilvolumina; Evaluation nach 24 Monaten; Anpassungen per Verordnung. Bestandsvintages bleiben unberührt.

{\footnotesize\begin{quote}
Hinweis für Kommunen: Der Kommunalbonus (KAV-Neu) wird zentral eingezogen und nach Art. 17.17/17.6 verteilt und veröffentlicht. Projektnahes Reporting erfolgt im AR-Dashboard (Art. 17.14).
\end{quote}}



%% ----------------------------------------
\section[VO 5a: Carbon Rewards im Elektrizitätssektor]{VO 5a: Carbon Rewards im Elektrizitätssektor}
\noindent\textit{ID: 19 -- vo5a (meta)}

\begin{quote}
Chemiepark R., Rheinland. Früh morgens sieht man, nur die Silhouetten der Kühltürme. Um sieben schiebt die Leitwarte den Tag an: Der Elektrolyseur taktet hoch, wenn der Wind die Leitungen füllt, senkt sich in windstillen Stunden und atmet aus dem Batteriespeicher nach. Aus Wasser wird Wasserstoff, aus Wasserstoff und CO\$\_2\$ entsteht Methanol. (In einem Pilotprojekt hat die Brauerei/Zementwerk der Stadt durch die Innovationsförderung die CO2-Pipeline gebaut.) Keine große Geste, sondern Taktarbeit zum Puls des Strommarktes. Neu ist, was im Hintergrund mitläuft: Jede MWh, die in sauberen Stunden genutzt wird, zählt. Jede flexible Pause entlastet das Netz. Die Vergütung kommt monatlich -- für die tatsächlich vermiedenen Emissionen und für die bereitgestellte Systemflexibilität. Der Werksleiter sagt bei der Schichtübergabe nur: „Wir fahren nach Wetter und Netz -- und werden dafür bezahlt, es klug zu tun.“ Am Gate verlässt die erste Charge E-Methanol das Gelände. Auf dem Lieferschein steht mehr als eine Nummer: Der Klimabeitrag fährt mit -- messbar, nachvollziehbar. Die Abnehmer mögen das; nicht als Freibrief, sondern als Zusatzwert, der den Einkauf leichter macht. Im Park selbst bleibt die Gewissheit: Hier liegt eine Schlagader des Klimaschutzes. CO\$\_2\$e vorher: 28.000 t/Jahr CO\$\_2\$e nachher: 5.000 t/Jahr - für  die Jahresproduktion von 20.000 t Methanol.

  Wichtig für die Energiewende: günstige Strompreise und garantierter, erneuerbarer Zubau. Wer hier deckelt, der richtet Deutschland zugrunde. 

  Strom erhält klare Methodenblöcke: AVG (durchschnittsbasiert) oder RT (zeitaufgelöst) für Erzeugung, FLEX für Speicher/Lastverschiebung sowie einen separaten NDZ-Zusatz für netzdienliche Verfügbarkeiten. PtX ist abbildbar; Doppelzählung wird ausgeschlossen.
\end{quote}

- Geltung \& Datenquellen. Diese VO konkretisiert Art. 5 für Strom. Emissionsfaktoren (EF\_avg/EF\_marg), Zeitraster, Netzdienstekategorien (NDZ) und § 14a-Ereignislogs werden ausschließlich aus der EN-VO/Datenplattform gespeist (Art. 17.17/17.14). Zuständigkeiten der BNetzA bleiben unberührt.

- Erzeugung (Tracks „AVG“/„RT“). AVG: CR = EF\_avg,y,z × MWh\_inj. RT: CR = Σ EF\_marg,t,z × MWh\_inj(t). Trackwahl bei Inbetriebnahme; Bindung ≥ 36 Monate; einmaliger Wechsel zulässig (VO-Bedingungen).

- Speicher \& Lastmanagement (FLEX). CR = Σ[ MWh\_out(t\$\_2\$) × EF\_marg,t\$\_2\$,z − (MWh\_in(t₁)/η) × EF\_marg,t₁,z ]; nur positive Salden sind rewardfähig. Baselines, Mess- und Nachweisregeln regelt die VO.

- Netzdienlichkeit (NDZ). Zusatz-Reward je zugesagter/gelieferter Systemdienst: CR\_NDZ = Σ[ MW\_avail,service(t) × W\_service(t,z) × κ\_service ]. Doppelte Anrechnung mit E-MWh/FLEX ist ausgeschlossen.

- Power-to-X (PtX). Anrechenbarkeit produkt- oder stromseitig; Zusätzlichkeit, Zeitabgleich, Systemgrenzen und Vermeidung der Doppelverwertung regelt die VO in Anbindung an EN-VO.

- Preisleitplanken \& Settlement. Linien-/Track-spezifische Floors/Höchstpreise und optionaler Stabilitätskorridor; Methodik (WACC/CAPEX/OPEX, Profil-/Vermarktungskosten) offenlegen. Automatisierte Meldungen (kWh/kW/MW, Zeitstempel, Zone, Dienstekategorie) an das Register; Abrechnung mind. monatlich. QES/QSeal/ Zeitstempel nach ID-VO.

- Anti-Gaming \& Beihilfe. Rewards nur für zusätzliche Wirkung; parallele Inanspruchnahme von E-MWh, FLEX, NDZ für dieselbe Einheit ausgeschlossen (AR-VO-Prüflogik). Beihilfekonformität (CEEAG/GBER) sicherstellen.

- Übergang \& Evaluation. Pilot 12--24 Monate; schrittweise Erweiterung; Bestandsvintages unberührt; regelmäßige Evaluation mit Anpassungen per VO.


%% ----------------------------------------
\section[VO 5b: Carbon Rewards für Landnutzung \& Kreislaufwirtschaft]{VO 5b: Carbon Rewards für Landnutzung \& Kreislaufwirtschaft}
\noindent\textit{ID: 20 -- vol5b (meta)}

\begin{quote}
Privatwald S., Odenwald. Morgens liegt Nebel zwischen alten Buchen, Spechte klopfen, irgendwo schlüpft ein Reh ins Dickicht. S. geht den vertrauten Weg, sieht an den hohen Stämmen hinauf. Früher rollte hier alle paar Jahre schweres Gerät und die besten Stämme gingen früh heraus. Jetzt soll der Wald alt werden -- mit Ruhe, dicken Kronen, Totholz und Schatten über dem Bach. Das Neue ist angenehm unspektakulär: Keine Auktion, kein Papierstapel. Einmal im Jahr kommt die Zahlung -- als verbindliche Jahresrate für die wachsende Kohlenstoffspeicherung und den Verzicht auf den schnellen Hieb. Ein Teil fließt in Pfadpflege und eine Bank am Aussichtspunkt, der Rest in das, was S. schon lange wollte: eine Hecke am Feldrand, Obstbäume für die Kinder. Der Förster markiert Sturmfurchen, S. macht sich weniger Sorgen: Fällt ein Baum, springt der gemeinsame Puffer ein; das Risiko teilt die Gemeinschaft. Der Wald dankt auf seine Weise: mehr Schatten in Hitzeperioden, weniger Starkregenrinnen, mehr Vogelstimmen. An Wochenenden kommen Nachbarn zum Pilzesuchen; am Rand steht ein schlichtes Schild: „Dieser Wald darf alt werden.“ CO\$\_2\$e vorher: 0 t/Jahr CO\$\_2\$e nachher: 350 t/Jahr - für die 60 Hektar Wald.

  Dieses Modul präzisiert Methoden, Daten und Sicherungen für Minderungen und Abscheidungen durch Landnutzung sowie Kreislaufwirtschaft. Ziele: vergleichbare Baselines, robuste Permanenz, klare Produktzuordnung und kein Greenwashing.
\end{quote}

- Geltung \& Sektoren. LULUCF (z. B. Aufforstung, Walderhalt/-management, Moorwiedervernässung, Humusaufbau, Agroforst, Biochar/Boden-C, langlebige Holzprodukte) sowie Kreislaufwirtschaft (Abfallvermeidung, Re-Use, hochwertiges Recycling, Material-/Produktsubstitution, Design für Langlebigkeit/Modularität).

- Methodischer Verweis. Baselines, Permanenzklassen (D100/D40), Puffer/Versicherung, Leakage, Qualitätsfaktoren, Massenbilanz und Prüflevel werden ausschließlich in LMRV (Art. 17.8) und KMRV (Art. 17.11) geregelt und dort fortgeschrieben.

- Preis-/Volumenleitplanken. Caps, Floor-Korridor und Zahlungsmodi (inkl. Annuitätenlaufzeiten für D-Klassen) setzt diese VO für Tätigkeitsgruppen fest; Bestandsvintages bleiben unberührt.

- Produktpass \& Claims. Produktgebundene Rewards werden im Digitalen Produktpass (DPP) gemappt (Serie, Vintage, Linie, Dauerhaftigkeitsklasse, Menge). Claims sind nur nach Retirement und nach den Regeln des Art. 17.5 (VC/Green-Claims) zulässig.

- Biodiversität \& Sozialstandards. Do-No-Harm-Mindeststandards (Lebensräume/Arten, Landrechte, Beteiligung Eigentümer) gelten; Dokumentation nach LMRV/KMRV. Regionale Gemeinwohlverträge sind empfohlen.

- Abrechnung \& Register. Datenfelder, Assurance, Serienstati, Zeitstempel, Zustellung und Open-Data folgen AR-VO/ID-VO; Schnittstellen zu ETS/CBAM werden „once-only“ gespiegelt (Art. 17.14).

- Übergang \& Pilot. 24 Monate Pilot mit Parametertabellen; Evaluation und VO-Anpassung; Versionierung der Faktoren/Leitfäden im Register.


%% ----------------------------------------
\section[Artikel 5c: Klima-Biodiversitäts-Korridore (KBK) -- Top-Up im Carbon-Reward]{Artikel 5c: Klima-Biodiversitäts-Korridore (KBK) -- Top-Up im Carbon-Reward}
\noindent\textit{ID: 21 -- artikel5c (meta)}

\begin{quote}
Korridore verbinden Lebensräume, kühlen Städte und puffern Hochwasser -- und verstärken die Klimawirkung vieler Projekte. KBK vergütet solche zusätzlichen Klima-/Resilienz-Co-Benefits optional als Top-Up zum Carbon-Reward, ähnlich der Netzdienstlichkeits-Logik im Strom.
\end{quote}

- Zweck \& Andockung: KBK ist ein Top-Up-Modul zu Art. 5 (Carbon Reward) für Projekte in den Linien A/B, die in ausgewiesenen Korridoren liegen oder diese funktionsfähig machen (z. B. Entsiegelung, Auen, Hecken-/Baumzüge, Moor-/Gewässervernetzung, urbane Blaue-Grüne Infrastruktur).

- Korridor-Kulisse: Die Agentur veröffentlicht eine nationale KBK-Karte (Bund/Länder/Kommunen, Natura-2000/RL-Öko-Netze integriert) mit Kategorien und Mindestbreiten; Pflege als Open-Data-Layer (Art. 17.14).

- Eligibility \& MRV: Top-Up nur bei Additionalität über die CO\$\_2\$-Baseline hinaus, nachgewiesener Konnektivität (z. B. Durchgängigkeit/Index), Resilienzfunktion (Hitze/Hochwasser) und Do-No-Harm (Arten/Habitate). Permanenz-/Pufferregeln aus Art. 5b gelten entsprechend.

- Vergütung: KBK-Top-Up als pro-tCO\$\_2\$e-Zuschlag oder als flächen-/horizontbezogene Annuität (D40/D100) mit Cap [VO] €/Projekt·Jahr. Die CO\$\_2\$-Mengen bleiben unverändert; der Top-Up vergütet ausschließlich den zusätzlichen Korridor-Nutzen (keine Doppelzählung).

- Finanzierung \& Schnittstellen: Mittel aus Transformationsförderung (Art. 3) und optional KKR-Top-Up (Art. 10c). Keine Anrechnung auf produktseitige Claims; Marketing-Claims folgen Art. 5 (No-Offset, Retirement).

- Verfahren: Standard-Onboarding über Art. 8a; Safe-Harbor-Kits (Typkorridore) und Muster-Verpflichtungen für Pflege/Sicherung (z. B. Grundbuch/Vertragsnaturschutz).

- Verordnungsermächtigung: Festlegung KBK-Kategorien, Konnektivitäts-/Resilienzmetriken, Mindestbreiten, Caps, Prüftiefe (limited/reasonable), Datenraster und Pflegepflichten; EU-/Naturschutzrecht bleibt unberührt (Art. 17).

{\footnotesize\begin{quote}
Ziel: CO\$\_2\$-Wirkung und ökologische/resiliente Funktionen gleichzeitig skalieren, ohne das No-Offset-Prinzip zu unterlaufen.
\end{quote}}



%% ----------------------------------------
\section[VP 5c: KBK -- Kulisse, Metriken, Top-Up \& Verfahren]{VP 5c: KBK -- Kulisse, Metriken, Top-Up \& Verfahren}
\noindent\textit{ID: 22 -- vo5c (meta)}

\begin{quote}
Parameterisierung des KBK-Top-Ups zum Carbon-Reward (Art. 5): Kulisse/Kategorien, Konnektivitäts- und Resilienzmetriken, Vergütungslogik, Permanenz/Sicherung, Anti-Gaming und Transparenz.
\end{quote}

- Kulisse \& Kategorien:
- Kartenwerk: Nationale KBK-Karte mit Kategorien: Fluss-/Auenkorridore, Moor-/Feuchtgebietsachsen, Wald-/Hecken-/Agroforst-Bänder, Urbane Blau-Grün-Korridore.

- Mindestbreiten: je Kategorie [VO]--[VO] m; städtisch abweichende Bandbreiten zulässig bei nachgewiesenem Resilienznutzen.

- Eligibility \& Metriken:
- Konnektivität (KC): Nachweis per anerkannter Metrik (z. B. Least-Cost-Path/Probability-of-Connectivity) ≥[VO] gegenüber Baseline.

- Resilienz-Score (R): Hitze-/Überflutungskühlleistung, Retentionsbeitrag, Verdunstung; Schwelle R ≥[VO].

- Do-No-Harm: Mindeststandards Arten/Habitate; Konfliktprüfung mit Schutzgebieten.

- Additionalität: Über geltende Pflichtstandards hinaus; keine Doppelzählung mit reinen Naturschutzfinanzierungen für identische Wirkung.

- Vergütungslogik (Top-Up):
- Variante 1 -- tCO\$\_2\$e-Zuschlag: σKBK €/t als Zuschlag auf die jeweils anerkannten tCO\$\_2\$e des Basis-Rewards; Cap [VO] €/Projekt·Jahr.

- Variante 2 -- Flächen-Annuität: Annuität je ha·Jahr in Dauerklassen D40/D100; Cap [VO] €/ha·Jahr; CO\$\_2\$-Mengen bleiben unverändert.

- Finanzierung: Transformationsförderung (Art. 3) und optional KKR-Top-Up (Art. 10c); keine Anrechnung auf produktseitige Claims (Art. 5 -- No-Offset).

- Permanenz \& Sicherung:
- Dauerklassen: D40/D100 mit Puffer 5--20 \% (risikobasiert); Reversal-Management analog Art. 5b.

- Rechtsbindung: Pflege-/Erhaltungsverpflichtungen (z. B. Grundbuch-Dienstbarkeit/Vertragsnaturschutz) über die Vergütungsdauer hinaus.

- Safe-Harbor \& Verfahren:
- Standard-Kits: Typkorridore (Heckenband, Auenfenster, Moor-Stepping-Stones, urbane Kühlachse) inkl. MRV-Bausteinen; Übernahme = vermutungsweise konform.

- Onboarding/Abrechnung: über Klima-Serviceplattform (Art. 8a); Registerfelder: Seriennummer, Vintage, Kategorie, KC/R-Werte, Dauerklasse.

- Anti-Gaming \& Schnittstellen:
- Keine Doppelförderung derselben Funktion (z. B. identische Biodiversitätsleistung) aus anderen Programmen; kombinierte Finanzierung zulässig bei trennscharfer Wirkung.

- Abgrenzung zu LULUCF-Belohnungen (Art. 5b): KBK vergütet nur zusätzliche Konnektivitäts/Resilienz-Co-Benefits.

- Transparenz \& Evaluation: Offene Karten/Parameter, Projektliste mit KC/R-Werten; 24-Monats-Evaluation (Konnektivitätsgewinn, Kühlleistung, Hochwasserpuffer, Kosten-Nutzen); Anpassung per VO.

- EU-/Naturschutzrecht: Kohärenz mit EU-Naturschutz-/Wasserrecht; erforderliche Notifizierungen nach Art. 17; Vergabe/Beihilfe nach Harmonisierungskapiteln.


%% ----------------------------------------
\section[Artikel 6: Strompreise \& Netzentgelte -- Leitplanken]{Artikel 6: Strompreise \& Netzentgelte -- Leitplanken}
\noindent\textit{ID: 23 -- artikel6 (gesetz)}

\begin{quote}
Elektrizität ist die zentrale Energieform der Transformation. Sie entscheidet über Wettbewerbsfähigkeit, bezahlbares Wohnen und den Erfolg der Dekarbonisierung. Dieses Gesetz schafft verlässliche, planbare und klimaverträgliche Strompreise: Verzerrende Umlagen werden geordnet abgebaut, Netzentgelte werden zeit- und leistungsvariabel, Flexibilität wird belohnt. Marktsignale bleiben erhalten; Preisspitzen werden gezielt gedämpft. Die operative Ausgestaltung erfolgt in Artikel 17.17 (EN-VO) und den dort benannten Verordnungen/Regelwerken.
\end{quote}

- Ziel \& Verhältnis. Dieser Artikel setzt die Leitplanken für die Preisbildung im Stromsektor, einschließlich der Wechselwirkungen mit den Carbon Rewards (Art. 5/5a). Technische Parameter, Datenpfade, Sicherungen und Fristen werden nach Art. 17.17 (EN-VO) festgelegt; Monitoring/Transparenz nach Art. 17.14 (AR-VO), Ident/Signaturen nach Art. 17.12 (ID-VO).

- Abgaben-/Umlagen-Neuzuschnitt.
- Konzessionsabgaben. Die verbrauchsgebundenen Konzessionsabgaben auf Strom laufen geordnet aus; die kommunale Beteiligung erfolgt künftig nach Maßgabe des Kommunalbonus (vgl. Art. 17.17 § 2 KAV-Neu). Übergangsregelungen regelt die EN-VO.

- EEG \& Bestandsförderungen. Neuanlagen erhalten keine EEG-Vergütung; Förderung erfolgt über Carbon Rewards (Art. 5/5a). Bestandsanlagen können zu festgelegten Stichtagen einmalig und irreversibel wechseln; Doppelförderung ist ausgeschlossen. Wechselpfade und Datenanschlüsse regelt die EN-VO.

- Besondere Umlagen. Offshore-Finanzierung und weitere Sonderumlagen (z. B. § 19 StromNEV, AbLaV) werden verursachergerecht neu geordnet und -- wo sachgerecht -- haushalterisch/auktionär finanziert; reine Endkundenabwälzung wird vermieden. Details in der EN-VO.

- Stromsteuer. Die Stromsteuer wird auf das unionsrechtliche Minimum abgesenkt. Wirkungs-/Verteilungsanalysen erfolgen im Ampel-Reporting (Art. 17.14).

- Netzentgelte \& Flexibilität. Netzentgelte werden zeit- und leistungsvariabel ausgestaltet (ToU/LP), mit standardisierten § 14a-Profilen (S/M/L) und Aggregator-Zugang. Engpass-/Zonenmodelle, Open-Data-Pflichten (Fenster/Zonen/Preise) und ein Preiskorridor zur Abfederung außergewöhnlicher Kostenanstiege werden in der EN-VO definiert. Doppelbelastungen werden vermieden.

- Dynamische Tarife \& Lieferantenpflichten. Für SMGW-fähige Kunden ist mindestens ein dynamischer Tarif (Spot-Index oder ToU) anzubieten; Preisbestandteile sind T-1 maschinenlesbar zu veröffentlichen. Fallback-Regeln, Standard-Preisblatt und Abrechnungsschema folgen der EN-VO.

- Temporäre Preisstabilisierung (USt-Stabilisator). Zur Dämpfung außergewöhnlicher Preisspitzen kann ein preisschwellenbasierter, zeitlich befristeter Umsatzsteuer-Stabilisator für Strom aktiviert werden, ausschließlich im Rahmen des Unionsrechts. Schwellen, Bandbreiten, Aktivierungslogik, Dauer und Rückkehr zum Normalsatz legt die EN-VO fest; jede Aktivierung wird im Ampel-Reporting transparent gemacht.

- Gleichbehandlung \& Beihilfe. Entlastungen gelten grundsätzlich symmetrisch für Haushalte und Unternehmen. Sektor-/standortspezifische Entlastungen sind nur zulässig, wenn sie zeitlich befristet, beihilfekonform (CEEAG/GBER) und klimapolitisch begründet sind (EN-VO).

- Transparenz \& Monitoring. Alle Preisbestandteile, Netzdaten und Kennzahlen (z. B. Redispatch-Kosten, Lastverschiebung, Speicherstunden) werden nach Art. 17.14 maschinenlesbar veröffentlicht (Dashboard, Open-Data). Outcome-KPIs werden jährlich berichtet; Ampel-Trigger können Anpassungen in EN-VO/Netzentgeltregulierung auslösen.

- Übergang \& Bestand. Umstellung in Stufen mit klaren Terminen; Bestandsrechte bleiben unberührt, soweit nicht wirksam optiert. Reservierungen/Übergangsbescheide nach EN-VO; Evaluation und ggf. Nachsteuerung binnen 6 Monaten nach Jahresbericht (Art. 17.14).

{\footnotesize\begin{quote}
Leitbild: Erneuerbare senken langfristig das Preisniveau; smarte Netzentgelte und dynamische Tarife sorgen für Effizienz; soziale/regionale Effekte werden außerhalb der Stromrechnung ausgeglichen; Klimawirkung wird über Carbon Rewards vergütet.
\end{quote}}



%% ----------------------------------------
\section[Artikel 7: CO\$\_2\$-Steuer (Framing als THG-Preisstabilisator auch möglich)]{Artikel 7: CO\$\_2\$-Steuer (Framing als THG-Preisstabilisator auch möglich)}
\noindent\textit{ID: 24 -- artikel7 (gesetz)}

\begin{quote}
2031 - Das Ende des Klimaschutzes: Ein Zusammenschluss aus Förderländern und Trittbrettfahrern versucht, Europa aus der Spur zu bringen. Erst werden Rohstoffpreise mit Ankündigungen und überraschenden Liefermengen verschoben, dann kursieren Gerüchte über angebliche Überschüsse im Emissionshandel. Parallel fallen in manchen Ländern zeitweise die Datenkanäle großer Unternehmen aus. Die Folge: Der CO2‑Preis rutscht für zwei Monate deutlich ab, springt kurz darauf auf ein Hoch. In Frankreich und Schweden geraten bekannte Industrieunternehmen ins Straucheln, Hedging scheitert, Aufträge brechen weg, Insolvenzmeldungen folgen. In Deutschland bleibt das Signal stabil, Projekte laufen weiter, Banken halten ihre Linien.
 Im Sommer wird in Brüssel eine wichtige Anpassung verschoben -- befeuert von öffentlichem Druck und klug gesetzten Botschaften. Mehr Zertifikate kommen kurzfristig in den Markt, der Preis fällt über drei Quartale. Förderinstitute frieren Zusagen ein, Planungen geraten ins Stocken. In Deutschland bleiben Mindestpreise und verbindliche Zusagen für Klimaleistung bestehen. Unternehmen rechnen mit festen Cashflows, Finanzierungen schließen sich trotzdem. Die Pipeline bleibt intakt, während anderswo die Ketten reißen.
 Im Herbst trifft eine Welle gezielter IT‑Störungen mehrere große Firmen jenseits der Grenze: Produktdaten, Nachhaltigkeitsberichte und Nachweise sind plötzlich nicht mehr vertrauenswürdig. Der Handel reagiert mit Nervosität, Sicherheiten werden erhöht, Preise zerren in beide Richtungen. In Deutschland greift die Vorsorge: Systeme laufen geordnet weiter, Zwischenmaßnahmen treten automatisch in Kraft, Zahlungen und Projekte bleiben verlässlich. Die Hektik flacht an der Landesgrenze ab.
 Als sich am Kap der Guten Hoffnung ein Engpass zuspitzt, wird daraus ein globales Drama erzählt. Spot‑Lieferungen werden reduziert, Gaspreise ziehen an, Strompreise springen in mehreren EU‑Ländern. Produktionslinien stehen, Notrufe nach Hilfspaketen übertönen die Nachrichten. In Deutschland fangen Speicher, flexible Lasten und klare Vergabe­regeln die Spitzen ab. Haushalte spüren den Ausschlag, aber nicht den Schlag; Betriebe sichern sich über feste Klimaleistungs‑Einnahmen und bleiben lieferfähig.
 Zum Jahresende setzt der Druck auf die nationale Steuerung ein: Institute und Kampagnen warnen vor „Standortverlust“ und fordern die Abschaffung fester Preisleitplanken. Mehrere Staaten lassen nationale Maßnahmen auslaufen, der CO2‑Preis rutscht, Exporte wanken. In Deutschland gilt der veröffentlichte Pfad weiter, automatische Gegenmaßnahmen greifen, der Grenzausgleich bleibt scharf. Unternehmen nutzen verlässliche Mindestpreise und garantierte Vergütungen, um Aufträge zu halten und Investitionen zu schließen. Während Europa ein „Schlussjahr“ für den Klimaschutz zu erleben scheint, wird Deutschland zum Stresstest‑Beleg dafür, dass klare Regeln, verlässliche Preise und planbare Klimaleistung den Unterschied machen.

  Ein klares Preissignal lenkt Investitionen -- verlässlich und planbar. Die CO\$\_2\$-Steuer ergänzt den EU-ETS mit einem nationalen Mindestpreis (Floor) und einem vorausgekündigten Preispfad. Das schützt vor Preisschwankungen, schafft Fairness über Sektoren und sichert Investitionen ab.
\end{quote}

- Steuertatbestand und Bemessung: Für jede emittierte Tonne CO\$\_2\$e aus (a) der Verbrennung fossiler Energieträger und (b) prozessbedingten Emissionen wird eine Steuer erhoben. Bemessungsgrundlage sind tCO\$\_2\$e nach Art. 2/2a (GWP100); methodische Updates wirken ambitionsneutral.

- Steuerschuldnerschaft und Entstehung:
- Upstream-Prinzip: Steuerschuldner sind die Inverkehrbringer der erfassten Brenn- und Kraftstoffe zum Zeitpunkt der steuerlichen Entnahme zum Verbrauch/Abgabe.

- Prozess-THG: Soweit Emissionen nicht upstream erfasst werden, ist Steuerschuldner der Anlagenbetreiber.

- Haftungskette: Gesamtschuld/Haftung nach Abgabenordnung; Nachweis-/Dokumentationspflichten folgen aus der VO.

- Mindestpreis \& Preiskorridor: Der Preisplan umfasst (a) eine verbindliche Festzone für die nächsten 5 Jahre, (b) einen Korridor (±20 \%) für weitere 5 Jahre sowie (c) einen unverbindlichen Ausblick. Der Floor ist nicht-sinkend.

- ETS-Interplay (Differenzprinzip): In Tätigkeiten, die dem EU-ETS/ETS2 unterliegen, wird auf Verbrennungs-CO\$\_2\$ nur die positive Differenz aus nationalem Mindestpreis minus maßgeblichem ETS-Preis erhoben. Liegt der ETS-Preis ≥ Floor, entfällt dieser Baustein. Prozess-Emissionen, die im ETS bepreist sind, sind nicht steuerbar.

- Backstop-Automatik (Budgetpfad): Weicht die projizierte Emissionsentwicklung vom Budgetpfad (Art. 2) ab, erfolgt eine regelbasierte Anpassung des Steuersatzes für Jahr N+1. Die VO legt Parameter (α/β/γ, Schwellen, Kappungen) fest. Bekanntgabe bis 1. Oktober; unterjährige Änderungen sind ausgeschlossen.

- Marktstabilisierung (national): Zur Dämpfung außergewöhnlicher Preisschocks kann ausschließlich für den steuerlichen Anteil eine befristete Glättung angewandt werden. Parlamentarische Kontrolle, volle Transparenz (Art. 17.14). EU-Instrumente (MSR) bleiben unberührt.

- Verwendung der Einnahmen: Einnahmen werden vorrangig verwendet: 90 \% Klimadividende (Art. 8), 10 \% Transformationsförderung (Art. 4). Abweichungen sind zulässig, werden jedoch jährlich begründet und veröffentlicht (Art. 17.14).

- Behörden \& Verfahren: Erhebung durch Zoll/BMF; Daten/Reporting/Backstop-Monitoring über die Transformationsagentur (Art. 8, 17.14). Abgabenordnung gilt entsprechend.

- Transparenz \& Open Data: Preispläne, Einnahmen/Verwendung, Budgetabweichung und Backstop-Status werden jährlich veröffentlicht (Art. 17.14).

- Verordnungsermächtigung: Die Bundesregierung wird ermächtigt, durch VO/VwV festzulegen: Nachweis-/Erhebungsverfahren, Befreiungen/Erstattungen, Backstop-Parameter, Publizitäts-/Datenstandards, BEHG/ETS-Schnittstellen, Bagatellgrenzen sowie IT-/Signaturvorgaben (Art. 17.12/17.14).


%% ----------------------------------------
\section[VO 7a: Risikoaufschlag für Vorketten-THG]{VO 7a: Risikoaufschlag für Vorketten-THG}
\noindent\textit{ID: 25 -- vo7a (meta)}

\begin{quote}
Ein Teil der Klimawirkung entsteht vorkettig (CH₄/N\$\_2\$O). Der Risikoaufschlag sorgt für ein vollständiges Preissignal und kann bei sauberen Lieferketten abgesenkt werden. ToDo: "s wäre sinnvoll, explizit zu sagen, dass die Faktoren R und R* auf Basis von best available science (z. B. IPCC-Emissionsfaktoren, peer-reviewte LCA-Datenbanken) und regelmäßigen Reviews kalibriert werden müssen."
\end{quote}

- Anwendungsbereich: Für definierte Brennstoffe/Routen werden Default-Risikofaktoren Rf,k angesetzt, soweit dieselben Mengen nicht bereits verbindlich bepreist sind.

- Bemessung: tCO\$\_2\$e\_steuer = tCO\$\_2\$,Verbrennung × (1 + Rf,k). Der Risikoaufschlag gilt sektorenübergreifend und unabhängig vom ETS-Preis.

- Individuelle Absenkung: Prospektiv möglich bei mess-/prüfbaren Nachweisen für geringere Vorketten-Emissionen (R*). Fehlangaben → Rückforderung + Malus.

- CBAM-Zertifikat: Bescheinigung „carbon price effectively paid“ weist den Risikoaufschlag separat aus (Art. 9a).

- Transparenz/Review: Jährliche Aktualisierung der Default-Listen, Quellen/Unsicherheiten, Konsultation; Open-Data-Veröffentlichung (Art. 17.14).


%% ----------------------------------------
\section[VO 7b: Durchführungsregeln CO\$\_2\$-Steuer]{VO 7b: Durchführungsregeln CO\$\_2\$-Steuer}
\noindent\textit{ID: 26 -- vo7b (meta)}

- Steuerschuldner \& Entstehung: Inverkehrbringer; Entstehung bei Überführung in den freien Verkehr. Schnittstellen zu Energie-/BEHG-Registern werden genutzt (Art. 17).

- Bemessungsgrundlage: Standardfaktoren Verbrennungs-CO\$\_2\$; Mess-/Schätzhierarchie; Prozess-CO\$\_2\$, soweit nicht ETS-bepreist.

- Risikofaktoren: Listenführung/Update; Verfahren zur Individualabsenkung (Geltungsdauer, Audit-Pflichten).

- Ausnahmen/Gutschriften: Re-Export, internationaler Bunker; dauerhafte Abscheidung/Bindung (CCS/CCU) anrechenbar; biogener Anteil steuerfrei bei RED-Nachweis.

- ETS/ETS2-Schnittstelle: Differenzprinzip nur für Verbrennungs-CO\$\_2\$; Prozess-ETS nicht steuerbar; Risikoaufschlag zusätzlich.

- Compliance-Kalender: Vorauszahlungen, Jahreserklärung, Sicherheiten, Verzugszinsen, Sanktionen; Bagatellgrenzen [VO].

- Open Data \& IT: Veröffentlichung, Versionierung; digitale Verfahren, Signaturen (Art. 17.12/17.14).


%% ----------------------------------------
\section[Artikel 8: Klimadividende]{Artikel 8: Klimadividende}
\noindent\textit{ID: 27 -- artikel8 (gesetz)}

\begin{quote}
Die Klimadividende macht den CO\$\_2\$-Preis sozial fair: Alle dauerhaft in Deutschland lebenden Menschen erhalten den gleichen Betrag pro Kopf -- einfach, digital, regelmäßig. Ziel der Klimadividende ist es, die Klimadividende regressiv wirkende Effekte der THG-Bepreisung zu kompensieren.
\end{quote}

- Zweck \& Prinzip: Pro-Kopf-Rückverteilung klimapolitischer Einnahmen; jede anspruchsberechtigte Person erhält den gleichen Jahresbetrag.

- Finanzierung: Vorrangige Zuführung der Einnahmen aus Art. 7; ergänzende Mittel gemäß Harmonisierung (Art. 17). Quoten/Verwendung siehe Art. 7.

- Anspruch: Personen mit dauerhaftem gewöhnlichen Aufenthalt im Bundesgebiet; Näheres in Art. 8a (VO/VwV).

- Bemessung: Jahresbetrag = Poolmittel / Anspruchsberechtigte abzgl. Glättungsrücklage. Mechanismus in Art. 8a.

- Auszahlung: Digital/barrierefrei in regelmäßigen Intervallen; Details in Art. 8a.

- Minderjährige: Zahlung auf ein zweckgebundenes Zukunftskonto mit definierten zulässigen Verwendungen; Details Art. 8a.

- Schutz: Unpfändbar, nicht abtretbar, nicht anrechenbar auf bedarfsgeprüfte Sozialleistungen; Fachgesetze werden harmonisiert (Art. 17).

- Transparenz: Jährlicher Bericht zu Einnahmen, Auszahlungen, Reserve, Empfängerkreisen; Open Data (Art. 17.14).


%% ----------------------------------------
\section[VO 8a: Durchführungsregeln Klimadividende]{VO 8a: Durchführungsregeln Klimadividende}
\noindent\textit{ID: 28 -- vo8a (meta)}

- Anspruchsprüfung: Gewöhnlicher Aufenthalt (Melderegister/äquivalente Nachweise); Teiljahre zeitanteilig.

- Registrierung \& Ident: Klima-Portal mit eID/bundID; analoge Verfahren; Vertreterregeln; Steuer-ID als Primärschlüssel.

- Auszahlung: Standard quartalsweise; optional monatlich. Kanäle: IBAN, Zahlungskarte, Basiskonto -- gebührenfrei.

- Glättung \& Startbonus: Dividendenpool mit Rücklage; Startbonus möglich; Auszahlungsformel/Trigger in VO.

- Zukunftskonto Minderjährige: Standardanlage in klimawirksame, risikoarme Instrumente; definierte vorzeitige Entnahmen; Kosten-Caps; Produktkriterien.

- Steuer-/Sozialrecht: Einkommensteuerfrei; keine Anrechnung auf SGB II/XII/BAföG/UVG; Harmonisierung per Art. 17.

- Fehlbeträge/Überschüsse: Rücklage hoch/runter; kein Nachschusszwang; Nachmeldungen bis [VO] Jahre; Bagatellgrenzen.

- Missbrauchsschutz: Rückforderung, Bußgeld, ggf. Straftat; eng zweckgebundene Datenabgleiche.

- Open Data \& IT: Kennzahlen, Versionierung; EGovG/NIS2; Barrierefreiheit.


%% ----------------------------------------
\section[Artikel 9: CO\$\_2\$-Grenzausgleich]{Artikel 9: CO\$\_2\$-Grenzausgleich}
\noindent\textit{ID: 29 -- artikel9 (gesetz)}

\begin{quote}
Fairer Wettbewerb braucht faire Klimakosten. Diese Regel ergänzt den EU-CBAM nur dort, wo Lücken bestehen, spiegelt das inländische CO\$\_2\$-Preissystem (inkl. Risikoaufschlag) und wahrt WTO-Grundsätze. Einnahmen fließen vollständig in die Klimadividende.
\end{quote}

- Vorrang EU-CBAM \& Zweck: Der EU-CBAM hat Vorrang. Das nationale CBAM-Komplement gilt nur für Waren/Sektoren/Emissionsbestandteile, die (noch) nicht vom EU-CBAM erfasst sind, um Inlands- und Importwaren beim CO\$\_2\$-Preis gleichzustellen.

- Rechtsnatur \& Gleichbehandlung: Das CBAM-Komplement ist eine inländische Verbrauchsteuer auf eingebettete Emissionen beim Inverkehrbringen. Gleichartige inländische und importierte Waren werden gleich behandelt (Art. 110 AEUV).

- Steuertatbestand \& Emissionsumfang: Steuergegenstand ist die eingebettete Treibhausgasmenge (GWP100) der eingeführten Ware p. Ein Risikofaktor (Art. 7a-Spiegel) berücksichtigt Vorketten-Emissionen. Details in Art. 9a.

- Bemessung \& Anrechnung: Abgabe je Sendung = max(0, Floory × Esteuer − effektiv gezahlter ausländischer Carbon-Preis). Anrechenbar sind prüfbare ausländische Preise (Steuern/ETS/Pflichtabgaben) für dieselbe Emissionsmenge. Doppelbelastung ausgeschlossen.

- WTO-/EU-Konformität \& Sunset: Technologieneutral; Notifizierung/Anpassung an EU-/WTO-Recht fortlaufend. Soweit EU-CBAM neue Waren/Emissionen abdeckt, entfallen überlappende nationale Elemente binnen 12 Monaten (Mapping-Liste, Art. 9a).

- Transparenz \& Rechtsschutz: Import-Register, Jahresbericht (Mengen, Default-Nutzung, Anrechnungen, Umgehung); offene Datenschnittstellen (Art. 17.14); Rechtsbehelfe mit fachlicher Anhörung.

- Mittelverwendung: Einnahmen fließen vollständig in den Dividendenpool (Art. 8).


%% ----------------------------------------
\section[VO 9a: Durchführungsregeln CO\$\_2\$-Grenzausgleich]{VO 9a: Durchführungsregeln CO\$\_2\$-Grenzausgleich}
\noindent\textit{ID: 30 -- vo9a (meta)}

- Nachweis \& MRV: Emissionswerte nach PCF/MRV (ISO 14067, GHG Product Standard, EN-Normen); Zulassung unabhängiger Prüfer; Fristen; Datenformate (API).

- Default-Werte \& Zonen: Konservative Default-Werte je HS-Position (BAC/BAT-Quartile) und behördlich veröffentlichte Strom-EF je Zone (u. a. für Strom/Wasserstoff/eFuels). Jährliche Aktualisierung.

- Risikofaktoren: Spiegel der Art. 7a-Listen (Lieferketten/Routen); individuelle Absenkung prospektiv bei prüfbaren Nachweisen; Übergangswerte befristet möglich.

- Verfahren \& Abrechnung: Registrierung vor Einfuhr, elektronische Zoll-Anmeldung mit Emissionsdaten/Belegen, Bescheid/Anrechnung EU-CBAM und ausländischer Carbon-Preise, Zahlung/Sicherheiten, Stundung.

- De-minimis \& KMU: Schwellen und vereinfachte Verfahren bei gewahrter Klimawirkung.

- Äquivalenzanerkennung: Verfahren zur Anerkennung ausländischer Systeme (Breite, Preisniveau, MRV-Qualität, Vorketten-Abdeckung) mit Wirkung auf Anrechnung/Default-Wahl.

- Anti-Umgehung: Tests für geringfügige Verarbeitung/Tarifwechsel/Re-Invoice, Ursprungs-/Wertschöpfungsprüfungen; Sanktionen/Sperren.

- Sunset-Mapping: Fortschreibung der Liste, welche nationalen Elemente bei EU-Abdeckung binnen 12 Monaten entfallen.

- Transparenz \& IT: Offene Tabellen/Parameter, Jahresbericht; Schnittstellen zu CBAM/Zoll-IT; Sicherheit nach Art. 17.12/17.14.


%% ----------------------------------------
\section[Artikel 10: Institutionen]{Artikel 10: Institutionen}
\noindent\textit{ID: 31 -- artikel10 (gesetz)}

\begin{quote}
Deutschland bekommt für die Klimatransformation drei klare Rollen: umsetzen, steuern, entscheiden. Die Klima-Transformationsagentur ist der One-Stop-Vollzug (Kredit, Carbon Rewards, CO\$\_2\$-Preis, Dividende, Grenzausgleich, Daten). Die Wissenschaftliche Klimakommission setzt Leitplanken nach Stand der Forschung (Ziele, Budget, Floors) und aktiviert Backstops bei Zielverfehlung. Eine spezialisierte Klimagerichtsbarkeit sorgt für zügigen Rechtsschutz und klare Auslegung. Alles ist mit den Harmonisierungsvorschriften (Art. 17 ff.) verdrahtet -- doppelte Zuständigkeiten werden abgebaut, Daten laufen „once-only“.
\end{quote}

- Klima-Transformationsagentur (Vollzug \& Service):
- Mandat/One-Stop: Zentrale Vollzugsbehörde für Transformationskredite, Innovationsförderung, Carbon Rewards, CO\$\_2\$-Steuer/Dividende, nationalen CO\$\_2\$-Grenzausgleich, Monitoring/Open Data.

- Harmonisierungsvollzug: Umsetzung der in Art. 17 ff. übertragenen Aufgaben; bestehende Parallelverfahren werden geordnet überführt.

- Digital \& Fristen: eID, Standardkataloge, APIs; Standardfall ≤ 30 Tage, komplex ≤ 60 Tage; Genehmigungsfiktion für Standardfälle (Ausnahmen Hochrisiko).

- Daten \& Transparenz: Datenführung ohne Monopol, Open-Data (schutzwürdige Daten ausgenommen); monatliche Durchlauf-/Wirkungsdashboards.

- Kooperation: Schnittstellenabkommen mit DEHSt/UBA, BZSt/Zoll, BNetzA sowie Ländern/Kommunen; Fachzuständigkeiten anderer Behörden bleiben unberührt.

- Wissenschaftliche Klimakommission (Steuerung \& Backstop):
- Aufgaben: Vorschlag nationaler Ziele/Budget/Jahresobergrenzen (Art. 1), Monitoring und Krisenfeststellungen; Preisleitplanken/Mindestpreise für Carbon Rewards (Art. 4) im Rahmen parlamentarischer Leitplanken; jährliche Berichte.

- Mandat (Prüfe ob hier richtig platziert). Die Kommission bewertet regelmäßig, ob der Anteil der Restemissionen und Removals im Einklang mit dem Stand der Wissenschaft zu CDR-Risiken steht und kann Leitplanken vorschlagen.

- Konsistenzprüfung (Prüfen ob hier richtig platziert). Neue fossile Großvorhaben und wesentliche Lebensdauerverlängerungen bestehender fossiler Anlagen setzen eine vorherige Stellungnahme der Klimakommission zur Budgetkompatibilität voraus.

- Backstop: Feststellung von Pfadabweichungen und Backstop-Verordnung; Wirksamkeit, wenn der Bundestag nicht binnen 3 Monaten gleichwirksame Alternativen beschließt.

- Integrität: Unabhängige Berufung, Amtszeit, Befangenheit/Karenz; Veröffentlichung von Modellen, Daten, Minderheitsvoten.

- Klimagerichtsbarkeit (Schutz \& Klarheit):
- Zuständigkeit: Vorrangig für Streitigkeiten aus dem KlimaGG, einschließlich Preis/Steuer/Reward/Grenzausgleich; Grundrechtsschutz und Auslegung.

- Schnellverfahren: Eilrechtsschutz mit Ziel 6 Wochen, Musterverfahren; offene Justizdaten.

- Föderale Einbindung:
- Länderbeirat (konsultativ): Vor Rechtsverordnungen, Safe-Harbor-Mustern und NDZ-Standardverträgen ist der Länderbeirat anzuhören; Stellungnahmen werden veröffentlicht. Kein Vetorecht.

- Bund-Länder-Verwaltungsvereinbarung: Gemeinsame Nutzung von Mustern/Vollzugshilfen; in Landesvollzügen gelten sie als antizipierte Regeln der Verwaltung -- Übernahme bleibt freiwillig.

- Abgrenzung \& Zusammenarbeit: Vollzug/Datenführung bei der Agentur; Zielsetzung/Backstop bei der Kommission; Kontrolle bei der Gerichtsbarkeit. Koordination über geregelte Schnittstellen, gemeinsame Datenstandards und jährliche Berichte.

{\footnotesize\begin{quote}
To-dos Übergang: (1) Schnittstellenvereinbarungen (DEHSt/BEHG, BZSt/Zoll, BNetzA), (2) Harmonisierungskatalog (Art. 17) live, (3) Überleitung EEG/KWKG/EnFG/KAV inkl. IT/Personal, (4) gemeinsamer Jahresbericht zur Anrechnung BEHG/TEHG und Monitoring-Harmonisierung.
\end{quote}}



%% ----------------------------------------
\section[VO 10a: Klima-Serviceplattform \& Safe-Harbor]{VO 10a: Klima-Serviceplattform \& Safe-Harbor}
\noindent\textit{ID: 32 -- vo10a (meta)}

\begin{quote}
Kritische Projekte scheitern oft an Personalengpässen und Komplexität. Die Klima-Serviceplattform macht aus vielen Einzelfragen einen einfachen Pfad: Einsatzpool vor Ort, geprüfte Muster mit Safe-Harbor-Wirkung, technischer Baukasten (Daten, MRV, Rahmenverträge) und klare Fristen -- ohne neue Verfahren.
\end{quote}

- Grundsätze \& Kein-Mehrverfahren: Desks für Kommunen, Bürger:innen, Unternehmen (inkl. KMU). Anforderungen werden in bestehende Verfahren integriert; keine Zusatz-Genehmigungsschritte. Keine neuen Pflichtaufgaben für Kommunen.

- Einsatzpool: Bundesweiter Personalpool (Vergabe/Beihilfe, MRV/IT, Bau/Technik, Wärmeplanung, NDZ/Flex, Partizipation) mit befristeter Entsendung (max. 24 Monate) und Fallmanagement („Projekt-Nurse“). Haftung/Compliance bei der Agentur.

- Muster mit Safe-Harbor: Muster-Satzungen, städtebauliche Verträge, Vergabe-Bausteine (LCC/CO\$\_2\$), Beihilfe-Steckbriefe (CEEAG/GBER), Standard-MRV-Kits. Unveränderte Übernahme = vermutungsweise rechtssicher/beihilfekonform. Good-Faith-Milderung bei Formfehlern.

- Technischer Baukasten \& Register: Standardisierte Datenpakete (EF\_avg/EF\_marg, Klima-Risiko-Layer, Baselines) als GIS/API; Plug-\&-Play-MRV; Beitritts-Rahmenverträge (PV/Speicher/WP, Holz-Modulbau, Dämmung, Retention; NDZ-Sammelverträge). Produkt-/Gebäudepässe werden mit dem Register gemappt.

- Service \& SLAs: 48-Stunden-Kurzgutachten; 15 Arbeitstage bis Standard-Entscheidung Kredit/CR-Onboarding; 30 Tage bis Einsatzbeginn; 90 Tage bis „Start of Works“ für Standardmaßnahmen. Clearingstelle Kommunal (30 Tage).

- Finanzierung \& Anreize: Basisfinanzierung mind. 5 \% des TFK-Wirtschaftsplans; Enablement-Pauschale bis 3 \% förderfähige Kosten; optional Tempo-Bonus bei SLA-Einhaltung [VO].

- NDZ-Standardisierung: Standardverträge mit VNB (Aggregations-/Abrufregeln, Vergütung, Bilanzkreis-/Haftung); Kommunen können als Bündel-Aggregator auftreten oder die Agentur beauftragen.

- Once-Only \& KMU-Fit: Bereits eingereichte Daten dürfen nicht erneut verlangt werden; Länderportale via Schnittstelle. Rahmenverträge losweise; Zielquoten/Losobergrenzen [VO].

- Monitoring \& Sunset: Jährliche KPIs (Durchlaufzeiten, Anteil Rahmenverträge, Soft-Cost-Quote, zusätzliche CR-Mengen, Zufriedenheit). Muster/SLAs/Standardwerke: Evaluation alle 24 Monate; ohne Verlängerung Sunset nach 36 Monaten.

- IT-Sicherheit \& Datenschutz: DSGVO-konforme Verarbeitung; offene APIs; staatliche IT-Sicherheitsstandards.

- Übergang: Erste Muster-/Rahmenwerks-Edition binnen 6 Monaten; Einsatzpool organisatorisch einsatzfähig binnen 90 Tagen, erste Entsendungen binnen 120 Tagen.


%% ----------------------------------------
\section[VO 10b: Beschleunigungsrahmen Genehmigungen]{VO 10b: Beschleunigungsrahmen Genehmigungen}
\noindent\textit{ID: 33 -- artikel10b (meta)}

\begin{quote}
Einheitliche, schnelle Verfahren für Vorhaben nach diesem Gesetz: Kategorien, Typauflagen, harte Fristen mit enger Stop-the-clock, Fiktion in Standardfällen und klare Eskalation -- Aarhus-konform und föderal fair.
\end{quote}

- Kategorien \& Geltung: Per VO-Einstufung in C (Anzeige + Typauflagen), B (vereinfachte Genehmigung) und A (komplex mit Konzentrationswirkung); digitale Verfahrensführung verpflichtend.

- Typauflagen: Immissions-/Artenschutz, Lärm/Blend, Baustellenlogistik, Mess-/MRV-Pflichten; Gleichsachverhalte werden harmonisiert.

- Fristen \& Fiktionen: Vollständigkeitsprüfung C/B 14 AT, A 30 T. Entscheidungsfristen: C 30 T (Fiktion), B 4 Monate (Fiktion light), A 6 Monate. Stop-the-clock nur für fehlende Unterlagen/EU-Beteiligungen, max. 30 + 30 T.

- Konzentrationswirkung \& Eskalation (A): Federführung durch Landesbehörde mit Konzentrationswirkung; Mitwirkungsbehörden 6 Wochen, ausbleibend = keine Einwände. Bei Verzug Task-Force/Übernahme; Fristen-Dashboard (Open Data).

- Repowering: Prüfumfang auf neue/erhöhte Auswirkungen beschränkt; standardisierte Artenschutzpakete zulässig.

- Beteiligung \& Rechtsschutz: Digitale + analoge Auslegung (mind. 1 Monat), hybrider Erörterungstermin; gerichtliche Konzentration/Bündelung gleichgelagerter Klagen nach Begleitänderungen.

- Verordnungsermächtigung: Bundesregierung legt Kategorien, Schwellen, Typauflagen, Standard-Nebenbestimmungen, Fristen, Formate fest; EU-Beihilfe-Gate, Sunset/Review und Kein-Mehrverfahren gelten entsprechend.


%% ----------------------------------------
\section[Artikel 11: Bund-Länder-Pakt Klimatransformation]{Artikel 11: Bund-Länder-Pakt Klimatransformation}
\noindent\textit{ID: 34 -- artikel11 (gesetz)}

\begin{quote}
Der Bund-Länder-Pakt bündelt Kapazitäten und sorgt für einheitlichen, schnellen Vollzug der Klimatransformation -- ohne neue Pflichtaufgaben für Länder/Kommunen. Er knüpft an die Harmonisierungsvorschriften in Artikel 17 ff. an (Once-Only, Open-Data, IT-Sicherheit, Monitoring).
\end{quote}

- Zweck \& Prinzipien. Der Pakt dient der Beschleunigung, Entlastung und Vereinheitlichung im föderalen Vollzug. Er respektiert die Zuständigkeiten der Länder; er arbeitet mit Anreizen (Safe-Harbor, Pauschalen, Boni), nicht mit Zwang.

- Rechtsform \& Abschluss. Die Bundesregierung wird ermächtigt, mit den Ländern eine Verwaltungsvereinbarung zu schließen. Mindestinhalte nach Abs. 3--9 sind verbindlich; Details werden in der Vereinbarung und in nachgeordneten Verwaltungsvorschriften präzisiert.

- Safe-Harbor-Muster \& Standards. Übernahme der Muster und Standardwerke aus Art. 8a/8b und Art. 17 ff. (z. B. LCC/CO\$\_2\$-Vergabebausteine, MRV-Kits, NDZ-Rahmenverträge). Unveränderte Anwendung entfaltet eine widerlegliche Vermutung der Recht-/Beihilfekonformität (Safe-Harbor).

- Personal-/Einsatzpool. Einrichtung eines gemeinsamen Expertenpools (Vergabe/Beihilfe, Genehmigung, MRV/IT, Wärmeplanung, NDZ/Flex). Entsendungen sind befristet; Einsatz nach einheitlichen SLAs. Finanzierung und Kostenerstattung nach Abs. 7.

- Fristen \& SLAs. Länder/Kommunen, die am Pakt teilnehmen, wenden die Fristenlogik aus Art. 8b an (Vollständigkeits- und Entscheidungsfristen, enge Stop-the-Clock). Ein öffentliches Fristen-Dashboard wird über Art. 17.14 (AR-VO) gespeist.

- Once-Only, Daten \& IT. Verbindliches Daten-/Schnittstellenprofil gemäß Art. 17.12 (ID-VO) und 17.14 (AR-VO): eID/bundID, offene APIs, QES/QSeal, einheitliche Datenmodelle/Versionierung. Bereits übermittelte Daten dürfen nicht erneut angefordert werden (Once-Only).

- Mitfinanzierung \& Pauschalen. Der Bund leistet pauschale Beiträge zu Vollzugskosten und Einsätzen aus dem Wirtschaftsplan des TFK (Art. 3) bzw. Haushaltsmitteln. Die Vereinbarung legt quotal fest: Einsatzpauschalen, Projekt-Enablement-Pauschalen (z. B. bis 3 \%), sowie optionale Bonuszahlungen bei SLA-Einhaltung. Kommunalbuchung: nach Art. 17.7 (KH).

- Monitoring \& Transparenz. Einheitliche KPIs (Durchlaufzeiten, Genehmigungsquoten, Soft-Cost-Anteil, zusätzliche CR-Mengen) werden vierteljährlich als Open-Data gemäß Art. 17.14 veröffentlicht; Ampel-Trigger bei Abweichungen.

- Beteiligung \& Länderbeirat. Ein Länderbeirat wird vor Erlass/Änderung einschlägiger Verordnungen und Muster (Art. 8a/8b, 17.12, 17.17, 17.14) angehört; Stellungnahmen werden veröffentlicht. Kein Vetorecht.

- Streitbeilegung. Schlichtungsgremium (Bund/Länder, unabhängige Fachperson) mit 30-Tage-Frist; Ergebnisse werden dokumentiert. Rechtsschutz im Übrigen nach den allgemeinen Regeln.

- Sunset \& Review. Der Pakt wird jährlich evaluiert (Art. 17.14). Ohne Erneuerung endet er nach 36 Monaten (Sunset); Anpassungen sind im Einvernehmen möglich.

{\footnotesize\begin{quote}
- Kohärenz: Dieser Artikel verweist für Technik/IT/Monitoring auf Art. 17.12 (ID-VO), 17.17 (EN-VO) und 17.14 (AR-VO); doppelte Detailregeln werden vermieden.

- Anreize statt Zwang: Teilnahme bringt Safe-Harbor, Einsatzzugriff und Pauschalen; Nichtteilnahme bleibt ohne Sanktion.
\end{quote}}



%% ----------------------------------------
\section[Artikel 12: Strukturwandel, Weiterbildung und Austausch]{Artikel 12: Strukturwandel, Weiterbildung und Austausch}
\noindent\textit{ID: 35 -- artikel12 (gesetz)}

\begin{quote}
Strukturwandel soll Chancen öffnen statt Brüche vertiefen: verlässliche Qualifizierung, regionale Transformationsverträge, einfache Zugänge für KMU -- verzahnt mit Transformationskredit (Art. 3) und Transformationsförderung (Art. 4).
\end{quote}

- Frühwarn- \& Beteiligungspflichten. Geplante Stilllegungen oder wesentliche Transformationsschritte werden frühzeitig angezeigt; Beschäftigte und Sozialpartner sind einzubeziehen. Verfahrenswege, Schnittstellen und Datenformate richten sich nach Art. 17 ff.

- Rechtsanspruch auf Qualifizierung \& Vermittlung. Anspruchsberechtigte erhalten Zugang zu Um- und Weiterqualifizierung mit Ziel der zügigen Neuplatzierung. Umsetzung durch die Klima-Transformationsagentur (Art. 10) in Kooperation mit BA und Ländern; Primärfinanzierung aus dem TFK (Art. 3); Prozess/IT nach Art. 17 ff.

- Qualitätskennzahlen. Verbindliche KPIs (u. a. Abschlussquoten, Zeit bis Neuplatzierung, Lohnentwicklung) werden festgelegt; Messmethodik, Veröffentlichung als Open Data und Versionierung gemäß Art. 17.14.

- Just-Transition-Regionen. Transformationsverträge zwischen Agentur und Regionen (Ziele, Meilensteine, Mittelbündel aus Art. 3--5, Infrastrukturprioritäten). Standardmuster, Datenschnittstellen, Berichtswesen nach Art. 17 ff.

- KMU-/Handwerkszugang. Vereinfachte Nachweise und priorisierte Beratung; die Service-Level (z. B. 30-Tage-SLA) gelten entsprechend Art. 10. Safe-Harbor-Muster und Standardkataloge nach Art. 17 ff.

- Internationaler Austausch. „Klima-ERASMUS“ für Schule, Ausbildung, Studium, berufliche Weiterbildung; Kopplung an Innovationscluster/Reallabore (Art. 4). Programmrichtlinien und Beihilfenachweise nach Art. 17 ff.

- Keine Bailouts. Kein Entschädigungsanspruch für Kapitalverluste aus der Umsetzung dieses Gesetzes; keine Schuldenübernahme/Fehlinvestitionsrettung durch die öffentliche Hand.

- Öffentliche Beteiligungen. Erhöhte Transparenz; keine Zuschüsse für nicht transformierbare Geschäftsmodelle; öffentliche Beteiligungen werden dekarbonisiert im Einklang mit Art. 14.

- Vergaben \& TCO. Öffentliche Investitionen sind an Gesamtkosten (TCO) einschließlich CO\$\_2\$-Wirkung auszurichten; Bewertungsmethodik, Datenquellen und Dokumentation nach Art. 17 ff.

- Berichtspflichten. Jahresbericht zu sozialen/ökonomischen Wirkungen; Publikation als Open Data gemäß Art. 17.14.


%% ----------------------------------------
\section[Artikel 13: Fossil-to-Climate Shift (FCS) -- Umwidmung fossiler Begünstigungen]{Artikel 13: Fossil-to-Climate Shift (FCS) -- Umwidmung fossiler Begünstigungen}
\noindent\textit{ID: 36 -- artikel13 (gesetz)}

\begin{quote}
Schrumpfende fossile Begünstigungen werden automatisch in Klimamittel verwandelt: Was nicht mehr als Vorteil gewährt wird, fließt planbar in den Transformationsfonds Klima (TFK) -- ohne neue Bürokratie.
\end{quote}

- Begriffe \& Abgrenzung. Einbezogen sind im Bundeshaushalt ausgebrachte Vergünstigungen für Diesel/Kerosin u. a. gemäß Anlage. Subventionsäquivalent nach BMF-Methodik; Referenzbetrag (R) = höchster tatsächlich realisierter Jahresgesamtwert der letzten fünf Haushaltsjahre. Ermittlung/Datengrundlagen: Art. 17 ff.

- Obergrenze \& Kürzungsautomatik. Der Jahresgesamtwert darf R nicht überschreiten. Droht Überschreitung, lineare Kürzung per Rechtsverordnung mit Zustimmung des Bundestages. Haushalts-/Abgabenharmonisierung nach Art. 17 ff.

- Automatische Zuführung an den TFK. Z = R − B (mind. 0) wird als gebundene Ausgabe monatlich an den TFK (Art. 3) abgeführt; Endabrechnung bis 30.06. des Folgejahres. IT-Vollzug und Open-Data-Bericht nach Art. 17.14.

- Methodik \& Transparenz. Das BMF legt die Ermittlung per VO fest; quartalsweise Prognosen und jährlicher FCS-Bericht (R, B, Z, etwaige Kürzungsfaktoren). Format/Versionierung gemäß Art. 17.14.

- Haushalts- \& Beihilferecht. Zweckgebundene, formelgebundene Zuführung zum Sondervermögen TFK; Mittelverwendung beihilfekonform. Kohärenz-/Notifizierungspflichten nach Art. 17 ff.

- Übergang. Zuführung mind. 80 \%/90 \% in Jahr 1/2; ab Jahr 3: 100 \%. Fortschreibung von R alle drei Jahre.

- Kerosin-Sonderlage. Internationale Steuerbefreiungen bleiben unberührt; sinken Verkehrsmenge oder Begünstigungsumfang (z. B. EU-Reformen/SAF-Quoten), steigt Z automatisch.


%% ----------------------------------------
\section[Artikel 14: Divestment]{Artikel 14: Divestment}
\noindent\textit{ID: 37 -- artikel14 (gesetz)}

\begin{quote}
Öffentliche Portfolios verfolgen ein Zero-Exposure-Prinzip gegenüber fossilen Aktivitäten und richten Neuanlagen auf klimawirksame Investitionen aus. Bei staatlich kontrollierten Unternehmen tritt an die Stelle eines Abstoßens die steuernde Dekarbonisierung über verbindliche Transformationsfahrpläne (TFP). Operative Screens/Look-through/Reporting regelt die VO 14a; TFP-Details die TFP-VO. EU-Kohärenz bleibt gewahrt.
\end{quote}

- Geltung \& Begriffe
- Erfasst sind Portfolios und Sondervermögen von Bund/Ländern/Kommunen einschließlich Pensions-/Versorgungswerken sowie öffentliche Unternehmen, soweit keine staatliche Kontrolle besteht.

- Fossile Aktivitäten im Sinne dieses Artikels sind die in der VO 14a (Anlagen FOSS-A/FOSS-S) gelisteten Tätigkeiten und gekoppelten Services (NACE/NAICS-Mapping).

- Zero-Exposure für öffentliche Portfolios
- Neuanlagen in Instrumente/Emittenten mit fossilen Aktivitäten sind ab Inkrafttreten unzulässig.

- Run-off/Verkauf: bestehende Exposures werden geordnet ohne Re-Invest abgebaut; Fristen gemäß Gesetz/VO.

- Look-through: Fonds/ETFs sind durchgängig zu durchleuchten; ex-Fossil-Indizes binnen [VO 6--12] Monaten.

- Eng begrenzte Übergänge/Sicherungen bedürfen Begründungsbeschluss, Auslaufplan und Veröffentlichung (VO 14a).

- Staatlich kontrollierte Unternehmen -- Transformationsfahrplan (TFP)
- Besteht beherrschender Einfluss der öffentlichen Hand (insb. ≥ 50 \% Stimmrechte oder gleichwertige Kontrolle), gilt statt Divestment ein verbindlicher TFP mit Zielwerten: −70 \% bis 2030, −90 \% bis 2033, −100 \% bis 2035 (Scope 1+2, intensity-bereinigt; Residualmengen nur mit Linie-B-D100-Dauerhaftigkeit).

- CAPEX-Gates: ≥ [VO 85 \%] jährliche Sachinvestitionen klimakompatibel; Fossil-Ersatz nur bei Sicherheits-/Systemkritik mit Auslaufplan.

- Governance \& Vergütung: Vorstands-/Geschäftsführerziele an TFP-Meilensteine gekoppelt; variabler Anteil nur bei Pfaderfüllung.

- Transparenz \& Prüfung: jährlicher TFP-Statusbericht (Open Data), Assurance; Schnittstellen nach Art. 17.14.

- Fallback: Verfehlt das Unternehmen wesentliche Meilensteine ohne Abhilfeplan, kann die zuständige Stelle Divestment anordnen.

- Als klimawirksam gelten Anlagen, die messbar zur Zielerreichung beitragen (u. a. CR-fähige Projekte/Emittenten, Elektrifizierung/Effizienz, naturbasierte Senken) oder EU-Taxonomie-konform sind; Doppelzählung ausgeschlossen.

- Der Transformationsfonds Klima (Art. 2) bietet ein standardisiertes Anlagevehikel; keine Staatsgarantie über gesetzliche Mechanismen hinaus.

- Transparenz, Register \& Green-Claims
- Quartals-/Jahresoffenlegung der Exposures und Fortschritte gemäß VO 14a; Daten maschinenlesbar (Art. 17.14).

- Produkt-/Marketing-Claims müssen Registry-/DPP-belegt sein; Greenwashing führt zu Clawback/Sanktionen (Art. Art. 17.12/Art. 17.6).

- Klimarisiko-Pflichten für Finanzakteure
- Szenarioanalysen/Stresstests und Übergangspläne kompatibel mit Art. 1; aufsichtliche Umsetzung im EU-Rahmen.

- Beschaffung \& Bankbeziehungen der öffentlichen Hand
- Divestment-/Übergangspläne sind Zuschlagskriterium; Hausbanken/Depotstellen ohne Zero-Exposure werden grundsätzlich nicht beauftragt (Ausnahmen mit Begründung/Auslaufplan).

- EU-/Völkerrechts-Kohärenz
- Umsetzung im Einklang mit EU-Taxonomie, SFDR/CSRD und Aufsichtsrecht; Notifizierungen/Abstimmungen nach Art. 17.


%% ----------------------------------------
\section[VO 14a: Durchführungs- und Transparenzregeln Divestment]{VO 14a: Durchführungs- und Transparenzregeln Divestment}
\noindent\textit{ID: 38 -- vo14a (meta)}

\begin{quote}
Setzt das Zero-Exposure-Prinzip für öffentliche Portfolios um: keine Neuinvestitionen in fossile Aktivitäten, geordneter Run-off bestehender Exposures, strikte Look-through-Pflichten, standardisierte Offenlegung. Staatlich kontrollierte Beteiligungen folgen der TFP-VO (Transformationsfahrpläne) und sind hiernach gesondert zu berichten.
\end{quote}

{\footnotesize\begin{quote}
Klare Binär-Screens statt Prozentdebatten; strenge Look-through-Pflichten; TFP-Pfad für staatlich kontrollierte Unternehmen als separate Schiene.
\end{quote}}


- Geltung \& Zero-Exposure-Prinzip
- Gilt für Portfolios von Bund/Ländern/Kommunen und deren Sonder-/Treuhandvermögen, inkl. Pensions-/Versorgungswerken und öffentlichen Unternehmen ohne staatliche Kontrolle.

- Zero-Exposure: Neuanlagen in Emittenten/Instrumente mit fossiler Aktivität sind unzulässig; bestehende Exposures laufen geordnet aus (Run-off/Verkauf bis Stichtagen gem. Gesetz).

- Ausnahme staatliche Kontrolle: Unternehmen mit beherrschendem Einfluss des Staates unterfallen der TFP-VO (Transformationsfahrpläne); deren Papiere zählen nicht als fossile Exposure, sind jedoch separat zu berichten.

- Aktivitätskatalog \& Verbotsliste
- Verbotene fossile Aktivitäten (Anlage FOSS-A): Kohlebergbau/-verstromung, Öl-/Gas-E\&P, Ölsande/Schieferöl, LNG-Liquefaction, Raffinerien, Ferntransport-Pipelines, fossile Strom-/Wärmeerzeugung ohne Abscheidung.

- Verwandte Services (Anlage FOSS-S) wie Bohrdienstleister, Offshore-Services, spezialisierte Ausrüster gelten als inaktivierbar fossile Aktivität, sofern ≥ [VO]\% Umsatz/CAPEX an FOSS-A gekoppelt sind.

- Liste wird quartalsweise aktualisiert (NACE/NAICS-Mapping) und als Open-Data veröffentlicht (Art. 17.14).

- Look-through \& Passivmandate
- Fonds/ETFs sind look-through zu bewerten. Fossile Anteile sind unzulässig; Umstieg auf ex-Fossil-Indizes binnen [VO 6--12] Monaten.

- Technische De-minimis-Toleranz für passives Tracking ≤ 0,5 \% des Portfoliowerts p. a. während der Umstellung, mit Plan zur vollständigen Eliminierung.

- Run-off, Sicherungen \& enge Ausnahmen
- Run-off: Veräußerung bei Liquidität; ansonsten Laufzeitende ohne Re-Invest. Keine Zeichnung von Folgeemissionen fossil-bezogener Emittenten.

- Sicherungs-/Übergangsfinanzierungen nur bei dokumentierter Systemkritik (Versorgung/Sicherheit), befristet, mit Auslaufplan und Veröffentlichung; Genehmigung durch das zuständige Organ.

- Green/Transition Bonds von in der Verbotsliste geführten Emittenten sind grundsätzlich unzulässig; zulässig nur bei ring-fenced, insolvenzfesten Strukturen und TFP-Pfad des Emittenten im Einklang mit 1,5 \$\textasciicircum{}\textbackslash\{\}circ\$C; Entscheidung per begründetem Beschluss.

- Schwellen \& Mapping
- „Fossil-dominant“-Schwellen entfallen; stattdessen gilt der Aktivitätskatalog (Anlage FOSS-A/-S) als Binär-Screen.

- Mapping zu EU-Taxonomie/SFDR/CSRD erfolgt über Verweise in Art. 17; Inkonsistenzen sind im Bericht zu begründen.

- Berichts-Templates \& IT-Formate
- Standardvorlagen: Bestände/Flüsse, Exposures nach FOSS-A/-S, Übergangspläne, Szenarioannahmen, De-minimis-Nutzung.

- Maschinenlesbar (API/JSON/CSV) gem. Art. 17.14; QSeal-gesiegelte Quartals-/Jahresberichte.

- Durchsetzung \& Kohärenz
- Prüfrechte der Agentur; automatisierte Abgleiche mit externen Datenprovidern/Indizes; Stichprobenrevision.

- Sanktionen: Rüge, Veröffentlichung, Mittelbindung; Clawback bei Verstößen; Rechtsmittel nach DV/AR (Art. 17.4/17.14).

- Sunset/Review \& Versionierung
- Evaluation alle 24 Monate; Anpassung an EU-Rechtsfortschritt (Taxonomie/SFDR/CSRD).

- Versionierung/Änderungsjournal über die Klima-Datenplattform; historische Screens reproduzierbar.


%% ----------------------------------------
\section[VO 14b: Transformationsfahrpläne staatlicher Beteiligungen]{VO 14b: Transformationsfahrpläne staatlicher Beteiligungen}
\noindent\textit{ID: 39 -- vo14b (meta)}

\begin{quote}
Konkretisiert die in Art. 14 (Divestment) vorgesehenen Transformationsfahrpläne (TFP) für staatliche Mehrheits-/Kontrollbeteiligungen: Geltung, KPIs/MRV, Vergütungskopplung, Investitions-Gates, Reporting, Transparenz und Eskalations-Trigger. Schnittstellen zu DV/ID/AR (Art. 17.4/17.12/17.14).
\end{quote}

{\footnotesize\begin{quote}
Schlankes Vollzugsgerüst: harte Ziele, prüfbare KPIs, Vergütungskopplung, Invest-Gates und klare Eskalation -- mit AR/ID/DV-Verweis statt Doppelnormen.
\end{quote}}


- Geltung \& Kontrolle
- Gilt für Unternehmen, an denen Bund/Land/öffentliche Träger unmittelbar oder mittelbar die Mehrheit der Stimmrechte halten oder beherrschenden Einfluss ausüben (§ 17 AktG-Maßstab).

- Minderheitsbeteiligungen: Stewardship-Pfad (Abstimmungsverhalten, Engagementbericht) gem. Abs. 10.

- TFP-Inhalte (Mindeststandard)
- Zielsystem: Scope 1+2 absolut ≥ 70 \% bis 2030, 90 \% bis 2033, 100 \% bis 2035; material relevante Scope-3-Kategorien mit sektoradäquatem Pfad.

- Residualgrenze: 2035 ≤ 5 \%; nur hoch-permanente Entnahmen (D100) ohne Doppelzählung.

- CAPEX-Pfad: Mindest-Taxonomie-CAPEX-Quote (≥ [VO] \% ab 2028; ≥ [VO] \% ab 2031).

- Fossil-Run-off: Inventar fossiler Vermögenswerte mit Auslaufplan (Inbetriebnahmejahr, Restlaufzeit, Rückbau/Umrüstung).

- Versorgung/Netzdienlichkeit: Sicherungsmaßnahmen und ggf. Übergangskapazitäten mit befristeten Auflagen.

- MRV \& KPIs
- Emissions-MRV an ETS/DEHSt „once-only“ angebunden; Register-Anker (Projekt-/Serien-ID) für Klima-KPIs.

- Kern-KPIs: tCO\$\_2\$e absolut (S1/2/3*), Intensität (z. B. tCO\$\_2\$e/MWh, tCO\$\_2\$e/t Produkt), Taxonomie-CAPEX-Quote, EE-/H\$\_2\$-Anteil, Methan-Leckage/Abfackelrate (wo einschlägig), Fossil-Run-off-Fortschritt, Invest-Stop-Verstöße (= 0).

- Externe Assurance: zunächst limited, nach 2 Jahren reasonable; Berichtsformat nach AR-Spezifikation (Art. 17.14).

- (*Scope 3: nur für material relevante Kategorien je Sektor, definiert in sektoralen MRV-Handbüchern.)

- Vergütung \& Governance
- Variable Vergütung von Vorstand/GF mit mind. 40 \% an TFP-Meilensteine gekoppelt; Malus/Clawback bei Fehlangaben/Zielverfehlung (≥ Gelb, vgl. Abs. 8).

- Eigentümerstrategie/Gesellschafterweisungen setzen TFP verpflichtend; Aufsichtsgremien erhalten Klima-Fachausschuss.

- Investitions-Gates \& Stop-Lines
- No new unabated fossil: keine neuen unabated Förder-/Erzeugungskapazitäten oder Lebensdauerverlängerungen > [VO] Jahre.

- Carbon-Schattenpreis: interne Bewertung ≥ nationaler CO\$\_2\$-Preisplan (Art. 7) + Risikoaufschlag; dokumentierter TCO/NPV mit CO\$\_2\$-Pfad-Sensitivität.

- Ausnahmen nur für sicherheits-/systemkritische Übergänge, befristet und mit Auslaufplan.

- Reporting \& Transparenz
- Jährlicher TFP-Bericht (Open-Dokument, geschäftsgeheimniswahrend) + Quartals-KPIs an Ampel-Dashboard (Art. 17.14).

- Alle statusändernden Events QES/QSeal und qualifiziert zeitgestempelt (Art. 17.12); Daten „once-only“ (Art. 17.4/17.14).

- Eskalations-Trigger
- Gelb (−1 Jahr ggü. Pfad): Bonus-Kappung, 90-Tage-Nachsteuerungsplan.

- Orange (−2 Jahre): Investitionsstopp fossil-relevanter CAPEX, Sonderprüfung, Vergütungs-Malus ≥ [VO] \%.

- Rot (−3 Jahre): Run-off/Carve-out Plan; Ultima Ratio: Divestment gem. Art. 12.

- Just-Transition-Pfad
- Verknüpfung zu SA-Kapitel (Art. 17.7): TKUG/Weiterbildung, Standort-Wandelpläne, Sozialplan-Koordination; Berichtspflicht über Beschäftigungs-KPIs.

- Stewardship bei Minderheitsbeteiligungen
- Abstimmungsleitlinien (Say-on-Climate, Net-Zero-Pläne, Vergütungskopplung) und jährlicher Engagement-Report; Eskalation: Investitionsauflagen → Stimmrechtsausübung → Divestment-Prüfung.

- Beihilfe/Haushalt \& Compliance
- Eigentümervorgaben sind keine Beihilfe; Förderungen/Garantien nur über EU-Beihilfe-Gate (CEEAG/GBER) und haushaltsrechtliche Titel (Art. 17.2).

- Verstöße gegen Reporting/MRV/Stop-Lines ⇒ Art. 17.4/17.14 und Sanktionsrahmen (Clawback/Ausschluss/Veröffentlichung).

- Härtefallklausel \& Review
- Bei nachweislich fehlender Infrastruktur (z. B. H\$\_2\$-Netz, Netzanschluss) kann die Zielzeit um max. 24 Monate verschoben werden; Kompensationsmaßnahmenpflicht (zusätzliche Minderungen anderorts im Konzern).

- Evaluation der VO nach 24 Monaten; Anpassung an Sektor-MRV-Handbücher.


%% ----------------------------------------
\section[Artikel 15: Katastrophenschutz, Vorsorge und Resilienzförderung]{Artikel 15: Katastrophenschutz, Vorsorge und Resilienzförderung}
\noindent\textit{ID: 40 -- artikel15 (gesetz)}

\begin{quote}
Prävention zuerst, klarer Umgang mit Risiken und schneller Wiederaufbau nach Build-Back-Better. Öffentliche Mittel folgen transparenten Risiko­zonen; Wiederaufbau an unveränderten Hochrisiko­standorten wird vermieden. Übungen, Beteiligung und offene Daten erhöhen die Handlungsfähigkeit vor Ort.
\end{quote}

- Monitoring \& Warnung. Klima- und Extremwetter-Monitoring durch die Klima-Transformationsagentur; barrierefreie Warnungen. Daten-/Schnittstellenstandards nach Art. 17 ff.

- Risikozonen \& Raumordnung (Bund-Länder). Der Bund setzt einheitliche Rahmenkriterien für Klima-Risikozonen (Hochwasser, Hitze, Sturm, Waldbrand, Dürre); die Länder weisen Zonen aus und integrieren sie in Raum-/Bauordnungsrecht. Methodik/Geodaten via Art. 17 ff/VO.

- Förderkonditionen (Build-Back-Better). Öffentliche Wiederaufbau- und Vorsorgemittel sind an Verlagerung, Anpassung oder geordneten Rückzug geknüpft, wenn ein Standort in Hochrisikozone liegt. Ausnahmen/Härtefälle per VO mit Begründungs- und Sicherungspflichten.

- Resilienz-Standards. Wiederaufbau erfolgt mit resilienzsteigernden Bau-/Infrastrukturstandards (z. B. Retention, Schwammstadt, Hitzeschutz, Redundanz). Standardkatalog/Musterlösungen nach Art. 17 ff/VO.

- Versicherungs- und Risikopool. Einrichtung eines öffentlich-privaten Risikopools (Rückversicherungsfenster) mit Versicherern/Sozialversicherungen; Staatliche Mittel dürfen keinen Wiederaufbau am identischen Hochrisiko-Standort erzwingen. Governance/Datenaustausch/Kompatibilität (z. B. Solvency) per VO; IT/Transparenz nach Art. 17 ff.

- Lernen nach Ereignissen. Ereignis- und Lernberichte (Schäden, Wirkung der Vorsorge, Anpassungsbedarf); ohne Maßnahmenbeschluss keine weitere Bundesförderung für den identischen Risikokonflikt am selben Standort.

- Übungen \& Beteiligung. Jährliche sektorübergreifende Übungen; Kommunen bieten mindestens zwei öffentliche Beteiligungsformate pro Jahr (hybrid, barrierefrei). Mustern/SLAs siehe Art. 17 ff.

- Digitale Tools \& Qualifizierung. Förderung von Simulations-/Planungstools, Risiko-GIS, offenen Lehrmaterialien; Schulungen für Einsatz-/Bauverwaltungen. Open-Data/IT nach Art. 17 ff.

- Verzahnung mit Instrumenten. Transformationskredit (Art. 3) und Transformationsförderung (Art. 4) sind für Resilienz-/Anpassungsinvestitionen nutzbar. Carbon-Rewards (Art. 5) nur bei zusätzlichem Klimanutzen; Doppelzählung ausgeschlossen.

- Verordnungsermächtigung. Bundesregierung konkretisiert Rahmendefinitionen, Zonen-Methodik, Standardkatalog sowie Parameter des Risikopools. Notifizierungen/Koordination und Open-Data-Pflichten gem. Art. 17 ff.


%% ----------------------------------------
\section[VO 15a: KRITIS-Klimastresstest (KIKS)]{VO 15a: KRITIS-Klimastresstest (KIKS)}
\noindent\textit{ID: 41 -- vo15a (meta)}

\begin{quote}
Betreiber kritischer und wesentlicher Infrastrukturen prüfen ihre Klima-Resilienz regelmäßig in standardisierten Stresstests. Die Verfahren sind sektorkompatibel (NIS2/Kritis-Vorgaben), beschleunigen Priorisierungen und verknüpfen Maßnahmen mit den Finanzierungsinstrumenten des KlimaGG -- ohne Doppelstrukturen.
\end{quote}

{\footnotesize\begin{quote}
Verzahnung ohne Doppelung: Methodik/Daten via Art. 17 (Datenplattform/Once-Only), Finanzierung über Art. 2/KKR, Aufsicht in Kooperation mit den Sektorregulierern; KIKS ergänzt diese um den Klima-Stresstest-Baustein.
\end{quote}}


- Geltung \& Abgrenzung. Verpflichtet sind Betreiber kritischer und wesentlicher Dienste in den Sektoren Energie (Erzeugung/Netze/Speicher), Wasser/Abwasser, Gesundheit/Pflege, Verkehr/Logistik, IKT/Cloud/Rechenzentren, öffentliche Verwaltung/Sicherheitsorganisationen sowie Finanzmarkt-Infrastruktur ab sektor-spezifischen Schwellen [VO]. Fachaufsichten (z. B. BNetzA, BSI, BBK, Länder) bleiben unberührt; KIKS ergänzt ausschließlich um Klima-Risikokomponenten.

- Methodik \& Szenarien. Grundlage sind einheitliche Gefahrenbibliotheken (Hitze, Dürre, Hochwasser/Überflutung, Sturm, Waldbrand, Kombinationsereignisse) aus amtlichen Datenquellen (z. B. DWD/HLNUG) mit Design-Basis und Extrem-Szenarien (Perzentile/Return-Level). Zeitachsen: 5/10/20 Jahre. Abhängigkeiten (Strom-/Daten-/Wasser-/Zuliefer-Ketten) sind verbindlich zu modellieren; Unsicherheiten werden als Bandbreiten ausgewiesen. Templates/Parameter liefert die Verordnung.

- Turnus \& Anlässe. Erstprüfung binnen 12 Monaten nach Inkrafttreten; danach alle 24 Monate. Außerordentliche Tests bei maßgeblichen Anlagenänderungen, nach Großereignissen oder bei Indikator-Triggern [VO].

- Kennzahlen \& Schwellen. Zu berichten sind u. a.: erwartete Versorgungsdefizite (Energie/Wasser/IT-Dienst), Mean-Time-to-Recover, Black-Start/Insellast-Fähigkeiten, Backup-Dauer, Hitze-/Kälte-Robustheit kritischer Bereiche (z. B. OP-Säle), Überflutungs-/Brandexposition, Single-Point-of-Failure. Ampel-Schwellen (grün/gelb/rot) setzt die Verordnung.

- Resilienzpläne. Innerhalb 6 Monaten nach jedem KIKS: priorisierter Maßnahmenplan (No-Regret zuerst) mit Zeitplan, Kosten-/Nutzen, Abhängigkeiten und CO\$\_2\$-Wirkung (Synergien bevorzugt). Maßnahmen sind -- wo passend -- mit Transformationskredit und KKR verknüpfbar; Doppelförderung wird ausgeschlossen.

- Übungen \& Nachweise. Jährliche sektorübergreifende Übungen (Tisch/Sim/Live), für Netzbetreiber inkl. Nachweis von Black-Start/Insellast nach Standard [VO]. Protokolle, Lessons-Learned und Umsetzungsstand sind zu hinterlegen.

- Meldewege \& Transparenz. Einreichung digital bei der Klima-Transformationsagentur mit gleichzeitiger Unterrichtung der zuständigen Fachaufsichten. Executive Summary wird veröffentlicht;  sicherheitsrelevante Details unterliegen abgestufter Vertraulichkeit. Datenformate/APIs gem. Art. 17 (Once-Only, Open-Data-Regeln).

- Aufsicht \& Abhilfe. Stellt die Agentur mit der Fachaufsicht eine akute Gefährdung fest, kann sie einen verbindlichen Abhilfe-Fahrplan vorgeben [VO: Fristen/Eskalation]. Wiederholte Nichtbefolgung führt zu Sanktionsmechanismen nach dem einschlägigen Sektorrecht; gerichtlicher Rechtsschutz nach Art. 8.

- Finanzierung \& Beihilfe. KIKS-Maßnahmen gelten als förderfähig im Rahmen Transformationskredit/KKR. Beihilfekonformität wird über das Harmonisierungskapitel (Art. 17 Beihilfe) gewährleistet.

- Verordnungsermächtigung. Die Bundesregierung legt per VO fest: Schwellen/Erfasste, Szenario-Parameter, Methodik-Templates, KPI-Definitionen, Ampel-Schwellen, Übungsstandards, Fristen/Eskalation, Datenformate und Schnittstellen zu NIS2/BSIG/Branchenvorgaben.

- Evaluation \& Sunset. Erste Evaluation nach 24 Monaten (Wirksamkeit/Kosten/Koordination), danach alle 36 Monate. VO-Parameter unterliegen der Review-/Sunset-Pflicht.


%% ----------------------------------------
\section[Artikel 15a: Hitzeschutz in öffentlichen Gebäuden]{Artikel 15a: Hitzeschutz in öffentlichen Gebäuden}
\noindent\textit{ID: 42 -- artikel15a (meta)}

\begin{quote}
->nachbessern!
\end{quote}

- Pflicht zur Vorsorge. Bei Neubau, Umbau oder Sanierung sind wirksame Hitzeschutz-Systeme vorzusehen; reine Provisorien genügen nicht.

- Solar-Kopplung als Regel. Klimatisierung ist grundsätzlich mit PV zu koppeln; Ausnahmen bei technischer/baulicher Unzumutbarkeit (z. B. Statik, Denkmalschutz) gegen Begründung.

- Kritische Infrastruktur. Krankenhäuser, Pflegeeinrichtungen, Schulen/Kitas erhalten Mindest-Autarkiezeiten (Solar+Speicher) für Netzausfälle; Parameter per VO.

- Finanzierung. Transformationskredit (Art. 3) und Haushaltsmittel der Träger; Vergabe/Beihilfe nach Art. 17 ff.

- Aufklärung. Bund/Länder informieren über Hitzeg efahren und Vorsorge; Materialien offen lizenziert. Publikationsstandards: Art. 17.14.


%% ----------------------------------------
\section[Artikel 15b: Verbraucherschutz und Transparenz bei Energietechnologien]{Artikel 15b: Verbraucherschutz und Transparenz bei Energietechnologien}
\noindent\textit{ID: 43 -- artikel15b (meta)}

\begin{quote}
Gute Entscheidungen brauchen klare Infos: Klimawirkung, Lebensdauer, Folgekosten und Auslaufrisiken müssen verständlich sein -- ohne Technologieverbote, mit ehrlicher Beratung.
\end{quote}

- Transparenzpflichten. Staatliche Stellen veröffentlichen verständliche Informationen zu Klimawirkung, Lebensdauer, TCO/Folgekosten und Auslaufrisiken betroffener Technologien. Formate/Barrierefreiheit gem. Art. 17.14.

- Hinweis- und Beratungspflichten. Anbieter müssen bei erkennbaren Auslaufrisiken oder erhöhten Betriebskosten transparent hinweisen und klimakompatible Alternativen darstellen; neutrale staatliche Beratung ergänzt. Leitfäden/Standards nach Art. 17 ff.

- Schutzklausel (eng). Zur Abwehr irreführender Vermarktung oder erheblicher Sicherheits-/Schadensrisiken kann die Bundesregierung zeitlich befristete Vertriebsbeschränkungen per VO erlassen; Verhältnismäßigkeit, Übergangsfristen und EU-Kohärenz gem. Art. 17 ff.


%% ----------------------------------------
\section[Artikel 15c: Kommunaler Klima- \& Resilienzfonds (KKR)]{Artikel 15c: Kommunaler Klima- \& Resilienzfonds (KKR)}
\noindent\textit{ID: 44 -- artikel15c (meta)}

\begin{quote}
Katastrophenschutz und Vorsorge passieren vor Ort. Der KKR schafft planbare, zweckgebundene Mittel für Kommunen -- mit klarer Formel und Standardmodulen. Ein Top-Up kann Carbon-Rewards finanziell verstärken, ohne Mengen doppelt zu zählen.
\end{quote}

- Zweck \& Gegenstand. Stärkung von Vorsorge, Katastrophenschutz, Strukturwandel und Innovation; der KKR ist ein zweckgebundenes Fenster des Transformationsfonds Klima (Art. 3).

- Mittelhöhe. Jährlicher Bedarfsrahmen per VO (€/Einwohner-Basis, Risiko-Index, Sozial-Faktor; Referenzdaten, Ober-/Untergrenzen). Nicht verausgabte Mittel werden in das Folgejahr übertragen.

- Verteilung. Länderzuweisung als formelgebundene Mittel (50 \% Einwohner, 30 \% Risiko-Index, 20 \% Sozial-Faktor); 1:1-Durchleitung an Kommunen nach identischem Schlüssel.

- Zweckbindung \& Maßnahmen. Mind. 50 \% Vorsorge/Katastrophenschutz (z. B. Retention, Deichrückverlegung, Schwammstadt, Warnnetze, Trinkwasser-Resilienz, Feuerwehr-Ausstattung, naturnahe Maßnahmen); bis 30 \% Strukturwandel; bis 20 \% Innovation. Standardmaßnahmen-Katalog per VO; Abweichungen bei Wirknachweis möglich.

- Top-Up für Carbon-Rewards. Für Projekte nach Art. 5 (Linie A/B) kann ein finanzieller Top-Up von [VO] \% auf CR-Auszahlungen gewährt werden (Cap [VO] €/Projekt·Jahr); keine Anrechnung auf tCO\$\_2\$e-Mengen; Doppelzählung ausgeschlossen.

- Mindestsicherung. Sinkt der kommunale Dividendefluss/KAV-Ersatz unter [VO] €/Einwohner·Jahr, gleicht der KKR auf.

- Verfahren \& Transparenz. Antrag über die Klima-Serviceplattform (Art. 10a), Once-Only-Datenübernahme, SLAs (30/90 Tage) für Standardmodule; Open-Data-Dashboard nach Art. 17.14.

- Haushalt \& Beihilfe. Zweck-/formelgebundene Bundesmittel an Länder; fachliche Zuständigkeiten bleiben unberührt; Beihilfe-/Vergabekohärenz nach Art. 17 ff.

- Rückforderung \& Roll-Over. Zweckwidrige Verwendung → Rückforderung/Sperrfristen; nicht verausgabte Kommunalmittel werden automatisch übertragen.


%% ----------------------------------------
\section[VO 15c: Durchleitung \& Regionaltopf]{VO 15c: Durchleitung \& Regionaltopf}
\noindent\textit{ID: 45 -- vo15c (meta)}

\begin{quote}
Planbare Mittelströme, faire Raumverteilung.
\end{quote}

- Durchleitung. Länder leiten KKR-Mittel binnen 30 Tagen weiter; schuldhafte Verzögerung: Verzugszinsen [VO] und Veröffentlichung im Quartalsbericht der Agentur.

- Regionaltopf. Mindestens [VO] \% der KKR-Jahresmittel sind für ländliche/struktur­schwache Räume zu reservieren (einheitliche Kriterien, Raumtypenklassifikation).

- Transparenz. Quartalsweise Open-Data-Berichte zu Zuweisungen/Abflüssen je Land/Region gemäß Art. 17.14.


%% ----------------------------------------
\section[Artikel 15d: Adaptives Wohnen \& Freiwilliger Umsiedlungsfonds (AWF)]{Artikel 15d: Adaptives Wohnen \& Freiwilliger Umsiedlungsfonds (AWF)}
\noindent\textit{ID: 46 -- artikel15d (meta)}

\begin{quote}
Wo Klimarisiken dauerhaft hoch bleiben, braucht es faire, freiwillige Wege in sichere Quartiere. Der AWF schafft klare Angebote: Kauf-/Tauschoptionen, Miet-Neutralität, Ersatzstandorte -- gekoppelt an die Risikozonen aus dem Katastrophenschutz und die Build-Back-Better-Regeln. Ziel: Schutz von Menschen und Vermögen, ohne Zwang und ohne endlose Einzelfall-Verhandlungen.
\end{quote}

- Zweck \& Grundsätze. Der AWF ermöglicht freiwillige Umsiedlungen, Standortverlagerungen und wohnortnahe Ersatzlösungen in ausgewiesenen Klima-Risikozonen (vgl. Katastrophenschutz-Artikel), nach den Prinzipien Freiwilligkeit, Sozialverträglichkeit, Transparenz, Build-Back-Better und Keine Zwangsverlagerung. Der Schutz besonders vulnerabler Gruppen hat Vorrang.

- Anspruchsberechtigte. Natürliche Personen mit Hauptwohnsitz, soziale Einrichtungen und Kleinst-/KMU mit Betriebsstätte in einer als hoch eingestuften Risikozone; Details (Zonenklassen, Schwellen, Nachweise) regelt die Verordnung.

- Instrumente (Wohnen).
- Kauf-/Tauschoption \& Landbank: Kommunen/Länder können über eine Landbank freiwillige Aufkäufe und Wohnungstausche zu fairen Verkehrswerten mit definiertem Zuschlag (z. B. Umzugs-, Transaktions-, Ausstattungspauschalen) anbieten. Rückbau/Entsiegelung des Altstandorts ist Regelfall.

- Miet-Neutralität: Haushalte erhalten einen befristeten Miet-Neutralitätszuschuss beim Umzug in sichere Lagen; Kappungs- und Belegungsbindungen können mitfinanziert werden.

- Umzugs- \& Anpassungspauschalen: Standardisierte Pauschalen für Umzug, Grundausstattung und besondere Bedarfe (Pflege, Barrierefreiheit).

- Instrumente (Gewerbe/Allgemeinbedarf).
- Ersatzstandortförderung für KMU/Kritische Versorgung (Apotheken, Pflege, Handwerk, Nahversorgung) inklusive Planungs- und Genehmigungs-Fast-Track.

- Gemeinbedarf (Kita, Schule, Pflege, Rettung): priorisierte Ersatzbau-Programme mit Build-Back-Better-Standard; Doppelförderung ist auszuschließen.

- Planungs- \& Genehmigungsprivilegien. Für Ersatzstandorte gelten die Beschleunigungsregeln, Standard-Nebenbestimmungen und Muster der Klima-Serviceplattform; Innenentwicklung und sichere Höhenlagen werden planungsrechtlich priorisiert.

- Finanzierung. Der AWF wird aus dem Kommunalen Klima- \& Resilienzfonds (KKR) gespeist und kann durch Mittel des Transformationsfonds sowie zweckgebundene Drittmittel ergänzt werden. Die Kofinanzierungsanteile und Deckel (€/Fall, €/m²) werden per Verordnung festgelegt. Haushalts- und Beihilferechtliche Vorgaben folgen den Harmonisierungsvorschriften.

- Sozialer Schutz \& Anti-Verdrängung. Vorrang für ortsnahe Lösungen; Härtefallklauseln; besonderer Schutz für Mieter:innen, Pflegebedürftige, Alleinerziehende und Menschen mit Behinderungen. Keine Negativ-Bonitierung in der Wohnungsvergabe durch AWF-Inanspruchnahme.

- Flächen \& Rückbau. Ankaufte Flächen werden renaturiert, entsiegelt oder risikomindernd umgenutzt (Retention, Grünzüge). Nutzungsbindungen und Dienstbarkeiten werden im Register erfasst; Doppelförderung ist auszuschließen.

- Daten, Transparenz, Monitoring. Offene Kennzahlen (Anträge, Bewilligungen, Kosten je Fall, Risikominderung, Sozialindikatoren) werden maschinenlesbar veröffentlicht. Der Datenschutz bleibt gewahrt; bereits eingereichte Daten dürfen nicht erneut verlangt werden.

- Rechtsschutz. Entscheidungen im AWF-Verfahren unterliegen beschleunigtem Rechtsschutz; Muster- und Sammelverfahren sind zulässig.

- Verordnungsermächtigung. Die Bundesregierung wird ermächtigt, durch Rechtsverordnung festzulegen: Zonen-Schwellen/Trigger, Bewertungs- und Zuschlagsformeln, Miet-Neutralitätsparameter, Pauschalen (Umzug/Anpassung), Landbank-Musterverträge, Kofinanzierungsregeln, De-minimis-Erleichterungen für KMU, Schnittstellen zu Katastrophenschutz und KKR sowie Evaluationskriterien. EU-Notifizierungen erfolgen gemäß Harmonisierung.

- Evaluation \& Sunset-Disziplin. Erste Evaluation nach 24 Monaten (Wirksamkeit, Verteilung, Kosten), danach alle 48 Monate; Parameter-VO mit Review/Sunset, sofern nicht verlängert.

{\footnotesize\begin{quote}
Der AWF verzahnt sich mit: Risikozonen/Build-Back-Better (Katastrophenschutz), Klima-Serviceplattform (Muster/SLAs), Kommunalem Klima- \& Resilienzfonds (KKR) und Transformationskredit (für Ersatz- und Anpassungsinvestitionen). Allgemeine Daten-, Beihilfe- und Haushaltsregeln stehen im Harmonisierungskapitel; hier keine Doppelung.
\end{quote}}



%% ----------------------------------------
\section[VO 15d: Durchführungsregeln AWF]{VO 15d: Durchführungsregeln AWF}
\noindent\textit{ID: 47 -- vo15d (meta)}

- Zonierung \& Trigger. Verbindliche Klassen (z. B. RZ-I hoch, RZ-II mittel) mit Datengrundlagen, Aktualisierungszyklus und Veröffentlichung.

- Bewertung \& Zuschläge. Formel zur Ermittlung fairer Verkehrswerte + Zuschläge (Umzug, Transaktion, Ausstattung, Barrierefreiheit); Ober-/Untergrenzen; regionale Korrekturfaktoren.

- Miet-Neutralität. Parameter (Höchstbetrag/Monat, Dauer, Bindungen), Schnittstelle zur Wohnungswirtschaft, Missbrauchsschutz.

- Landbank \& Verträge. Musterkauf-/Tauschverträge, Dienstbarkeiten, Rückbaupflichten, Renaturierungs-/Nutzungsbindungen; Register-Mapping.

- KMU \& Gemeinbedarf. Pauschalen/Quoten, De-minimis-Regeln, Beschleunigungs-SLA, Safe-Harbor-Beschaffung.

- Finanzierung. Kofinanzierungs- und Deckelparameter; Roll-over-Regeln; Verknüpfung mit KKR/TFK.

- Monitoring. Pflicht-KPIs (Zeit bis Angebot/Umzug, €/geretteter Haushalt, EAD-Minderung), Open-Data-Formate, Evaluationsrhythmus.


%% ----------------------------------------
\section[Artikel 15e: Nationaler Wasser- und Dürre-Resilienzrahmen (NWR)]{Artikel 15e: Nationaler Wasser- und Dürre-Resilienzrahmen (NWR)}
\noindent\textit{ID: 48 -- artikel15e (meta)}

\begin{quote}
Sabine, Binnenschifferin bei Kaub. Früher hieß Sommer: Bangen um den Pegel. Heute fährt ihr flachgehendes Schiff zuverlässig -- finanziert über den Transformationskredit, die Ladeplanung stützt sich auf offene Pegel-/Engpassdaten. Im Chemiepark am Ziel puffern Elektrolyseur und Speicher das Netz und verdienen Carbon Rewards (FLEX/RT). Kommunale Retentionsräume wurden über den KKR bezahlt; die Bahn hat Notkorridore ausgebaut. Wirkung: Ausfalltage: 45 → 6/Jahr. Tonne-km per Schiff/Schiene: +22 \%. Logistik-CO\$\_2\$ pro Tonne: −28 \%.

  Wasserknappheit und Starkregen nehmen zu. Der NWR schafft einheitliche Regeln für Vorsorge, Dürremanagement und Speicherausbau -- mit klaren Triggern, Vorrangregeln und offenen Daten. Er ergänzt Artikel 15 (Katastrophenschutz) und nutzt die Instrumente aus Artikel 2 (Transformationskredit) und 10c (KKR).
\end{quote}

- Zweck \& Geltung: Bund legt einen Rahmen für Dürre- und Wassermanagement fest: einheitliche Trigger, Vorrangregeln bei Knappheit, Speicher-/Rückhalteausbau, Wasserwiederverwendung. Länder konkretisieren in Einzugsgebieten; kommunale Satzungen bleiben möglich (Art. 17 ff.).

- Wasserkonten \& Open Data: Für Flussgebiete werden Wasserkonten geführt (Zufluss, Speicher, Entnahmen, Mindestökologie). Monatliche, maschinenlesbare Veröffentlichung; Schnittstellen zu Raumordnung, Landwirtschaft und Industrie (Art. 17.14).

- Dürrestufen (Trigger): Bundesweit einheitliche Stufen (Info/Alarm/Krise/Extrem) auf Basis kombinierter Indizes (Bodenfeuchte, Abfluss, Grundwasser, Temperatur). Stufen lösen definierte Maßnahmenpakete aus; Details per Verordnung.

- Vorrangregeln bei Knappheit: Priorität gilt (1) Trinkwasser/öffentliche Gesundheit, (2) kritische Infrastruktur/Grundversorgung, (3) Mindestökologie, (4) Produktion/Landwirtschaft nach Effizienz-/Substitutionskriterien. Ausnahmen nur gegen dokumentierte Zieläquivalenz.

- Speicher, Rückhalt, MAR: Förderung von Retentionsflächen, Auenreaktivierung, dezentralen Schwammstadt-Maßnahmen und Managed Aquifer Recharge (MAR). Standard-Genehmigungspfade über Art. 8a/8b; Transformationskredit (Art. 2) und KKR (Art. 10c) sind nutzbar.

- Wasserwiederverwendung: Einheitliche Mindeststandards für industriell/kommunal aufbereitetes Wasser; schnelle Zulassung für definierte Anwendungen; EU-Konformität (Art. 17).

- Leitungsnetze \& Verluste: Zielwerte und Transparenzpflichten für Wasserverluste; Sanierung ist förderfähig (Art. 2/10c). Jahresberichte als Open Data.

- Verzahnung mit Risiko-/Raumordnung: NWR-Karten sind in Planung/Bauleitplanung zu berücksichtigen; Build-Back-Better (Art. 15) bleibt unberührt.

- Verordnungsermächtigung: Bundesregierung legt Trigger-Indizes, Vorrang-/Ausnahmekriterien, Datenformate, Mindeststandards Wiederverwendung/MAR und Monitoringzyklen fest; Notifizierung nach Art. 17.

{\footnotesize\begin{quote}
Fokus auf Governance (Trigger/Vorrang/Daten) statt Einzelfalltechnik; Finanzierung über TFK/KKR, Verfahren über Art. 8a/8b.
\end{quote}}



%% ----------------------------------------
\section[VO 15e: NWR -- Schwellen, Vorrang, Daten \& Verfahren]{VO 15e: NWR -- Schwellen, Vorrang, Daten \& Verfahren}
\noindent\textit{ID: 49 -- vo15e (meta)}

\begin{quote}
Operative Parameter für Trigger-Stufen, Vorrangregeln, Speicher/Wiederverwendung, Datenformate und Vollzug. IT/Datenschnittstellen, Beihilfe-/EU-Kohärenz und Vergabe richten sich nach den Harmonisierungskapiteln (Art. 17 ff.).
\end{quote}

- Trigger-System Dürre/Knappheit:
- Indizes: Bodenfeuchte-Index, Abfluss-Index, Grundwasser-Index, Niederschlags-Anomalie; Quellen/Laufzeiten werden als offene Referenzdatensätze veröffentlicht (Art. 17.14).

- Stufen \& Schwellen: Info (>[VO] \% Abweichung), Alarm (>[VO] \%), Krise (>[VO] \%), Extrem (>[VO] \%); Mindestkriterien kombinierter Indizes je Flussgebiet.

- Aktualisierung: Wöchentliche Bewertung; Stufen-Bekanntgabe maschinenlesbar (API/GeoJSON) mit Zeitstempel und Gebietskulisse.

- Maßnahmenpakete je Stufe (bundeseinheitlicher Rahmen):
- Info: Transparenzpflichten, Hinweis-Bewässerungspläne, freiwillige Verbrauchsreduktion in Verwaltung/öffentlichen Betrieben.

- Alarm: Priorisierung nach Ziff. 3; neue wasserintensive Sondernutzungen nur mit Nachweis Ersatz-/Wiederverwendung; Prüfung von Lastverschiebung/Prozesswasser.

- Krise: Temporäre Nutzungsbeschränkungen für nicht-kritische Anwendungen; Anordnung Wasserwiederverwendung bei geeigneten Anlagen; Auflagen an Kühlwasser, Reinigung/Rezirkulation.

- Extrem: Kontingentierung nach Vorrangreihenfolge, Sicherstellung Trinkwasser/öffentliche Gesundheit; Freigabe von Nothilfetransporten/Speichern; zügige Ersatzwassererschließung.

- Vorrang- und Ausnahmeprüfung bei Knappheit:
- Vorrangreihenfolge: (1) Trinkwasser/öffentliche Gesundheit; (2) kritische Infrastruktur/Grundversorgung; (3) Mindestökologie; (4) übrige Nutzungen nach Effizienz/Substitutionsfähigkeit.

- Ausnahmen: Nur bei dokumentierter Zieläquivalenz (gleiche/geringere ökologische Belastung, gleiche Versorgungssicherheit); Befristung/Monitoring zwingend.

- Speicher, Rückhalt, MAR \& Schwammstadt:
- Standardmodule: Retentionsflächen, Auenreaktivierung, Deichrückverlegung, Versickerungsfelder, Zisternenverbünde, Gründächer/Entsiegelung, Managed Aquifer Recharge (MAR).

- Genehmigung: Kategorie-Zuweisung und Typauflagen nach Art. 8b; Standard-MRV-Kits (Volumen, Infiltration, ökologische Begleitwerte) mit Safe-Harbor.

- Wasserwiederverwendung (Re-Use):
- Qualitätsklassen: Klassen A--C je nach Verwendung (z. B. Bewässerung Industrie/Grün, Reinigungswasser, Kühlkreisläufe); Mindeststandards, Monitoring, Rückverfolgbarkeit.

- Zulassung: Fast-Track für Umrüstung bestehender Klär-/Industrieanlagen bei Einhaltung Standardmodule; Schnittstelle Gesundheits-/Lebensmittelrecht bleibt unberührt.

- Leitungsnetze \& Verluste:
- Zielkurve: Reduktion realer Netzverluste auf ≤[VO] \% in [VO] Jahren; jährliche Zwischenziele; Veröffentlichung je Versorger.

- Förderfähigkeit: Sanierungen/Leckage-Monitoring über Transformationskredit (Art. 2) und KKR (Art. 10c) als Standardmodule.

- Daten, Karten \& Open-Interfaces:
- Wasserkonten: Monatsbilanzen (Zufluss, Entnahme, Speicher, Mindestabgabe) je Einzugsgebiet; offene APIs/Layer (Art. 17.14).

- Projekt-Layer: Verbindliche Eintragung von NWR-Projekten (Standort, Modul, Wirkung, Status).

- Finanzierung \& Verzahnung: Transformationskredit (Art. 2), KKR (Art. 10c) und Transformationsförderung (Art. 3) mit Vorrang für Retention/MAR/Re-Use-Module; Beihilfe-Kohärenz nach Art. 17.

- Übergang \& Evaluation: 12-Monats-Pilot in ausgewählten Einzugsgebieten; jährliche Evaluation (Wirkung, Kosten, soziale Effekte); parameterisierte Anpassung per VO.


%% ----------------------------------------
\section[Artikel 16: Klima-Notfall]{Artikel 16: Klima-Notfall}
\noindent\textit{ID: 50 -- artikel16 (gesetz)}

\begin{quote}
Für seltene, gravierende Lagen: schnelle, zielgebundene Eingriffe mit parlamentarischer Kontrolle, klaren Sunset-Fristen und Make-Up-Pflicht. Der Klima-Notfall nach diesem Gesetz ergänzt das allgemeine Katastrophenrecht; die Ziele aus Art. 2 bleiben maßgeblich. Pauschale Preisdeckel oder Ausgleichszahlungen ohne Klimawirkung sind ausgeschlossen.
\end{quote}

- Zweck \& Abgrenzung. Ein Klima-Notfall liegt nur vor, wenn eine atypische, zeitkritische Lage die Zielerreichung nach Art. 2 akut gefährdet oder außergewöhnliche Preis-/Versorgungsschocks drohen, die ohne Gegensteuerung die Transformation substanziell unterminieren (z. B. extreme Hitzewellen mit hoher Mortalität, großflächige Überflutung, großräumige Waldbrände). Grundsätze: Verhältnismäßigkeit, Zielbindung, Transparenz, Sunset, Make-Up.

- Notfall-Feststellung \& Trigger.
- Wissenschaftlicher Trigger: Die Klimakommission (Art. 8) stellt auf Basis der Verfahren nach Art. 2a eine akute Gefährdung von Budget/Pfad fest und veröffentlicht binnen 4 Wochen einen Maßnahmenkatalog mit Wirk-/Risikoanalyse.

- Ökonomischer/Versorgungs-Trigger: Bei außergewöhnlichen, nicht selbstverschuldeten Preisschocks (CO\$\_2\$/ Energie) oder kritischen Infrastrukturausfällen legt die Bundesregierung binnen 4 Wochen einen Stabilisierungsplan vor.

- Parlament: Der Bundestag entscheidet binnen weiterer 4 Wochen. Unterbleibt ein Beschluss, tritt die Backstop-Verordnung nach Art. 2a in Kraft.

- Zulässige Notfall-Instrumente (Katalog). Nur folgende, befristete Maßnahmen sind zulässig:
- Preisglättung im steuerlichen Anteil der CO\$\_2\$-Bepreisung (Art. 7) innerhalb des gesetzlich festgelegten Korridors (Make-Up-Pflicht; kein Absinken des Floors unter den Mindestpfad).

- Dividendenglättung durch Vorziehen/Glättung aus der Dividenden-Rücklage (Art. 8a), mit transparenter Nachholung.

- Beschleunigte Carbon-Reward-Fenster (Art. 5) für kurzfristig mobilisierbare Minderungs-/Entnahmemengen mit Safe-Harbor-MRV; No-Offset/Claim-Regeln bleiben unberührt.

- Priorisierung Transformationskredit (Art. 3) für kritische Infrastruktur (Gesundheit, Wasser, Kommunikation, Energie, Verkehr) und Hitzeschutz/Kühlfähigkeit in vulnerablen Einrichtungen; SLAs nach Art. 8a gelten im Notfallmodus verkürzt.

- Beschleunigungsrahmen (Art. 8b) im Notfallmodus: engere Fristen/Fiktionen innerhalb der dortigen Leitplanken; Aarhus-/Umweltrecht bleibt gewahrt.

- Anwendung der Risikozonen (Art. 15): Wiederaufbau in Hochrisikozonen nur mit Verlagerung/Anpassung; „Build-Back-Better“ verpflichtend.

- EU-ETS bleibt unberührt: Nationale Eingriffe in den EU-ETS (inkl. MSR) oder „Zertifikatsreserven“ finden nicht statt.

- Haushalt \& Beihilfe. Maßnahmen erfolgen innerhalb haushaltsrechtlicher Ermächtigungen und beihilferechtlich genehmigter Rahmen; die Verordnung legt jährliche Höchstvolumina und etwaige Notifizierungspflichten fest. Überplanmäßige Ausgaben sind ausgeschlossen.

- Sunset \& Evaluation. Notfallmaßnahmen gelten für 6 Monate; eine einmalige Verlängerung um 6 Monate ist mit Begründung möglich. Ex-post-Evaluation binnen 3 Monaten.

- Make-Up-Pflicht. Zeitweilige Preisglättung oder Verzögerung der Floor-Anhebung wird binnen 24 Monaten aufgeholt (Preis- und Mengenwirkung), u. a. durch steilere Pfadfortschreibung (Art. 7) und/oder zusätzliche CR-Mengen (Art. 5).

- Transparenz \& Rechtsschutz. Laufendes Notfall-Monitoring-/Dashboard, Veröffentlichungen gemäß Art. 17.14; Eilrechtsschutz vor der Klimagerichtsbarkeit (Art. 8).

- Missbrauchssperre. Wiederholung gleicher Maßnahmen innerhalb von 24 Monaten bedarf einer verschärften Begründung (Fortschritts-/Alternativnachweis). Strukturelle Politikänderungen dürfen nicht per Notfall-VO erfolgen.

{\footnotesize\begin{quote}
Koordination im Notfall über die Klima-Transformationsagentur (Art. 8) in Abstimmung mit Ländern/Kommunen; IT/Once-Only/Offenlegung nach Art. 17 ff.
\end{quote}}



%% ----------------------------------------
\section[VO 16a: Klima-Notfall -- Schwellen, Verfahren, Transparenz]{VO 16a: Klima-Notfall -- Schwellen, Verfahren, Transparenz}
\noindent\textit{ID: 51 -- vo16a (meta)}

\begin{quote}
Operative Schwellwerte, Berechnungslogiken und Fristen -- einheitlich, prüfbar, schlank -- mit vollständiger Daten-/IT-Anbindung nach Art. 17.14.
\end{quote}

{\footnotesize\begin{quote}
EU-ETS bleibt unberührt; nationale Eingriffe beschränken sich auf den steuerlichen Anteil (Art. 7), Dividenden-Rücklage (Art. 8a), CR-Fenster (Art. 5) und Vollzugsbeschleunigung (Art. 8a/8b).
\end{quote}}


- Schwellwerte (Trigger-Indikatoren).
- Budgetlücke (Art. 2/2a): prognostizierte Zielverfehlung > [VO] \% des nationalen Budgets oder > [VO] Mt im Jahr N+1.

- Preisschock (Art. 7): 30-Tage-Median relevanter Preise > [VO] \% über 180-Tage-Baseline und Gefahr der Zielunterminierung (Impact-Analyse).

- Versorgungs-/Infrastruktur-Shock: erwartete zusätzliche Emissionen > [VO] Mt oder Netzdienstlichkeitsdefizit > [VO] kritische Stunden/Monat; betroffene Schutzgüter (Gesundheit, Wasser, Energie) sind zu benennen.

- Verfahrensfristen \& Inhalte.
- T₀ Feststellung/Ankündigung; binnen 14 Tagen Notfall-Sachstand (Daten/Modelle/Unsicherheiten),

- T₀+28 Maßnahmenkatalog (Optionen, Wirk-/Kostenanalyse, Sunset, Make-Up-Pfad),

- T₀+56 Parlamentsentscheidung; bei Ausbleiben greift die Backstop-VO nach Art. 2a.

    Alle Datengrundlagen/Modelle werden gemäß Art. 17.14 versioniert, maschinenlesbar veröffentlicht.

- Beschleunigte CR-Fenster. Quartalsfenster für Zusatzmengen (Art. 5) mit Safe-Harbor-MRV-Kits, Volumenlimit [VO], beihilferechtskonform (CEEAG/GBER), keine Preisaufschläge über den Einheits-Floor; Claims/No-Offset bleiben unberührt.

- Preisglättung \& Glättungsparameter. Algorithmen für temporäre Anpassung des steuerlichen Anteils (Art. 7) und für Dividendenvorzug (Art. 8a); Obergrenzen, Rückführungs-/Make-Up-Pfad, Transparenz nach Art. 17.14.

- Beschleunigungsrahmen. Aktivierung der Notfall-SLAs nach Art. 8a/8b; Stop-the-clock nur für fehlende Unterlagen/EU-Beteiligungen, max. [VO] Tage; offene Muster/Typauflagen werden priorisiert.

- Haushalt/Beihilfe-Caps. Jährliche Höchstvolumina [VO] (Mrd. €) je Instrument; Notifizierungspflichten, Monitoring und Bericht an den Bundestag; keine überplanmäßigen Ausgaben.

- Dashboard \& Reporting. Notfall-Dashboard (Preis, Budgetlücke, CR-Mengen, Kredit-Priorisierungen, Zeitplan) gemäß Art. 17.14; Abschlussbericht binnen 3 Monaten nach Sunset.

- Sunset \& Review. VO-Parameter treten nach 24 Monaten automatisch außer Kraft, sofern sie nicht evaluiert/verlängert werden; Updates richten sich nach Art. 2a.


%% ----------------------------------------
\section[Artikel 17: Harmonisierung und Vorrang]{Artikel 17: Harmonisierung und Vorrang}
\noindent\textit{ID: 52 -- artikel17 (gesetz)}

\begin{quote}
Dieses Kapitel passt bestehende Rechtsmaterien an das KlimaGG an und verhindert Doppelstrukturen. Es bündelt querschnittliche Regeln, die alle KlimaGG-Teile verbinden. Maßstab sind: Zielvorrang (Art. 1), klare Zuständigkeiten (Art. 8), Beteiligung/Transparenz (Art. 8a), Datenharmonisierung (Art. 17.14) und EU-/Völkerrechtskonformität.
\end{quote}


%% ----------------------------------------
\section[Artikel 17.1: Teil A -- Allgemeiner Teil]{Artikel 17.1: Teil A -- Allgemeiner Teil}
\noindent\textit{ID: 53 -- art-17-1-allgemeiner-teil (meta)}

\begin{quote}
Legt Begriffe, Zielprinzipien und die technische/organisatorische Basis fest: Zusätzlichkeit, Doppelzählungssperre, Systemkohärenz (EU/ETS/CBAM/CORSIA), Once-Only, Identifikatoren \& API-Schnittstellen, Datenschutz sowie Verordnungsermächtigungen. Gilt querschnittlich für alle Instrumente und Register des KlimaGG.
\end{quote}

- Geltung \& Verhältnis
- Dieser Teil gilt für sämtliche Verfahren, Registervorgänge, Förder-/Preis­instrumente und Berichte nach dem KlimaGG.

- Vorrang des Unionsrechts bleibt unberührt. Fachrecht (z. B. Vergabe-, Beihilfe-, Außenwirtschafts-, Steuerrecht) gilt fort; dieser Teil setzt Konsistenz- und Schnittstellenregeln (lex specialis KlimaGG geht abweichendem Bundesrecht vor, soweit ausdrücklich vorgesehen).

- Begriffe \& Prinzipien
- Zusätzlichkeit: Förder- oder Reward-Fähigkeit nur oberhalb gesetzlicher Mindeststandards/Quoten; Pflichtmengen sind nicht rewardfähig.

- Doppelzählungssperre: Jede tCO\$\_2\$e darf nur einmal angereizt/beansprucht werden; Registry-Sperr-/Abgleichsmechanismen sind verbindlich.

- Systemkohärenz: Keine Verrechnung gegen EU-ETS/ETS2, CBAM oder CORSIA-Pflichten; Claims und No-Offset richten sich nach Artikel 5 (Carbon Reward) und den dortigen Leitfäden.

- Ambitionsneutralität: Methodik-Updates/Re-Basierungen verändern das Ambitionsniveau nicht (Zieläquivalenz).

- Do-No-Harm: Keine Verschlechterung von Biodiversität, Arbeit-/Gesundheitsschutz oder Sicherheit; Mindeststandards sind einzuhalten.

- IDs, Daten \& Once-Only
- Verbindliche Identifikatoren: Projekt-ID, Serien-ID (Vintage/Line), Register-ID; Mapping auf LEI/ISIN/EORI/HS-Code und Referenzen zum Digitalen Produktpass (DPP) sowie Anlagen-/Standort-IDs.

- Once-Only: Erforderliche Daten werden einmal erhoben und über definierte APIs mit Register, Klima-Daten-Plattform (KDP), Monitoring und Statistik geteilt; Mehrfachabfragen sind unzulässig.

- Events \& Zeit: Statusändernde Events sind qualifiziert zu signieren und mit ISO-8601-Zeitstempel (UTC) zu versehen; Versionierung und Änderungsjournal sind verpflichtend.

- Vertrauensdienste: Signaturen/Siegel und Zeitstempel entsprechen eIDAS-Anforderungen; Details in der Vertrauensdienste-/ID-VO dieses Kapitels.

- Status-Flags \& Sperren: Das Register kennt bindende Status-Flags: import-locked, export-claim-lock, aid-pending, redeem-locked. Einlösungen (redeem) sind bei aid-pending bis zur Beihilfe-Genehmigung gesperrt.

- Transparenz \& Kontrolle
- Öffentliches Verfahrens-/Monitoring-Dashboard (Ampel-Logik) mit maschinenlesbaren Aggregaten; personenbezogene Daten werden pseudonymisiert bereitgestellt.

- Prüf- und Einsichtsrechte der zuständigen Kontrollorgane (insb. Rechnungshof, Markt-/Finanzaufsicht, Stabilitätsrat) bleiben unberührt.

- Datenschutz (Querschnitt)
- Datensparsamkeit, Zweckbindung, Pseudonymisierung; Portabilität für DPP/Registry-Belege.

- Aufbewahrung grundsätzlich bis max. 10 Jahre nach letzter Einlösung/Bestandskraft, sofern spezialrechtlich keine längeren Fristen gelten; anschließend Löschung/Archiv nach DS-VO.

- Verordnungsermächtigungen (Teil-A-VOs)
- ID-API-VO: Datenmodelle (IDs/Events), Endpunkte, Webhooks, Fehlercodes, Sandbox/Conformance-Tests.

- DS-VO: Pseudonymisierungsprofile, Lösch-/Archivprozesse, Audit-Logs, Zugriffsklassen, Datenteilungsklauseln (Once-Only).

- KDP-Grundprofil-VO: Profile für Aktenführung/Zustellung, Sicherheitsniveaus (NIS2/BSI), Fallback-Verfahren bei IT-Ausfällen, Betriebs-/Verfügbarkeitsmetriken.

- Change-\&-Versioning-VO: SemVer, Mindestparallellaufzeiten, Deprecation-Policy (≥ 6 Monate), Release-Zyklen, Conformance-Zertifikate als Go-Live-Gate.

{\footnotesize\begin{quote}
Leitidee: eine Nummer pro Projekt, eine Quelle pro Fakt -- prüfbar, versioniert, maschinenlesbar.
\end{quote}}



%% ----------------------------------------
\section[Artikel 17.2.1: Teil B -- Haushalt \& Finanzverfassung]{Artikel 17.2.1: Teil B -- Haushalt \& Finanzverfassung}
\noindent\textit{ID: 54 -- art-17-2-haushalt (meta)}

\begin{quote}
Operationalisiert die verfassungsrechtlichen Klima-Ziele haushalts- und schuldenverfassungsfest: eng begrenzte Klima-Investitionsregel, Tilgung/Transparenz, TFK-Instrumente, Carbon-Reward-Floor als Eventualverbindlichkeit, ESA/Beihilfe-Kohärenz sowie klare Trennung zu Netzentgelten/Umlagen und eine formelgebundene kommunale Durchleitung.
\end{quote}

- Klima-Investitionsregel (Korridor)
- Investive Ausgaben mit nachweisbarer Emissionsminderung oder Resilienzsteigerung können -- innerhalb eines gesetzlich definierten Korridor-Caps (in \% BIP bzw. absolute Obergrenze) -- bei der strukturellen Nettokreditgrenze berücksichtigt werden, im Rahmen von Art. 109/115 GG und der in Artikel 1 dieses Gesetzes normierten Verfassungsgrundsätze.

- Anerkennungsfähig nur investive Ausgaben, MRV-basiert und zieläquivalent zum nationalen Pfad; konsumtive Ausgaben ausgeschlossen.

- Parlamentsvorbehalt: jährlicher TFK-Wirtschaftsplan (Volumen, Tilgungsplan, Rücklagen-Mindestniveaus) mit Begründung und Bericht.

- Haushaltsdisziplin \& Kontrolle
- Tilgungspfad und Cap-Ausnutzung sind zu beschließen; bei Ampel-Status „Rot“ (17.14) Maßnahmenplan binnen 30 Tagen.

- Prüfrechte des Rechnungshofs, Stabilitätsrat-Meldungen (quartalsweise), öffentliche Ampel-Berichte gem. 17.14.

- Transformationskredit (TFK) \& Förderlogik
- Revolvierende, investive Instrumente (Darlehen/Beteiligungen/Garantien) mit Expected-Loss-Rückstellungen und festem Tilgungsplan; Risiko-Reporting quartalsweise.

- Zuschussanteile nur bei nachgewiesener Förderlücke (CEEAG-Proportionalität); Beihilfe-Gate vor Wirksamwerden.

- Carbon-Reward-Floor (CRF)
- Staatliche Mindestauszahlungszusage wird als Eventualverbindlichkeit geführt; Finanzierung über Garantiegebühren und eine zweckgebundene Risikorücklage (Cap/Mindestniveau).

- Einlösungen ex post als gebundene Ausgabe; keine Marktankäufe/kein Eigenbestand des Bundes an CR-Ansprüchen.

- ESA-/EU-Konformität
- ESA-richtige Klassifizierung von Kreditaufnahmen, Darlehen/Beteiligungen und Garantien; Gebührenzuflüsse stärken Rücklagen ohne vorweg haushaltswirksam zu werden.

- EU-Beihilfe-Konformität (CEEAG/GBER) über Beihilfe-Gate; Veröffentlichung der genehmigungsrelevanten Parameter im Ampel-Reporting.

- Beihilfe-Gate \& Redeem-Lock
- Für beihilferelevante Rewards gilt: Ausgabe (issuance) zulässig; Einlösung (redeem) bis zur Genehmigung gesperrt (aid-pending/redeem-locked im Register nach 17.1/17.12).

- Spur 1: GBER/De-minimis für Kleinvolumina mit standardisierten Parametern. Spur 2: Pre-Notification parallel zum VO-Erlass; Spur 3: Full Notification mit Netting bei Auflagen.

- AR-Reporting: Kennzahlen „pending-Tage“, Volumenanteil aid-pending, Time-to-Decision.

- Trennung zu Netzentgelten/Umlagen
- Instrumente dieses Gesetzes (u. a. CRF, TFK, KKR-Zuweisungen, kommunaler Reward-Bonus, Klimadividende) werden weder über Netzentgelte noch über EEG/KWKG/sonstige Umlagen refinanziert (Abgrenzung: 17.17 EN).

- Kosten-/Rücklagenführung ausschließlich haushaltsrechtlich; keine kostenwirksame Erfassung in ARegV/StromNEV-Mechanismen.

- Kommunale Durchleitung \& Kofinanzierung
- KKR-Mittel als zweck- und formelgebundene Zuweisungen 1:1 an Kommunen (vgl. 17.6 KH); Weiterleitung binnen 30 Tagen.

- Eigene Haushaltsstellen/Titelgruppen; offene Projektlisten mit Projekt-ID und Registry-Abgleich (Once-Only nach 17.1).

- Keine Verrechnung mit anderen Förder-/Haushaltstöpfen; Kofinanzierung im Rahmen der vorgegebenen Zweckbindung.

- Transparenz, Monitoring \& Reporting
- Quartalsweise Offenlegung: Kreditbestände, Tilgung, Rücklagen, Risiken/Einlösungen (CRF), Beihilfe-Status; Veröffentlichung im Ampel-Dashboard (17.14).

- Jährlicher Wirkungsbericht: €/tCO\$\_2\$, regionale/soziale Verteilungseffekte, Risikoanalyse (Zins-/Credit-/Marktrisiken).

- IT-Kapazitäten \& Betrieb
- Für Register, KDP und AR-Plattform wird ein ringfenced IT-Budget (Betrieb/Weiterentwicklung/Conformance) samt Planstellen beschlossen; Mittel werden als investive Klima-IT dem Korridor zugerechnet.

- Der Bund stellt für zentrale digitale Bausteine dieses Gesetzes Open-Source-Referenzimplementierungen bereit oder fördert deren Entwicklung.

- Conformance- und Pen-Test-Kosten sind förderfähig; OPS-SLOs und Incident-Berichte werden nach 17.14 veröffentlicht.

- Sanktions- \& Rückforderungsbezug: Falschangaben, Doppelförderung, zweckwidrige Verwendung: Clawback, Bußgelder, Sperrvermerke im Register nach den jeweiligen Einzelartikeln und den einschlägigen Harmonisierungsvorschriften; Veröffentlichung aggregierter Maßnahmen im Ampel-Reporting.

- Verordnungsermächtigungen (Teil-B-VOs)
- HB-FinVO: Korridor-Cap, Tilgungsschemata, Rücklagen-Mindestniveaus, Ampel-Schwellen/Trigger, Berichtsfelder/Formate.

- ESA-Klass-VO: Detailabbildung ESA-Kategorien/Verbuchung, Risikomodelle, Offenlegungsanhänge.

- Verweis auf EN-EP-VO (17.17) zur Abgrenzung gegenüber Netzentgelten/Umlagen/Rechnungsausweisen.

{\footnotesize\begin{quote}
- Kern: Eng begrenzte, MRV-fundierte Investitionsschiene mit Cap, Tilgung und Transparenz -- Schuldenregel bleibt gewahrt.

- Planbarkeit: ESA-/Beihilfe-ready, kommunaltauglich, ohne Nebenwirkungen auf Strompreise.
\end{quote}}



%% ----------------------------------------
\section[Artikel 17.2.2: Überleitung nicht gebundener KTF-Mittel in die Klimadividende]{Artikel 17.2.2: Überleitung nicht gebundener KTF-Mittel in die Klimadividende}
\noindent\textit{ID: 55 -- art-17-2-ktfg (meta)}

\begin{quote}
Harmonisierung mit dem Klima- und Transformationsfonds-Gesetz (KTFG): Nicht verausgabte oder freiwerdende KTF-Mittel werden -- nachrangig zu EU-/Rechtsbindungen -- automatisch in die Klimadividende (Art. 8) überführt. Ziel: Doppelstrukturen vermeiden; ETS-Einnahmen fließen entweder in Transformation oder direkt an Bürgerinnen und Bürger.
\end{quote}

- Grundsatz \& Verhältnis
- Der KTF bleibt Sondervermögen nach KTFG; Mittelverwendung richtet sich nach diesem Gesetz und Art. 17.2 (Haushalt \& Finanzverfassung).

- Die Überleitung ergänzt die haushaltsrechtlichen Regeln; Art. 109/115 GG und beihilferechtliche Vorgaben bleiben unberührt.

- Definitionen
- Nicht gebunden sind Mittel, die zum Jahresende weder verausgabt noch durch rechtsverbindliche Verpflichtungen (Verträge/Bewilligungen/VE) belegt sind.

- Freiwerdend sind Mittel, die aus Rückflüssen, Aufhebungen oder Minderabrufen bewilligter Programme hervorgehen.

- Rangfolge
- Vorrangig zu bedienen sind völker-/unionsrechtliche Bindungen, genehmigte Beihilfemaßnahmen sowie rechtsverbindliche Bewilligungen.

- Die Überleitung erfolgt vorrangig vor neu aufgelegten, nachgeordneten Programmen ohne Rechtsbindung.

- Mechanik der Überleitung
- Abschläge: Ab Q3 des Haushaltsjahres werden monatliche Abschläge in den Dividendenpool (Art. 8) geleistet, bemessen an der konservativen Prognose der nicht gebundenen Jahresreste.

- Jahresabschluss: Endabrechnung und vollständige Überleitung des Restbetrags bis 30. Juni des Folgejahres.

- Die Mittel fließen in die Glättungsrücklage der Klimadividende (Art. 8a) und erhöhen den pro-Kopf-Betrag gemäß dortiger Formel.

- Anti-Umgehung
- Programmaufsetzungen/-erweiterungen nach dem 1. Oktober eines Haushaltsjahres dürfen allein nicht zur Vermeidung der Überleitung dienen; hierfür ist ein Bedarfstest (Zielbezug, Reifegrad, Beihilfe-Gate) zu dokumentieren.

- Unverhältnismäßige Vorgriffe (VEs ohne Realisierungsplan) bleiben bei der Definition „nicht gebunden“ unberücksichtigt.

- Transparenz \& Reporting
- Quartalsberichte (Prognose \& Ist) zu Zuflüssen, Bindungen, Überleitungen (Schaubild: R/B/Z) im Ampel-Dashboard nach Art. 17.14.

- Jahresbericht der Bundesregierung an den Bundestag: Zuflüsse, Verwendungen, Restbestände, Prognose und Auswirkungen auf die pro-Kopf-Dividende.

- ESA/Beihilfe \& Abgrenzung
- Die Überleitung in die Klimadividende ist nicht beihilferelevant; beihilferechtliche Bindungen des KTF sind vorab zu erfüllen.

- Keine Finanzierung über Netzentgelte/Umlagen; haushaltsrechtliche Abwicklung gemäß Art. 17.2.

- Verordnungsermächtigung
- KTF-Überleitungs-VO: Begriffs-/Abgrenzungsregeln (gebunden/freiwerdend), Prognosemethodik, Abschlagsstaffel, Endabrechnung, Schnittstelle zum Dividendenpool, Offenlegungsfelder/Formate (Art. 17.14), Prüfrechte.

- Übergang
- Im ersten Anwendungsjahr: Mindestüberleitung von 80 \% der als „nicht gebunden“ klassifizierten Mittel; ab dem zweiten Jahr 100 \%.

{\footnotesize\begin{quote}
Vermeidet „Parken“ von ETS-Einnahmen und stärkt die pro-Kopf-Rückverteilung, ohne Rechtsbindungen oder EU-Vorgaben zu verletzen.
\end{quote}}



%% ----------------------------------------
\section[Artikel 17.3: Teil C -- Außenwirtschaft \& Zoll (AX)]{Artikel 17.3: Teil C -- Außenwirtschaft \& Zoll (AX)}
\noindent\textit{ID: 56 -- art-17-3-teil-c-ax (meta)}

\begin{quote}
Stellt Rechts- und Datensicherheit für grenzüberschreitende Warenströme und KlimaGG-Registerwertrechte her: Vorrang von EU-Recht (CBAM), MRV-/Registry-Äquivalenz, Doppelzählungsschutz, AML/KYC-Pflichten und IT-Integration mit der Zollverwaltung.
\end{quote}

- Geltungsbereich \& Begriffe
- Gilt für (a) Ein- und Ausfuhr klimarelevanter Waren (einschl. CBAM-Sektoren) sowie (b) grenzüberschreitende Transfers von KlimaGG-Registerwertrechten.

- Begriffe/IDs gemäß Art. 17.1 (Projekt-ID, Serien-ID, Mapping auf EORI/HS/LEI/ISIN, DPP-Referenzen); Daten-/API-Grundsätze gemäß Art. 17.14.

- CBAM-Vorrang \& Differenzprinzip
- Der EU-CBAM hat Vorrang. Nationale Elemente nach Art. 9/9a wirken ausschließlich lückenschließend für nicht erfasste Waren/Emissionen und CBAM-kohärent.

- Carbon-Reward-Einheiten sind kein gezahlter CO\$\_2\$-Preis für CBAM/Zoll; eine Anrechnung auf CBAM/ETS-Pflichten ist ausgeschlossen (vgl. Art. 9).

- Import (Gate-Regeln)
- Zollanmeldungen enthalten ein maschinenlesbares CO\$\_2\$-Datenfeld (PCF/Prozess-ID; ggf. CBAM-Referenz). Schnittstellen/Formatspezifikationen nach Art. 17.14.

- Anerkennung ausländischer Klimaleistungen nur bei MRV- und Registry-Äquivalenz (White-List nach Abs. 7); andernfalls keine Nutzbarkeit für KlimaGG-Claims/Vergaben.

- Vor Freigabe: automatischer Registry-Abgleich (Status clear oder import-locked). Konfliktfälle → pending verification bis Klärung nach Abs. 7.

- Fallback: konservative Default-Emissionswerte werden in Art. 9a (CBAM-VO) festgelegt und jährlich aktualisiert.

- Export (Schutz vor Doppelverwendung)
- Waren mit produktgebundenen Rewards: vor Ausfuhr export-claim-lock im Register; parallele Nutzung der Einheiten im Inland/anderen Systemen ist ausgeschlossen (No-Offset gemäß Art. 5).

- Exportentlastungen für CO\$\_2\$-Kosten nur im unions-/WTO-rechtlich zulässigen Umfang; Rewards begründen keinen ETS/CBAM-Erstattungsanspruch (vgl. Art. 9).

- Freiwillige Wirkungs-Claims im Zielland setzen die Übermittlung von Metadaten (Vintage, Linie/Klasse, Serien, Prüfflaggen) aus dem Register voraus.

- Zoll- \& IT-Integration
- API-Kopplung KlimaGG-Register ↔ Zoll-IT (EORI, HS-Codes, Produktionscharge, CO\$\_2\$-Datenfeld) mit risikoorientierten Prüfpfaden; Spezifikationen nach Art. 17.14.

- Freizonen/Transit unterliegen identischen Gate-Regeln; Freeports begründen keine Ausnahmen („kein Freeport-Leakage“).

- AWG/AWV, Sanktionsrecht \& AML/KYC
- Grenzüberschreitende Transfers von Registerwertrechten unterliegen AWG/AWV und Sanktionslistenprüfung; Sperre bei gelisteten Personen/Ländern; Meldepflichten bleiben unberührt.

- Verpflichtende AML/KYC-Prüfungen bei Inhaberwechseln/Einlösungen; Intermediärsketten nur mit durchgängiger Sorgfaltspflicht. Ident-/Signaturstandards gemäß Art. 17.1 (eID/eIDAS-kompatibel).

- Anerkennungsverfahren \& Interoperabilität
- White-List-VO (unterjährige Pflege) für ausländische MRV-/Registry-Systeme mit Re-Validierung mind. alle 3 Jahre; Echtzeit-Serien-/Eigentumsabgleich als Voraussetzung.

- Konfliktfall: Status pending verification plus Informationspflicht an beide Registerbetreiber; Veröffentlichung aggregierter Konfliktstatistiken gemäß Art. 17.14.

- Green-Claims \& öffentliche Aufträge
- Produkt-/Contribution-Claims nur auf Basis von Registry-Daten und nach Retirement der jeweiligen Einheit (No-Offset-Leitfaden Art. 5); parallele Nutzung in mehreren Systemen ist verboten.

- Vergaben/PPAs dürfen importierte Rewards nur bei bestätigter Äquivalenz (Abs. 7) und gesperrten/zugewiesenen Einheiten berücksichtigen; Vorgaben des Vergaberechts bleiben unberührt.

- Sanktionen, Clawback \& Reporting
- Ordnungswidrigkeiten gestaffelt; systemisches Fehlverhalten: Zwangsgeld, Veröffentlichung, Ausschluss vom Fast-Track (Art. 8a), Registry-Sperre.

- Clawback unrechtmäßiger Vorteile inkl. Zinsen; Einträge/Flags im Register.

- Quartalsberichte: Volumina nach HS, Serien-Zuordnungen, Sperr-/Konfliktfälle, Sanktionen; Open-Data-Aggregate nach Art. 17.14.

- Verordnungsermächtigung (AX-VO)
- CO\$\_2\$-Datenfelder in der Zollanmeldung, White-List-Verfahren, API-Spezifikationen Registry↔Zoll, Prüf-/Sanktionsleitfäden und Fallback-Defaults (soweit nicht bereits in Art. 9a geregelt).

{\footnotesize\begin{quote}
Schnittstellen-Kapitel: MRV/Default-Details verbleiben in Art. 9/9a; hier werden Gate-/Daten-/Vorrangregeln, Registry-Locks und AML/KYC klargezogen.
\end{quote}}



%% ----------------------------------------
\section[Artikel 17.4: Teil D -- Durchsetzung \& Vollzug (DV)]{Artikel 17.4: Teil D -- Durchsetzung \& Vollzug (DV)}
\noindent\textit{ID: 57 -- art-17-4-teil-d-dv (meta)}

\begin{quote}
Einheitliche, digitale und fristenbasierte Vollzugsregeln für KlimaGG-Verfahren: eAkte, Once-Only, Standardfristen mit Genehmigungsfiktion (soweit zulässig), Leitbehörde/Konzentration, Ombudsstelle, IT-Sicherheit (NIS2/BSI) und öffentliches Verfahrens-Dashboard.
\end{quote}

- Geltung \& Verhältnis
- Gilt querschnittlich für KlimaGG-Verfahren und harmonisierte Fachgesetze, soweit Spezialrecht nichts Abweichendes bestimmt.

- Projektbezogene Beschleunigungsregeln (Kategorien, Typauflagen, konkrete Fristen, Fiktionen) ergeben sich aus Art. 8b und der dazu erlassenen VO; diese gehen vor.

- Digitale Verfahrensakte \& Once-Only
- Verbindliche eAkte gemäß EGovG-Standard; Anträge, Unterlagen, Bescheide und Bekanntmachungen werden digital geführt.

- Once-Only/Same-Data: API-Verknüpfung mit dem Register nach Art. 5, der Datenplattform/Monitoring nach Art. 17.14 und Statistik; Doppelmeldungen sind unzulässig.

- Zustellung/Signaturen gemäß Art. 17.12 (eIDAS-konform); qualifizierte Zeitstempel für statusändernde Events.

- Standardfristen \& Genehmigungsfiktion (querschnittlich)
- Vollständigkeitsprüfung binnen 10 Arbeitstagen; Nachforderungen digital und begründet.

- Entscheidungsfristen: Standard 60 Tage ab Vollständigkeit; komplexe UVP-/Planverfahren 120 Tage. Stop-the-Clock nur für gesetzlich definierte Tatbestände.

- Genehmigungsfiktion -- soweit fachrechtlich zulässig und keine UVP-Pflicht besteht -- mit automatisch greifenden Standardauflagen; Einzelheiten in der 8b-VO.

- SLA-Matrix \& Scoreboard: Fristtreue, Nachforderungen/Antrag und Fiktionsquote werden quartalsweise je Land/Stelle ausgewiesen (AR-VO). Bei Rot automatische Eskalation und Maßnahmenplan binnen 30 Tagen.

- Leitbehörde, Konzentration \& Eskalation
- Eine Leitbehörde koordiniert Fachstellen und erlässt einen Bescheid mit Konzentrationswirkung.

- Bei Fristversäumnis: automatische Übernahme durch die Aufsichtsebene; Gründe und Maßnahmen werden im Dashboard veröffentlicht (Art. 17.14).

- Ombudsstelle (angeschlossen an Art. 8a): Entscheidung über Verfahrenskonflikte/Fristüberschreitungen binnen 30 Tagen; digitale Akteneinsicht.

- Beteiligung \& Öffentlichkeit
- Digitale Auslegung/Erörterung (barrierefrei; Kurzfassungen in einfacher Sprache/mehrsprachig); elektronische Stellungnahmen mit automatischer Eingangsbestätigung.

- Öffentliches Verfahrens-Dashboard (Status, Fristen, Änderungen) als maschinenlesbare Open-Data-Aggregate; Personenbezug pseudonymisiert (Art. 17.14).

- Integration Register/MRV \& Monitoring
- Nachweise werden via API in Register (Art. 5) und Monitoring (Art. 17.14) gespiegelt; Mehrerfassungen sind ausgeschlossen.

- Verdacht auf Doppelzählung/Falschangaben → Registry-Sperrvermerk und Meldung an die zuständige Aufsicht; Sanktions-/Clawback-Regeln der Einzelartikel bleiben unberührt.

- IT-Sicherheit, NIS2/BSI \& Business Continuity
- Plattformen sind als wesentliche Einrichtungen abzusichern (ISO/IEC 27001/BSI-Grundschutz; jährliche Pen-Tests; Vier-Augen-Prinzip; HSM-Schlüsselverwaltung).

- BCM-Kennzahlen: RTO ≤ 4 h, RPO ≤ 1 h, Georedundanz; Notfallübungen halbjährlich; Incidents binnen 24 h melden; betroffene Serien/Vorgänge vorsorglich sperren.

- Telemedien-/TK-Betrieb ohne Beeinträchtigung amtlicher Dienste; aktiver Missbrauchsschutz (DDoS/Bot).

- Rechtsschutz
- Eilrechtsschutz vor der Klimagerichtsbarkeit nach Art. 8 mit Folgenabwägung (Emissionsbudget/Resilienzinteressen); Entscheidung binnen 6 Wochen ab Antragsreife.

- Verfahrensbündelung/Verbindung gleichgelagerter Klagen bleibt unberührt; Spezialzuweisungen/Sprungrevisionen nach Fachrecht.

- Aufsicht, Disziplin \& Sanktionen
- Fristen-/Qualitätsaudits halbjährlich; Organisationsverfügungen bei systemischen Verstößen.

- Manipulation, Falschangaben, Registry-Betrug → Clawback/Bußgelder/Ausschlüsse gemäß den Sanktionsregeln der jeweiligen Einzelartikel; Behörden melden Verdachtsfälle unverzüglich.

- Gebühren, Zahlungen \& Open Data
- Digitale Gebührenabwicklung; Pauschalen für Standardfälle; Transparenz der Zahlpläne/Caps.

- Verfahrensoffene Daten als Open Data mit Änderungsjournal; IFG-/Transparenzrechte bleiben unberührt (Art. 17.14).

- Übergang \& Evaluation
- Pilotphase eAkte/Schnittstellen 12 Monate; anschließend volle Anwendung; Altverfahren können optieren.

- Jährliche Evaluation (Tempo, Rechtsschutz, IT-Sicherheit, UX) mit öffentlichem Bericht und Maßnahmenplan.

- Verordnungsermächtigung (DV-VO)
- Fristenprofile, Fiktions-Standardauflagen, eAkte-/API-Profile, Sicherheitsniveaus, Dashboard-Spezifikation, Ombudsverfahren; im Übrigen Verweis auf Art. 8b-VO für projektbezogene Details.

{\footnotesize\begin{quote}
„Eine Akte, eine Datenquelle, klare Fristen“ -- projektbezogene Feinsteuerung bleibt in Art. 8b; 17.4 setzt die querschnittlichen Leitplanken.
\end{quote}}



%% ----------------------------------------
\section[Artikel 17.5: Teil E -- Verbraucherrecht \& Digitaler Produktpass (VC)]{Artikel 17.5: Teil E -- Verbraucherrecht \& Digitaler Produktpass (VC)}
\noindent\textit{ID: 58 -- art-17-5-teil-e-vc (meta)}

\begin{quote}
Verbindet DPP-Daten mit echten Verbraucherrechten: vorvertragliche Information, Label/QR, faire Reward-Bundles, registry-sichere Rückabwicklung, Gewährleistung/Reparatur, Green-Claims-Disziplin und Datenportabilität -- kohärent mit Art. 10b, 17.1, 17.4 und 17.14.
\end{quote}

- Anwendungsbereich \& Verhältnis
- Gilt für B2C-Verträge über Produkte/Dienstleistungen mit Klima-, Energie- oder Ressourcenbezug sowie für optionale Carbon-Reward-Bundles nach Art. 5.

- Ergänzt Art. 10b (Transparenz bei Energietechnologien); technische IDs/APIs nach Art. 17.1, Vollzugs-/Fristenregeln nach Art. 17.4, Daten/Monitoring nach Art. 17.14.

- Vorvertragliche Informationen (BGB/RL 2011/83/EU -- Harmonisierung)
- Bereitzustellen mindestens: aus dem DPP abgeleitete PCF, Energieprofil/-effizienz, erwartete Nutzungsdauer/Haltbarkeit (MTBF), Reparierbarkeit/Teileverfügbarkeit, Rezyklatgehalt, Update-Zusagen, voraussichtliche Entsorgung/EoL-Pfad.

- TCO-Kurzhinweis (CapEx/OpEx/CO\$\_2\$-Kosten); ausführliche TCO via QR/API abrufbar.

- Transformationshinweis bei voraussichtlichem Auslauf (Anknüpfung an Art. 10b Abs. 2).

- DPP-Zugang \& Label
- Freier Zugang zum Digitalen Produktpass (DPP) via QR/Link (barrierefrei); unveränderliche Kernfelder (Signatur/Hash) im Kaufbeleg zu hinterlegen (Art. 17.1 IDs/Signaturen).

- Klima-DPP-Label (QR + Piktogramme) zeigt mindestens Bandwerte für PCF, Haltbarkeit, Reparierbarkeit, Rezyklatanteil; Spezifikation und Schwellen durch VC-VO.

- Optionale Carbon-Reward-Bundles (No-Offset-Kohärenz)
- Separater, eindeutiger Preis; keine Voreinstellung, kein Kopplungszwang.

- Claim-/Eigentumsrechte werden dem Verbraucher im Register nach Art. 5 zugewiesen und gesperrt; Kaufbeleg enthält Serien-/Projekt-ID, Vintage, Menge, Linie/Klasse.

- Doppelzählungsschutz: Einheiten dürfen nicht zugleich für Anbieter-Claims genutzt werden; Sperrvermerk ist zwingend (Art. 5 Claims/No-Offset).

- Widerruf/Retouren \& Registry-Rückabwicklung
- Bei Widerruf/Retoure binnen 14 Tagen: kostenfreie Rückübertragung oder Annullierung der Rewards im Register; automatische Protokollierung über die API (Art. 17.1).

- Bereits veröffentlichte Claims sind zu kennzeichnen und zu korrigieren; Anbieter trägt Nachweis-/Dokumentationspflicht.

- Gewährleistung, Reparatur \& Updates
- Verlängerung der Mängelhaftung um 12 Monate nach erfolgreicher Reparatur; Sicherheits-/Funktionsupdates über die erwartete Nutzungsdauer.

- Reparaturfähigkeit: Zugang zu Ersatzteilen, Diagnosedaten, Reparaturanleitungen; unzulässige Obsoleszenz ist sanktionsbewehrt.

- Unlautere Praktiken \& Green-Claims (UWG-Harmonisierung)
- Unzulässig: pauschal „klimaneutral“ ohne DPP/Registry-Beleg; Verrechnung fremder Rewards; irreführende Zeiträume/Permanenzangaben; Breaching der No-Offset-Regel.

- Pflichtbelege für umweltbezogene Aussagen: Registry-Export, DPP-Snapshot, Prüfberichte. Sanktionen: Bußgeld, Clawback, Claim-Sperre; Veröffentlichung aggregierter Verstöße (Art. 17.14).

- Datenrechte \& Datenschutz
- DPP-Portabilität (maschinenlesbar, API); Einsicht in Registerbelege für verbundene Rewards.

- Personenbezug minimal; Verarbeitung gemäß Datenschutz-Teil in Art. 17.1; Marketing-Profiling nur mit ausdrücklicher Einwilligung.

- Durchsetzung \& Streitbeilegung
- Verbandsklagerecht gegen systemische Verstöße; Marktüberwachung mit stichprobenbasierten DPP-/Registry-Prüfungen.

- ADR: kostenfreie außergerichtliche Schlichtung; Register stellt Widerrufs-/Retourenlogs bereit (Art. 17.1 API-Events).

- Übergang \& Sanktionen
- Übergangsfrist 12--24 Monate für DPP-Integration; Pflichtfelder können stufenweise aktiviert werden.

- Bußgelder umsatz- und schwerebezogen; Wiederholungstäter unterliegen erweiterten Transparenzpflichten.

- Verordnungsermächtigung (VC-VO)
- Label-/DPP-Feldspezifikationen, TCO-Profile, Claim-Beleglisten, Rückabwicklungs-APIs, Aufsichts-/Sanktionsleitfäden; Kohärenz mit Art. 17.1, 17.4, 17.14 sicherzustellen.

{\footnotesize\begin{quote}
DPP wird vom Datensilo zum durchsetzbaren Verbraucherrecht -- prüfbar, rückabwickelbar, ohne Greenwashing.
\end{quote}}



%% ----------------------------------------
\section[Artikel 17.6: Teil F -- Kommunal- \& Haushaltsrecht (KH)]{Artikel 17.6: Teil F -- Kommunal- \& Haushaltsrecht (KH)}
\noindent\textit{ID: 59 -- art-17-6-teil-f-kh (meta)}

\begin{quote}
Sichert die rechtssichere Buchung, Verwendung und Transparenz bundesseitiger KKR-Zuweisungen, kommunaler Reward-Zuflüsse und ggf. kommunaler Durchleitungen -- mit Übertragbarkeit, Prüfkette und Open-Data.
\end{quote}

- Geltung \& Grundsätze
- Mittel des Kommunalen Klima- \& Resilienzfonds (KKR) sind zweck- und formelgebundene Zuweisungen; Länder leiten sie 1:1 an Kommunen weiter; Einbehalte/Verrechnungen sind unzulässig.

- Kommunale Reward-Zuflüsse (z. B. Kommunalbonus) werden nicht über Netzentgelte/Umlagen refinanziert; Abgrenzung nach Art. 17.17 (EN).

- Haushaltsdarstellung \& Titel
- Eigene Titelgruppen/Sachkonten für: (a) KKR-Zuweisungen (Ein-/Ausgabe), (b) kommunalen Reward-Bonus, (c) Kofinanzierungen zu TFK/Programmen.

- KKR-Ausgaben sind vorrangig investiv (Vorsorge/Resilienz); konsumtive Begleitkosten bis [VO] \% zulässig.

- Übertragbarkeit \& Roll-Over
- Nicht verausgabte Mittel werden zweckidentisch ins Folgejahr übertragen; mehrjährige Maßnahmen erhalten Verpflichtungsermächtigungen bis [VO] Jahre.

- Automatisches Roll-Over für bewilligte, noch nicht abgerechnete Projekte; Storno nur mit Begründung und öffentlicher Mitteilung.

- Kofinanzierung \& Matchings
- Kofinanzierungen aus kommunalem Reward-Bonus sind zulässig; Doppelförderung derselben Leistung/Tonne ist ausgeschlossen (Register-Sperrlogik nach Art. 17.1 / Art. 5).

- ÖPP/PPP sind möglich; Vergaben erfolgen nach GWB/VgV/KonzVgV mit TCO-Pflicht dieses Gesetzes; DPP-Daten fließen in die Wertung (vgl. Art. 17.5).

- Prüfung, Controlling \& Open Data
- Rechnungsprüfung nutzt Projekt-IDs (Registry-Mapping) und Georeferenzen; Abgleich mit Ampel-/Monitoring gemäß Art. 17.14.

- Quartalsweise offene Projektlisten (Betrag, Zweck, Status, Indikatoren) nach Datenschema der Klima-Datenplattform (17.14).

- Schulden- und Kassenwirtschaft
- Investitionen in Vorsorge/Resilienz gelten haushaltsrechtlich als investiv; Kreditaufnahme ist zulässig bei Tilgungspfad und TCO-Nachweis.

- Kassenkredite für Klima-Investitionen sind unzulässig; Überbrückung nur bis zur Auszahlung bewilligter KKR-Mittel.

- Vergabe \& TCO-Pflicht
- Alle KKR-finanzierten Beschaffungen verwenden TCO-Bewertung inkl. CO\$\_2\$-/Ressourcenkomponenten; DPP-Felder nach Art. 17.5 sind als Nachweisquelle zuzulassen.

- Rewards als Zuschlagskriterium nur, wenn Einheiten im Register zugunsten der Kommune gesperrt werden; Claims sind öffentlich zu dokumentieren (Register-Export, 17.1/17.14).

- Rückforderung \& Sanktionen
- Zweckwidrige Verwendung führt zu Rückforderung, Sperrfristen und ggf. Ausschluss vom Fast-Track; Wiederholfälle erhöhen die Prüfintensität.

- Standardisierter Haushaltsrahmen (Anlage KH-1)
- Musterkontenplan/Titelmatrix (z. B. 42.xx KKR-Einnahmen, 78.xx KKR-Ausgaben, 45.xx Reward-Bonus) per BMF/BMI-Verwaltungsvorschrift.

- Einheitliche Projekt-ID verknüpft Haushalt, Vergabe, Registry und Monitoring (Once-Only, vgl. 17.1).

- Föderale Schnittstellen
- Länder leiten binnen 30 Tagen durch; schuldhafte Verzögerung: Verzugszinsen [VO] und Nennung im Quartalsbericht der Agentur.

- Kommunalaufsicht berücksichtigt Investitionspfade; Genehmigungsfiktion bei Fristversäumnis nach [VO] Tagen, soweit landesrechtlich zulässig.

- Übergang
- Einführung binnen 12 Monaten; Alt-Titel ins KH-Schema überführen; Pilotkommunen erhalten Fast-Track-Support (vgl. 17.4).

- Verordnungsermächtigung (KH-VO)
- Kontenpläne, Projekt-ID-Schema, Datenformate für Open Data/Monitoring, Schwellen für investive Qualifizierung, Prüf-/Roll-Over-Regeln; Kohärenz mit 17.2, 17.4, 17.17, 17.14.

{\footnotesize\begin{quote}
Klarer Haushaltsrahmen + Übertragbarkeit = Planbarkeit. Offene Daten = Vertrauen. Keine versteckten Umlagen.
\end{quote}}



%% ----------------------------------------
\section[Artikel 17.7: Teil G -- Sozial- \& Arbeitsrecht (SA)]{Artikel 17.7: Teil G -- Sozial- \& Arbeitsrecht (SA)}
\noindent\textit{ID: 60 -- art-17-7-teil-g-sa (meta)}

\begin{quote}
Ziel: Soziale Absicherung, Qualifizierung und geordnete Übergänge in der Transformation. Beschäftigte erhalten einklagbare Weiterbildungs-/Übergangsrechte, Unternehmen klare Anzeige- und Beteiligungspflichten. Kurzarbeit wird zur Qualifizierungszeit, Gesundheitsschutz (u. a. Hitze) gestärkt.
\end{quote}

- Anwendungsbereich \& Verhältnis
- Gilt für arbeits- und sozialrechtliche Begleitregeln der Transformation; Spezialrecht (SGB/BGB/ArbZG u. a.) bleibt unberührt und wird durch diese Harmonisierung ergänzt.

- Datenflüsse, Identifikatoren und Monitoring richten sich nach Art. 17.1 (IDs/Once-Only) und Art. 17.14 (Daten/Monitoring).

- Anspruch auf Transformationsberatung \& -qualifizierung (SGB III, BBiG)
- Beschäftigte in Betrieben mit angezeigter Stilllegung/Umrüstung/Dekarbonisierung haben Anspruch auf Transformationsberatung der BA und priorisierte Weiterbildung.

- Förderfähig primär berufsqualifizierende Module/Abschlüsse in klima-relevanten Profilen (Netze/EE, Wärme/Speicher, Kreislauf/LULUCF, Digitalisierung/MRV); Anerkennung non-formaler Kompetenzen (BBiG-Ergänzung) wird erleichtert.

- Transformations-Kurzarbeitergeld (TKUG) mit Qualifizierungsquote (SGB III)
- TKUG wird gewährt, wenn mindestens [VO] \% der Ausfallzeit für zertifizierte Weiterbildung genutzt werden.

- Leistungshöhe analog KUG; Aufstockung um [VO] Prozentpunkte bei erfolgreichem Abschluss definierter Module; Zuschüsse zu Lehrgangs-/Prüfkosten für Arbeitgeber.

- Frühwarn- \& Beteiligungspflichten (SGB III; BetrVG-Bezug)
- Anzeige geplanter Stilllegungen/Umrüstungen an BA und Klima-Transformationsagentur mindestens 6 Monate vor Wirksamwerden (Once-Only-Schnittstelle nach 17.14).

- Verpflichtende Transformationsvereinbarung mit Betriebsrat/Belegschaft (Zeitplan, Qualifikationspfade, Vermittlung, Sozialplan-Bezug); Veröffentlichung nicht-personenbezogener Eckdaten.

- Arbeitszeit \& Lernzeit (ArbZG/BBiG)
- Weiterbildung im TKUG gilt als Arbeitszeit; Lernzeitkonten und modulare/online-Verteilung sind zulässig.

- Freistellungsansprüche für Prüfungen/Module; Nacht-/Zuschlagsregeln bei E-Learning gem. ArbZG.

- Mobilitäts-, Betreuungs- und Pflegehilfen (SGB II/III/IX/XI)
- Erstattung notwendiger Fahrt-/Umzugskosten zur Qualifizierung/Arbeitsaufnahme; Zuschüsse für Kinderbetreuung und Pflegevertretung während Kurszeiten.

- Inklusionsvorrang (SGB IX): Assistenz, barrierefreie Lernumgebungen, Nachteilsausgleich.

- Gesundheitsschutz \& Hitzeschutz (ArbSchG; SGB VII-Konnex)
- Arbeitsstätten sind bei Hitzewellen nach Stand der Technik zu schützen; Gefährdungsbeurteilungen müssen Hitzerisiken adressieren (techn./organ./pers. Maßnahmen; Mindeststandards per VO).

- Reha-/Präventionsleistungen der Träger werden für klimainduzierte Gesundheitsrisiken geöffnet.

- Sozialausgleich \& Härtefälle (SGB II/III)
- Mehrbelastungen infolge Transformationspfaden werden durch ergänzende Leistungen temporär abgefedert; die Klimadividende wird begünstigend berücksichtigt (Anrechnungsregeln per VO im Einklang mit Art. 8).

- Härtefallfonds für Regionen mit clusterartigen Beschäftigungsverschiebungen; Kofinanzierung aus KKR/TFK (vgl. 17.6/17.2).

- Vermittlung \& Matching (SGB III; eIDAS-Anbindung)
- Zentrale Matching-Plattform (Kompetenz-Badges, eIDAS-Signaturen) verknüpft Module mit Bedarf TFK-/KKR-geförderter Projekte; Datenübermittlung Once-Only; Datenschutz nach 17.1/DS-VO.

- Monitoring \& Evaluation (17.14)
- Jährlicher Bericht (BA/Agentur): Teilnahme, Abschlussquoten, Wiedereingliederung, Einkommenseffekte, regionale Verteilung; Veröffentlichung im Ampel-Reporting.

- Experimentierklausel für Reallabore (WissZeitVG/FuEuOG) mit Sunset (24 Monate) und Evaluation.

- Übergang \& Finanzierung
- Gestufte Einführung binnen 12 Monaten; TKUG startet als Pilot in besonders betroffenen Branchen.

- Finanzierung: Regelhaushalt SGB-Leistungen; ergänzend TFK/KKR für Kurskosten, Matching-Plattform, regionale Härtefälle (vgl. 17.2/17.6).

- Verordnungsermächtigung (SA-VO)
- Qualifikationskataloge/Badges, TKUG-Quote/Höhen, Anerkennung non-formaler Kompetenzen, Matching-APIs (IDs nach 17.1), Hitzeschutz-Mindeststandards, Berichtsschemata (17.14).

{\footnotesize\begin{quote}
Leitidee: „Qualifizieren statt kompensieren“ -- sozial fair, rechtsklar, digital schlank (Once-Only) und messbar.
\end{quote}}



%% ----------------------------------------
\section[Artikel 17.8: Teil H -- Landwirtschaft (Rahmen, LS)]{Artikel 17.8: Teil H -- Landwirtschaft (Rahmen, LS)}
\noindent\textit{ID: 61 -- art-17-8-teil-h-ls (meta)}

\begin{quote}
Integration landwirtschaftlicher Treibhausgaswirkungen (insb. CH4/N2O, Boden-C) in das KlimaGG. Leitplanken im Stammgesetz; technische Einzelheiten durch Landwirtschafts-MRV-Verordnung (LMRV) und ein fortzuschreibendes MRV-Handbuch.
\end{quote}

- Geltungsbereich \& Linien
- Erfassung landwirtschaftlicher Maßnahmen unter den Carbon-Reward-Linien: A (Vermeidung/Substitution), B (Entnahme/Speicherung) und C (Nutzungsumstellung/Stilllegung).

- Die LMRV legt den förderfähigen Maßnahmenkatalog fest (u. a. Güllelager-/Ausbringtechnik, Fütterungsmanagement, Moorwiedervernässung, Agroforst, Humusaufbau, Biochar/Boden-C, langlebige HWP, Fruchtfolge/Leguminosen, Präzisionsdüngung).

- Grundsätze
- Zusätzlichkeit: nur überobligatorische Leistungen jenseits unions-/nationaler Mindeststandards; Stacking mit CAP-Eco-Schemes/AUKM ist nur für zusätzliche Wirkkomponenten zulässig (Mehrleistungs-Bepreisung).

- Doppelzählungssperre: eine Tonne -- ein Reward/Claim; Register-Sperr-/Abgleichmechanismen verbindlich.

- Do-No-Harm: Biodiversität, Boden-/Gewässer- und Gesundheitsschutz dürfen nicht verschlechtert werden; Dünge-, Tierwohl-, Pflanzenschutz- und Gewässerschutzrecht bleiben unberührt.

- Systemkohärenz: ETS/ETS2-, CAP- und Fachregime bleiben unberührt; Wechselwirkungen regelt die LMRV.

- MRV-Leitplanken
- Baselines, Mess-/Schätz-Hierarchien (Tier 2/3), Datenfelder, Prüf-/QS-Stufen und Unsicherheitsgrenzen legt die LMRV fest; IPCC-Guidelines werden berücksichtigt.

- MRV-Handbuch Landwirtschaft als Bestandteil der LMRV; Fortschreibung mindestens zweijährlich (Open Data).

- Default-Faktoren (u. a. CH4 enterisch, N2O aus Düngung) mit Option auf farm-spezifische Faktoren bei Audit-fähigem Nachweis.

- Assurance-Level: AL-0 (vorläufig), AL-1 (prüfbereit), AL-2 (geprüft/limited), AL-3 (reasonable). Großprojekte ab Schwelle [VO] mind. AL-3; jährliche Re-Verifizierung, stichprobenbasiert ≥ 5 \% Volumen/Jahr.

- Permanenz \& Reversal
- Permanenzklassen: D100 (≥100 J.), D40 (≥40 J.), D-A (agrarische, managementgebundene Speicher mit jährlicher Verifizierung).

- Risiko-Buffer (5--20 \%, risikobasiert) und Versicherungsoptionen; Reversals → Registry-Korrektur, Clawback und Sanktionen nach der allgemeinen Sanktionsnorm dieses Gesetzes.

- Leakage \& Integrität
- Aktivitäts- und Markt-Leakage werden durch standardisierte Abschläge/Sicherungen adressiert; sektorale Templates veröffentlicht die Agentur.

- Keine Anrechnung von Wirkungen, die ausschließlich auf Produktionsverlagerung beruhen.

- Daten \& Schnittstellen
- Once-Only mit Anbindung an die Klima-Datenplattform (Art. 17.14) und bestehende Agrarsysteme (z. B. INVeKoS/IACS, Tierregister) über die LMRV.

- Identifikatoren nach Art. 17.1 (Projekt-/Serien-ID, Mapping auf Betriebs-/Flächen-IDs, DPP-Referenzen); Personenbezug minimal, Datenschutz gemäß DS-VO.

- Beihilfe-/EU-Konformität
- Beihilferechtliche Einordnung nach Maßnahme (u. a. ABER/GBER/CEEAG); Anwendung des Beihilfe-Gates gemäß Art. 17.2.

- KMU-Erleichterungen \& Zugang
- Vereinfachte Standardpfade (Sammel-MRV, Pauschalen) für Kleinbetriebe bis [VO] ha/VE; Beratung über Safe-Harbor-Kits der Klima-Serviceplattform (Art. 8a).

- Transaktionskosten-Caps und Sammelverträge (Genossenschaften/Erzeugergemeinschaften) zulässig.

- Verordnungsermächtigung (LMRV)
- Maßnahmenkataloge je Linie, Baselines/Methoden, Prüfanforderungen (Assurance), Permanenz/Buffer/Reversal, Leakage-Regeln, Datenmodelle/APIs, CAP-Stacking-Leitfäden, De-minimis/KMU-Pfade, Pilot-/Fast-Track, Evaluation.

- Frist: Erlass binnen 12 Monaten; Pilotregeln können vorläufig gelten; Evaluation zweijährlich.

- Übergang \& Evaluation
- Pilotphase 12--24 Monate mit begrenzten Volumina; Bestandsvintages unberührt.

- Wirkungs-/Verteilungseffekte jährlich im Ampel-Reporting (Art. 17.14); Anpassung der LMRV nach Evaluation.

{\footnotesize\begin{quote}
Stammgesetz setzt Leitplanken; Details flexibel per LMRV/MRV-Handbuch. Schnittstellen statt Doppelmeldungen, CAP-kompatibel, KMU-tauglich.
\end{quote}}



%% ----------------------------------------
\section[Artikel 17.9: Teil I -- Chemische Industrie (Rahmen, CH)]{Artikel 17.9: Teil I -- Chemische Industrie (Rahmen, CH)}
\noindent\textit{ID: 62 -- art-17-9-teil-i-ch (meta)}

\begin{quote}
Technologieoffene Integration von Prozess- und Rohstoffemissionen (u. a. NH₃, Methanol, Olefine, Chlor-Alkali, N\$\_2\$O/Prozess-CO\$\_2\$) in das KlimaGG. Leitplanken im Stammgesetz; Detailregeln über die Chemie-MRV-Verordnung (CMRV) und ein fortzuschreibendes MRV-Handbuch.
\end{quote}

- Geltungsbereich \& Linien
- Förderfähige Klimaleistungen werden unter den Carbon-Reward-Linien erfasst: A (Vermeidung/Substitution, z. B. E-Cracker, H\$\_2\$-Substitution, Prozesswärme-Elektrifizierung), B (Entnahme/Speicherung, z. B. CCS/Carbonatierung/Polymer-Bindung mit Dauerhaftigkeitsklasse) und C (Stilllegung/Cluster-Umstellung).

- ETS-Pflichten bleiben unberührt; Rewards sind nur für zusätzliche Minderungen über dem Compliance-Niveau zulässig.

- Grundsätze
- Zusätzlichkeit \& Doppelzählung: nach Art. 17.1; Register-Sperrlogik verbindlich (eine Tonne -- ein Claim).

- Systemkohärenz: Schnittstellen zu ETS/ETS2/CBAM, Energie- und Wasserstoffregimen bleiben gewahrt; CBAM-Fragen richten sich nach Art. 17.3.

- Do-No-Harm: Sicherheits-, Umwelt-, Arbeits- und Stoffrecht (u. a. IED/PRTR) bleiben unberührt.

- Systemgrenzen \& Baselines
- Systemgrenze: anlagen- oder clusterbezogen gate-to-gate inkl. relevanter Energieträger/Prozessgase; Elektrizitätsfaktoren nach Strom-Modul der Carbon Rewards (zeitaufgelöst „RT“ oder durchschnittlich „AVG“).

- Baseline-Bildung: produktspezifisch (tCO\$\_2\$e/t Produkt) auf BAT/ETS-Benchmarks bzw. verifizierte historische Performance; Kapazitäts-/Mix-Änderungen werden normalisiert.

- Freizuteilung: Rewards netten ex post um dokumentierte ETS-Freizuteilungs-Effekte, um Doppelvergütung zu vermeiden (CMRV-Formel).

- MRV-Leitplanken \& Daten
- Stack-/Prozessmessung (CO\$\_2\$, N\$\_2\$O), Stoffstrom- und Energiebilanzen; Prüflevels (limited/reasonable) und Unsicherheitsgrenzen legt die CMRV fest.

- Once-Only: ETS-MRV (DEHSt) wird über Schnittstellen angebunden; Mehrfachmeldungen sind ausgeschlossen (Art. 17.1/17.14).

- Produkt-Claims werden über den Digitalen Produktpass (DPP) abgebildet (Vintage, Linie/Klasse, Menge) -- keine parallele Nutzung ohne Register-Retirement.

- Assurance \& ETS-Netting: Großanlagen mind. AL-3 (reasonable); Re-Verifizierung ≥ 5 \% Volumen/Jahr; Freizuteilungs-Netting ex post gemäß CMRV-Formel.

- CCS/CCU -- Dauerhaftigkeit \& Reversal
- Dauerhaftigkeitsklassen: D100 (geologische Speicherung/Carbonatierung ≥100 J.), D40 (Polymer-/Mineral-Bindung ≥40 J.), D-A (anwendungsgebunden, jährliche Verifizierung).

- Risiko-Buffer/Versicherung (5--20 \%, risikobasiert); Reversals ⇒ Registry-Korrektur, Clawback und Sanktionen nach der allgemeinen Sanktionsnorm.

- CCU in kurzlebigen Energieträgern (E-Fuels) ist nicht dauerhaft reward-fähig; anrechenbar sind nur nettopositive Bindungen nach CMRV.

- Elektrifizierung \& flexible Lasten
- Elektrifizierung von Crackern, Dampf/Prozesswärme und Elektrolyse wird nach dem Strom-Modul (RT/AVG, FLEX) bilanziert; Doppelerfassung mit Linie A wird ausgeschlossen.

- Netzdienliche Verfügbarkeiten (NDZ) können als Zusatz-Reward genutzt werden, sofern keine Doppelvergütung entsteht.

- Cluster-Ansätze \& Infrastruktur
- Cluster-MRV (H\$\_2\$-Backbone, CO\$\_2\$-Sammlung/Transport, Dampf-Netze) ist zulässig; Allokation der Rewards nach verursachungsgerechten Schlüsseln (CMRV-Schema).

- Gemeinsame Anlagen dürfen Sammel-Projekte bilden; Haftungs-/Reversal-Regeln definieren anteilige Verantwortung.

- Beihilfe \& EU-Gate
- Beihilferelevante Komponenten unterliegen dem EU-Beihilfe-Gate (CEEAG/GBER); Prüfung nach Art. 17.2.

- Export/Import-Schnittstellen richten sich nach Art. 17.3; Rewards sind kein „carbon price effectively paid“.

- Verordnungsermächtigung (CMRV)
- Maßnahmenkataloge, Baselines/Methoden, ETS-Netting, Prüfanforderungen, Dauerhaftigkeits-/Buffer-Regeln, Cluster-Allokation, Datenmodelle/APIs, Pilotierung und Evaluation.

- Frist: Erlass binnen 12 Monaten; Pilotbetrieb bis 24 Monate möglich; Versionierung als Open Data.

- Übergang \& Evaluation
- 12--24 Monate Pilot mit begrenzten Volumina; Bestandsvintages unberührt.

- Jährliche Wirkungsevaluation (Kosten-/Emissions-Intensitäten, Wettbewerbsfähigkeit, Verteilungswirkungen) im Ampel-Reporting nach Art. 17.14.

{\footnotesize\begin{quote}
Kern: ETS-kompatibel, clusterfähig, MRV-sauber -- Zusatzwirkung wird vergütet, Doppelzählung ausgeschlossen.
\end{quote}}



%% ----------------------------------------
\section[Artikel 17.10: Teil J -- Luftfahrt (AV)]{Artikel 17.10: Teil J -- Luftfahrt (AV)}
\noindent\textit{ID: 63 -- art-17-10-teil-j-av (meta)}

\begin{quote}
Klimaleistungen in der Luftfahrt (SAF, Elektrifizierung/H\$\_2\$, Airside-Prozesse, operative Effizienz) werden einheitlich in das KlimaGG integriert. Leitplanken im Stammgesetz; Detailregeln über die Aviation-MRV-Verordnung (AMRV) und ein fortzuschreibendes MRV-Handbuch.
\end{quote}

- Geltungsbereich \& Linien
- Linie A -- Vermeidung/Substitution: u. a. nachhaltige Flugkraftstoffe (SAF) oberhalb unionsrechtlicher Pflichtquoten, Elektrifizierung von APU/Pushback/Taxi, Bodenprozesse (GPU/PCA), Effizienzmaßnahmen mit belastbarer Minderung.

- Linie B -- Entnahme/Speicherung: anerkannte CO\$\_2\$-Entnahmen (z. B. DAC/BECCS) zur Abdeckung verbleibender Emissionen mit Dauerhaftigkeitsklasse (D100/D40) gemäß Stammgesetz.

- Linie C -- Staatlich verhandelt: nur für öffentlich verhandelte Programme (z. B. strukturelle Slot-/Routen-Stilllegungen, Hub-Neuordnung) gemäß Grundsystem; private Einzelmaßnahmen fallen nicht unter C.

- Kompatibilität (EU-ETS, CORSIA, ReFuelEU)
- Keine Doppelanrechnung: Rewards nur für Mehrmengen über unionsrechtliche Pflichtniveaus (ReFuelEU) und über ETS/CORSIA-Compliance hinaus; parallele Nutzung derselben Tonne ist ausgeschlossen.

- Freizuteilungs-/Compliance-Effekte werden bei der Reward-Berechnung ex post angerechnet (Netting), um Doppelvergütung zu vermeiden.

- CBAM-Fragen sind nicht einschlägig; Außenwirtschafts-/Exportclaims richten sich nach Art. 17.3.

- MRV-Leitplanken \& Daten
- SAF-Nachweis: Chain-of-Custody/Mass-Balance nach unionsrechtlichen Spezifikationen; Anerkennung nur bei verifizierter Nachhaltigkeit, Doppelzählungssperre und dokumentierter Übererfüllung der Quote.

- Effizienz/Operation: Baselines je Flotte/Route (z. B. gCO\$\_2\$e/RPK) mit konservativen Korrekturen für Auslastung/Wetter; nur zusätzliche strukturelle Effekte reward-fähig (kein „business-as-usual“).

- Airside-Elektrifizierung: Strom-Emissionen nach dem Strom-Modul (AVG/RT/FLEX); Doppelvergütung mit Linie A ausgeschlossen.

- Once-Only: ETS/CORSIA/ReFuelEU-MRV werden über Schnittstellen angebunden; Mehrfachmeldungen unzulässig (Art. 17.1/17.14).

- Assurance-Level: SAF-Mehrmengen AL-3; operative Einsparungen mind. AL-2 mit konservativen Korrekturen; Re-Verifizierung ≥ 5 \% Volumen/Jahr.

- Nicht-CO\$\_2\$-Effekte (Kontrails) -- Pilot
- AMRV kann einen vorsichtigen Pilot für methodisch robuste Kontrail-Vermeidung vorsehen (z. B. prädiktive Routenanpassung über Eiskristall-Zonen) mit konservativer Bepreisungsäquivalenz, strengen Unsicherheitsgrenzen und externer Verifizierung.

- Kein Produkt-„Offset“: Nutzung nur als Contribution-Claim nach Art. 17.5; parallele Nutzung in Kompensationssystemen verboten.

- Permanenz \& Reversal (Linie B)
- Dauerhaftigkeit: D100 (geologisch/mineralisch ≥100 J.) voll reward-fähig; D40 (≥40 J.) mit Risiko-Buffer/Versicherung gemäß Stammgesetz.

- Reversals führen zu Registry-Korrektur und Clawback; Sanktionen nach den allgemeinen Sanktionsregeln des KlimaGG.

- Cluster \& Infrastruktur
- Flughafen-Cluster (SAF-Logistik, H\$\_2\$-Betankung, Airside-Netze) können als Sammelprojekte geführt werden; Allokation der Mengen nach verursachungsgerechten Schlüsseln (AMRV-Schema).

- Netzdienliche Verfügbarkeiten (z. B. Lastflex einzelner Airside-Anlagen) können als NDZ-Zusatz vergütet werden, sofern keine Doppelvergütung entsteht.

- Verbraucher-Claims \& DPP
- Produktbezogene Aussagen (z. B. pro Ticket) folgen Art. 17.5: keine „klimaneutral“-Aussagen für laufende Emissionen; zulässig sind DPP-gestützte Produktwerte und Contribution-Claims nach Retirement im Register.

- Bei optionalen Reward-Bundles für Passagiere: separate Preisstellung, Register-Zuweisung/Sperre, kostenfreie Rückabwicklung bei Widerruf (Art. 17.5).

- Beihilfe \& EU-Gate
- Beihilferelevante Komponenten (z. B. SAF-Mehrmengen) unterliegen dem EU-Beihilfe-Gate (CEEAG/GBER) nach Art. 17.2.

- Verordnungsermächtigung (AMRV)
- Maßnahmenkataloge, Methoden (SAF-Mass-Balance, Effizienz-Baselines, Airside-Strom), ETS/CORSIA/ReFuelEU-Abgleich, Pilot Nicht-CO\$\_2\$-Effekte, Daten/APIs, Prüflevels, Evaluation.

- Frist: Erlass binnen 12 Monaten; Pilotbetrieb bis 24 Monate; Versionierung als Open Data (Art. 17.14).

- Übergang \& Evaluation
- 12--24 Monate Pilot mit begrenzten Volumina; Bestandsvintages unberührt.

- Jährliche Wirkungsevaluation (Kosten-/Emissionsintensitäten, Wettbewerbsfähigkeit, Verteilungswirkungen) im Ampel-Reporting nach Art. 17.14.

{\footnotesize\begin{quote}
Kern: SAF-Mehrmengen statt Pflichtquoten-Recycling, ETS/CORSIA-sauber, Airside-Elektrifizierung integriert, Kontrails nur pilotiert -- Claims strikt nach 17.5.
\end{quote}}



%% ----------------------------------------
\section[Artikel 17.11: Teil K -- Kreislaufwirtschaft (CE)]{Artikel 17.11: Teil K -- Kreislaufwirtschaft (CE)}
\noindent\textit{ID: 64 -- art-17-11-teil-k-ce (meta)}

\begin{quote}
Kreislaufeffekte (Re-Use, hochwertiges Recycling, Substitution, Lebensdauerverlängerung, Design-for-Circularity) werden standardisiert in das KlimaGG integriert. Leitplanken im Stammgesetz; technische Details über die Kreislauf-MRV-Verordnung (KMRV), ein fortzuschreibendes MRV-Handbuch und die Anbindung an den Digitalen Produktpass (DPP/ESPR).
\end{quote}

- Geltungsbereich \& Linien
- Linie A -- Vermeidung/Substitution: Re-Use, Remanufacturing, hochwertiges Recycling, designbedingte Material-/Produktsubstitution, Lebensdauerverlängerung (Langlebigkeit/Repairability).

- Linie B -- Entnahme/Speicherung: z. B. Biochar in langlebigen Produkten/Materialien mit Dauerhaftigkeitsklasse (D100/D40) nach Stammgesetz.

- Linie C -- Staatlich verhandelt: nur für öffentlich verhandelte Strukturmaßnahmen (z. B. Stilllegung emissionsintensiver Abfallpfade, Cluster-Umstellung); private Einzelmaßnahmen fallen nicht unter C.

- Grundsätze (Art. 17.1, 17.5)
- Zusätzlichkeit: nur Leistungen über unions-/nationalen Mindeststandards (EPR/Abfall-/Produktrecht) hinaus; keine Doppel- oder Parallelförderung.

- Doppelzählungssperre: Registry-Sperrlogik verbindlich; klare Abgrenzung zu EPR-Gutschriften/Handelssystemen.

- Do-No-Harm: keine Verschlechterung von Biodiversität, Boden-/Gewässer- und Gesundheitsschutz.

- Claim-Regeln: Verbraucher-/Produktclaims ausschließlich gem. Art. 17.5 (keine pauschalen „klimaneutral“-Aussagen).

- MRV-Leitplanken \& Berechnungslogik (KMRV)
- Baseline \& Substitution: tCO\$\_2\$e = (EFprim − EFsek) × M × Q − Eproc, wobei EFprim/sek produkt-/materialspezifische Primär-/Sekundärfaktoren sind, M die Masse/Einheiten und Q ein Qualitäts-/Downcycling-/Lebensdauerfaktor; Eproc sind Prozess-/Energieemissionen.

- Lebensdauerverlängerung: Bilanz über vermeidenen Ersatz mit konservativen Nutzungsdauern, Reparaturquoten und Rebound-Korrektur.

- Massenbilanz/CoC: verpflichtende Massenbilanz/Chain-of-Custody; Herkunft, Qualität und Energieprofile nachzuweisen.

- Prüflevels: gestufte Assurance (limited/reasonable) nach Projektgröße/Risiko; Stichprobenpflicht.

- Assurance-Schwellen: Materialströme ≥ [VO] t/Jahr obligatorisch AL-3; Downcycling-Faktor Q nur mit AL-2+.

- Ausschlüsse \& Anti-Gaming
- Waste-to-Energy (WtE): thermische Verwertung/Fossilanteile sind grundsätzlich nicht reward-fähig; Ausnahmen nur bei klarer systemischer Netto-Minderung ggü. rechtsverbindlicher Baseline und ohne Doppelanrechnung.

- Up-/Downcycling: Downcycling wird über Q abgeschlagen; nur hochwertige Kreisläufe erhalten vollen Effekt.

- Vorverlagerung: keine Gutschrift für bloße Verlagerung in Drittstaaten; Export-/Leakage-Sperren gem. Art. 17.3.

- DPP/ESPR-Verknüpfung \& Daten (Art. 17.14)
- Pflichtfelder aus DPP/ESPR (z. B. Materialmix, Rezyklatanteil, Reparierbarkeit, Haltbarkeit) werden für MRV genutzt (Once-Only, API).

- DPP-Kernelemente (Hash/Signatur) im Kauf-/Projektbeleg; Datenportabilität und Open-Data-Aggregate gem. 17.14.

- Permanenz \& Reversal (Linie B)
- Dauerhaftigkeitsklassen D100/D40/D-A mit Buffer/Versicherungsmechanismen; Reversals ⇒ Registry-Korrektur und Clawback nach allgemeinem Sanktionsrahmen.

- Schnittstellen zu Abfall-/Produktrecht \& Export (Art. 17.3)
- EPR/Abfallrecht (u. a. Verpackung, ElektroG/BattG) bleibt unberührt; Rewards nur zusätzlich und ohne Doppelgutschrift.

- Grenzüberschreitende Stoffströme: Anerkennung nur bei MRV-/Registry-Äquivalenz; Default-Werte und Sperrregeln per KMRV; Freeport-Leakage ausgeschlossen.

- Beihilfe \& EU-Gate (Art. 17.2)
- Beihilferelevante Komponenten laufen über das EU-Beihilfe-Gate (CEEAG/GBER); Transparenz im Ampel-Reporting.

- Verordnungsermächtigung (KMRV)
- Maßnahmenkataloge je Produkt-/Materialgruppe, Faktoren-Tabellen (EFprim/sek, Q), MRV-Methoden/Assurance, DPP-/API-Profile, Export-Äquivalenz/Default-Tabellen, Pilot/Fast-Track, Evaluation.

- Erlass binnen 12 Monaten; Pilotbetrieb bis 24 Monate; Versionierung als Open Data (17.14).

- Übergang \& Monitoring
- 12--24 Monate Pilot mit priorisierten Materialströmen (z. B. Stahl, Alu, Kunststoffe, Batterien, Baustoffe); Bestandsvintages unberührt.

- Jährliche Wirkungsevaluation (€/tCO\$\_2\$, Materialkreislaufquoten, Qualitätsgrade) im Ampel-Reporting nach 17.14.

{\footnotesize\begin{quote}
Hochwertige Kreisläufe werden belohnt, Greenwashing und Doppelgutschriften verhindert -- mit DPP-Daten, Export-Sperren und klarer MRV-Formel.
\end{quote}}



%% ----------------------------------------
\section[Artikel 17.12: Teil L -- Vertrauensdienste, eIDAS \& Digitale Identitäten (ID)]{Artikel 17.12: Teil L -- Vertrauensdienste, eIDAS \& Digitale Identitäten (ID)}
\noindent\textit{ID: 65 -- art-17-12-teil-l-id (meta)}

\begin{quote}
EU-weit anerkannte Identifizierung, Signatur, Siegel, Zeitstempel und Zustellung für alle KlimaGG-Plattformen (Register, Klima-Datenplattform, Vergabe-/PPA-Zuweisungen, Bescheide). Integration der EU Digital Identity Wallets; starke Beweis- und Vollzugswirkung bei zugleich barrierefreien Fallback-Optionen.
\end{quote}

- Anwendungsbereich \& Verhältnis
- Gilt für Onboarding/KYC, Registry-Vorgänge (Ausgabe/Issuance, Transfer, Sperre/Unlock, Annullierung, Einlösung/Redeem), Zuweisungen in Vergaben/PPAs, Bescheide/Sanktionen, Offenlegungen/Reports sowie ADR-Vergleiche.

- Privatrechtliche Verträge können die Dienste nutzen; öffentlich-rechtliche Akte mit Registerbezug müssen sie nutzen. Querschnittsregeln aus 17.1 (IDs/Once-Only), 17.4 (Vollzug) und 17.14 (Datenplattform) bleiben unberührt.

- Identifizierung \& Rollen (eIDAS/EU-Wallet)
- Zugelassen sind notifizierte eIDs und EU Digital Identity Wallets; Sicherheitsniveau: hoch für juristische Personen/Rollen mit Vermögens-/Registerwirkung, substanziell für Retail-Nutzer ohne Drittwirkung; risikoadäquate Anhebung per VO möglich.

- Rollen: Inhaber:in, Vertretung/Bevollmächtigte:r, Prüfer:in/Assurance, Register-Operator, Behördenrolle; Rollen-Nachweise als Wallet-Credentials (organisationsgebunden, widerrufbar).

- Alternativen (Video-Ident/Vor-Ort) sind zulässig; Ergebnis wird in ein qualifiziertes Ident-Token überführt.

- Signatur, Siegel, Zeitstempel, Zustellung
- QES ist erforderlich für rechtsbegründende/wirkändernde Registry-Akte (Ausgabe, Eigentumsübertrag, Redeem, ADR-Vergleich); PAdES/XAdES mit Long-Term-Validation (LTV) verpflichtend.

- QSeal der Registerstelle für Serienausgaben, amtliche Registry-Auszüge, Referenzpreis-/Offenlegungsdokumente, Verwaltungsakte und Open-Data-Snapshots.

- Qualifizierte Zeitstempel (UTC) für alle statusändernden Events; Append-Only-Eventlog.

- QERDS (qualifizierter elektronischer Zustelldienst) für Bescheide/Fristläufe; Fallback: postalisch mit Register-Verknüpfung.

- Interoperabilität \& Kennungen
- Verbindliche Mappings auf LEI/ISIN/EORI/HS sowie DPP-Referenzen; technische Profile über 17.14 (APIs/Webhooks, Fehlercodes).

- Transport-/API-Sicherheit: QWACs, mTLS, OIDC mit rollenbasierten Claims; HSM-gestütztes Schlüsselmanagement, Vier-Augen-Prinzip für kritische Aktionen.

- API-Governance \& Versionierung
- SemVer mit Parallellaufzeit je Minor ≥ 18 Monate; Deprecation-Notices ≥ 6 Monate; versionierte Pfade (/vX/...).

- Conformance-Sandbox mit Pflicht-Testfällen; QSeal-signiertes Conformance-Zertifikat als Go-Live-Voraussetzung.

- Change Advisory Board (Registerstelle, Aufsichten, Ländervertretung) mit Beschlussfrist 10 AT; Protokolle öffentlich.

- Delegation \& Vertretung
- Digitale Vollmachten (Wallet-Credential/QES) mit Umfang, Dauer, Widerruf; Registry prüft kryptographisch und protokolliert Widerrufe mit Zeitstempel.

- Unterbevollmächtigungen nur, wenn die Ausgangsvollmacht dies vorsieht; Transparenz im Eventlog.

- AML/KYC \& Drittstaaten
- Risikoadäquate AML/KYC-Prüfungen ohne Doppelmeldungen (Once-Only über 17.14); Nutzung anerkannter Vertrauens-/Identanbieter. Außenwirtschaft/Embargo-Prüfungen nach 17.3, kein Widerspruch.

- Drittstaaten-Identitäten/-Vertrauensdienste nur bei gleichwertigem Schutzniveau (White-List-VO) oder gesichertem Fallback.

- Barrierefreiheit \& Fallback
- Barrierefreie Verfahren; analoge Zugänge (Bürgerbüro/Post-Ident) mit späterer Tokenisierung in EU-Wallet/eID; keine Benachteiligung durch Offline-Nutzung.

- Betrieb \& SLOs
- Monatliche SLOs: API-Uptime ≥ 99,5 \%, MTTR ≤ 8 h; Veröffentlichung über AR-Open-Data (/open-data/ops).

- Bei Unterschreitung automatische Feature-Freeze bis Stabilisierung; Paper/QERDS-Fallback aktiv.

- Datenschutz \& Aufbewahrung
- Datenminimierung, Zweckbindung; personenbezogene Daten pseudonymisiert in Offenlegungen; Audit-Logs manipulationssicher.

- Aufbewahrung: max. 10 Jahre nach letzter Einlösung/Bestandskraft, soweit spezialgesetzlich nichts anderes angeordnet ist.

- Durchsetzung \& Sanktionen
- Missbrauch/Möchtegern-Identitäten/Signaturfälschung ⇒ Sperrvermerke, Clawback, Bußgelder und Ausschluss vom Fast-Track nach allgemeinem Sanktionsrahmen des Stammgesetzes; Meldung an Aufsichten.

- Übergang \& Evaluation
- Übergangsphase: bis Monat 12 sind AdES + qualifizierter Zeitstempel zulässig; ab Monat 12 QES/QSeal verpflichtend für Registry-Akte. Wallet-Roll-out stufenweise.

- Evaluation nach 24 Monaten (Nutzbarkeit, Sicherheit, Interoperabilität); Anpassungen per VO.

- Verordnungsermächtigung (ID-VO)
- Technische Profile/Prozesse: akzeptierte eID/Wallet-Schemata, QES/QSeal/Time-Stamp/ERDS-Anbieterlisten, LTV-Anforderungen, Onboarding-Flows, Delegationsformate, API/Protokolle, Schwellen für Sicherheitsniveaus.

{\footnotesize\begin{quote}
Hohe Beweiswirkung bei geringer Reibung: Wallet-basierte Rollen, LTV-sichere Signaturen, QERDS-Zustellung -- und für alle ohne eID ein sauberer Fallback.
\end{quote}}



%% ----------------------------------------
\section[VO 17.12a: Verordnung zu Vertrauensdiensten, eIDAS \& Digitalen Identitäten im KlimaGG (ID-VO)]{VO 17.12a: Verordnung zu Vertrauensdiensten, eIDAS \& Digitalen Identitäten im KlimaGG (ID-VO)}
\noindent\textit{ID: 66 -- art-17-12-id-vo (meta)}

\begin{quote}
Konkretisiert Art. 17.12 (ID): Onboarding/Identifizierung, Signatur-/Siegelanforderungen, qualifizierte Zeitstempel, Zustellung, Transport-/Web-Sicherheit, Rollen-/Rechtemanagement, API-/Protokollprofile, Interoperabilität sowie Übergang. Querschnittsbezug: 17.1 (Begriffe/IDs/DS), 17.14 (Datenplattform/Monitoring), Verbraucherrechte/Claims: 17.5.
\end{quote}

- Geltung, Akteure \& Rollen
- Gilt für rechtserhebliche Vorgänge im KlimaGG-Register: Issuance, Transfer, Redeem, Sperre/Löschung, Prüf-/Auditvorgänge, Sanktions-/Clawback-Akte, Open-Data-Snapshots, PPA-/Vergabe-Zuweisungen.

- Akteure: Registerstelle, Emittenten/Projektträger, Prüfer/Auditoren, Verwahrer/Depotstellen, Behörden (u. a. Energie/Zoll/Beihilfe), Aggregatoren, Lieferanten/DSOs, Verbraucher.

- Rollen als Wallet-Credentials: registry-issuer, registry-verifier, registry-auditor, registry-custodian, registry-operator, participant.

- Identifizierung \& Onboarding
- Akzeptiert: notifizierte eIDs und EU Digital Identity Wallets (Sicherheitsniveau: hoch für juristische Personen/Institutionsrollen; substanziell für Retail, risikobasiert auf hoch anhebbar).

- Alternativen: VideoIdent/Vor-Ort-Ident; LEI/EORI-Mapping; UBO-Erhebung nach GwG.

- Delegation: organisationsbezogene VCs (signatory/officer), prüfbar; Nachweis/Bindung via QERDS.

- Signaturen \& Siegel
- QES-Pflicht für: Issuance-Anträge, Transfers, Redeems, ADR-Vergleiche/Clawbacks, kritische Rollenänderungen.

- QSeal-Pflicht der Registerstelle für: Serienausgaben, amtliche Registry-Auszüge, Referenzpreis-/Offenlegungsdokumente, Sanktions-/Verwaltungsakte, Open-Data-Pakete.

- Langzeitprofile: PAdES/XAdES mit LTA; OCSP/CRL-Belege beifügen.

- Qualifizierte Zeitstempel
- Alle statusändernden Events mit qualifiziertem elektronischem Zeitstempel (UTC) eines EU-listierten TSP.

- Event-IDs monoton/unverwechselbar (z. B. ULID/UUIDv7); Referenz im Registerjournal.

- Zustellung (QERDS)
- Bescheide, Sperr-/Sanktionsverfügungen, Fristmitteilungen via QERDS; Zustellnachweis löst Fristen aus.

- Fallback: hybride sichere Zustellung mit gleichwertiger Fristwirkung für Beteiligte ohne QERDS.

- Website-/Transport-Sicherheit
- QWACs für Web/API; TLS 1.3, mTLS für Partner-APIs, Perfect Forward Secrecy.

- HSM-gestütztes Schlüsselmanagement; Rotation [VO 90--180 Tage]; Vier-Augen-Prinzip.

- BSI-Grundschutz/ISO 27001, NIS2-Pflichten; jährliche Pen-Tests; Incident-Meldung binnen 24 h (vgl. 17.4/17.14).

- APIs, Protokolle \& Events
- AuthN/AuthZ: OIDC/OAuth2, signierte JWT (JWS), Scope-Minimalprinzip; Anti-Replay (nonce), Idempotency-Keys.

- Kern-Endpunkte: /v1/events, /v1/issuance, /v1/transfer, /v1/redeem, /v1/claims, /v1/keys/jwks, /v1/eidas/status.

- Event-Hülle: QSeal-signiert mit qualifiziertem Zeitstempel und Referenzen (Projekt-ID, Serien-ID, LEI/EORI); Detailformate siehe AR-VO (17.14).

\{
  "schema":"ID-VO/EventEnvelope/1.0",
  "event\_id":"01HXY...7F",
  "event\_type":"transfer|issuance|redeem|lock|claim",
  "occurred\_at":"2026-01-20T11:34:15Z",
  "subject\_refs":\{"project\_id":"PRJ-123","series\_id":"SER-456","lei":"529900..."\},
  "actor":\{"wallet\_did":"did:eu:wallet:abc","role":"registry-issuer","qes\_sig":"<base64>"\},
  "payload\_hash":"SHA256:...",
  "qts":"<qualified-timestamp-token>",
  "qseal":"<registry-seal-signature>"
\}

- Rollen, Delegation \& Mehrfachfreigaben
- Kritische Aktionen erfordern Mehrfachfreigaben (z. B. 2-von-3) [VO].

- Rollen-/Rechtejournal unveränderlich; Delegationen befristet und widerrufbar.

- Interoperabilität \& Anerkennung
- Anerkennung aller eIDAS-TSP aus EU-Vertrauensliste; Drittstaaten nur bei dokumentierter Gleichwertigkeit ([VO-Whitelist]), sonst Fallback.

- Re-Validierung der TSP-Profile mind. alle [VO 24] Monate.

- Versionierung \& Deprecation
- API-Versionen nach SemVer; Parallellaufzeit je Minor ≥ 18 Monate; Deprecation-Notice ≥ 6 Monate; Breaking Changes nur mit CAB-Beschluss.

- Versionierte JWKS-Endpunkte; Idempotency-Keys Pflicht für POST auf kritischen Pfaden.

- Conformance \& CAB
- Öffentliche Conformance-Suite; Bestehen der Pflicht-Testfälle als Voraussetzung für produktive Zertifikate.

- Change Advisory Board; Protokollierung und Veröffentlichung der Beschlüsse.

- OPS-Transparenz
- Pflichtfeed /open-data/ops (Uptime, Incidents, MTTR, Releases) -- QSeal-gesiegelt, monatlich.

- Feature-Freeze bei Unterschreitung der SLOs; Fallback-Verfahren (QERDS/Paper) verpflichtend auszurufen.

- Datenschutz \& Aufbewahrung
- Minimalprinzip/Zweckbindung, Pseudonymisierung; DPIA für Registry-Kernprozesse (Bezug 17.1).

- Aufbewahrung rechtserheblicher Nachweise max. [VO bis 10] Jahre nach letzter Einlösung; danach Löschung/Archiv mit Hash-Beweis.

- Business Continuity \& Fallback
- RTO ≤ [VO 4 h], RPO ≤ [VO 1 h], georedundant; Notfall-Wallet-Wiederherstellung via QERDS-gestützter Identbestätigung.

- Übergang: AdES + qualifizierter Zeitstempel bis Monat 12 zulässig; ab Monat 17 QES/QSeal verpflichtend.

- Sanktionen \& Vollzug
- Verstöße gegen Ident-, Signatur-, Zustell- oder Siegelpflichten sind Ordnungswidrigkeiten; Clawback/weitere Konsequenzen nach dem Sanktionsrahmen des KlimaGG.

- Registerseitige Sperre betroffener Serien/Transaktionen bis zur Heilung; veröffentlichte Aggregatberichte (vgl. 17.14).

- Übergang \& Evaluation
- Stufenplan: M0--M6 Onboarding/Wallets \& QERDS; M7--M12 QES/QSeal-Pflicht; M17 Vollumsetzung.

- Evaluation nach 24 Monaten (Fehlerquoten, Durchlaufzeiten, Interoperabilität); ggf. VO-Anpassung.


%% ----------------------------------------
\section[Artikel 17.13: Teil M -- Energiepreise \& Netzentgelte (EN)]{Artikel 17.13: Teil M -- Energiepreise \& Netzentgelte (EN)}
\noindent\textit{ID: 67 -- art-17-13-teil-m-en (meta)}

\begin{quote}
Ziel: Einheitlicher, transparenter und sozial verträglicher Energiemarktrahmen, der netzdienliche Flexibilität, Sektorkopplung und Investitionen unterstützt. Operative Details stehen in der EN-VO Verordnungsvorlage.
\end{quote}

- Grundsätze \& Verordnungsermächtigung (EN-VO)
- Grundsätze: Kostenreflexivität, Netzdienlichkeit, Technologieoffenheit, Datentransparenz, Sozialverträglichkeit.

- VO-Mandat: Bundesregierung erlässt die EN-VO (Zeitfenster/ToU/Leistung, §14a-Profile S/M/L, dynamische Tarife, Mess-/API-Profile, Preisblatt/Rechnung, Open-Data, Schutzmechanismen). Frist: 12 Monate nach Inkrafttreten.

- Daten \& Sicherheit: Schnittstellen zur Klima-Datenplattform (17.14); NIS2/IT-SiG-Konformität; QES/QSeal für abrechnungsrelevante Ereignisse.

- Market Governance: Lieferanten-/MSB-Pflichten, Fallback-Regeln, Standard-Preisblatt, Veröffentlichungspflichten gem. EN-VO.

- Gesetzesänderungen (Bündel)
- EnWG: Bereinigung obsoleter Förder-/Privilegienverweise; §14a-Brücke zur EN-VO (Standard-Flex S/M/L, Aggregator-Zugang, Nachweis über EN-VO-Logs); Doppelförderung ausgeschlossen (vgl. 17.1/Art. 4); MaStR↔Klima-Datenplattform (17.14, Once-Only).

- EEG: Neuanlagen ausschließlich KlimaGG-Instrumente; Bestandsanlagen Opt-In mit Stichtagen; Rückkehr ausgeschlossen. Doppelförderungsverbot (17.1/17.5); Repowering-Fast-Track; CfD/PPA kompatibel mit Transparenzbericht (17.14).

- KWKG: Neustopp für Neuförderung; Bestandsschutz. Stilllegung/Umrüstung über Carbon Rewards nur bei Additionalität/Permanenz; Jahresbericht „Transformation KWK“.

- EnFG: Auslaufplan binnen 6 Monaten, jährliche Fortschreibung; offene Berichte (Begünstigte, Volumina, Klimawirkung) via 17.14.

- StromStG: Absenkung auf EU-Minimum (Haushalte 0,10 ct/kWh; Industrie 0,05 ct/kWh); jährliche Wirkungsanalyse inkl. Netzentgelt-Interaktionen (17.14).

- KAV: Sunset der verbrauchsgebundenen Konzessionsabgaben (spätestens T+24). Ersatz durch 5 \% Quellenabzug auf alle Carbon-Reward-Auszahlungen: 50 \% Grundanteil (alle Kommunen, Einwohnerschlüssel), 50 \% Projektanteil (georeferenzierte tCO\$\_2\$e am Standort). Bei Projekten mit CR gr. X:  . Keine Umlage auf Letztverbraucher. Kommunalregister \& Quartalsauszahlung über 17.14; kommunale Buchung nach 17.6 (KH).

Übergangs- und Anwendungsdetails konkretisiert die EN-VO; Haushalts-/Kommunalbuchung vgl. 17.6 (KH).

- Netzentgelte \& Regulierungsverordnungen (StromNEV, ARegV, NABEG, NAV, StromNZV, MsTrV)
- Reformmandat BNetzA: Vorlage binnen 12 Monaten im Einklang mit EN-VO-Korridoren.

- Umlagen-Neuzuschnitt: §19 StromNEV, Offshore-Umlage, AbLaV verursachergerecht und -- wo sachgerecht -- alternativ finanziert; keine reine Endkunden-Abwälzung.

- Outcome-KPIs: u. a. Redispatch-Kosten, Lastverschiebung (MW/Slot), Speicherstunden, Härtefallquote; jährlicher BNetzA-Bericht ins Ampel-Reporting (17.14).

- Datenpfad: EN-VO-Preisblätter/Open-Data; NIS2/IT-SiG-Pflichten für Netzdaten-Plattformen.

- Kohleverstromungsbeendigungsgesetz (KVBG)
- Opt-In \& Verzicht: Anlagen-/Cluster-Erklärung; Verzicht auf KohleAusG-Ansprüche; KlimaGG-Förderung nur für zusätzliche Minderungen.

- Systemtest: ÜNB/BNetzA-Prüfung, Netzkosten-Neutralität; Reihung über transparente Auktionen.

- ETS-Zertifikate: Zuweisung an die Agentur; Stilllegung/Löschung nach Kriterienkatalog; Veröffentlichung im Ampel-Reporting (17.14).

- Kontrolle: Doppelförderungs- und Integritätsprüfungen; Sanktionen nach allgemeinem Sanktionsrahmen des Stammgesetzes.

- Übergang, Kohärenz \& Evaluation
- Meilensteine: EN-VO T+12; Netzentgelt-Reformvorschlag BNetzA T+12; KAV-Sunset T+24; StromStG-Absenkung ab T+1.

- Kohärenz: Doppelförderungs-Sperren (17.1/17.5), DV-Fristen (17.4), Open-Data/Monitoring (17.14), Kommunalbuchung (17.6).

- Evaluation: Quartals-Kurzberichte und jährlicher Detailbericht mit Ampel-Triggern (17.14); ggf. Anpassungs-VO binnen 6 Monaten nach Bericht.

{\footnotesize\begin{quote}
- Konzentration statt Zersplitterung: Ein EN-Artikel als normative Klammer; Detailvollzug in der EN-VO.

- Referenzanker: Einzelgesetz-Änderungen (EnWG/EEG/KWKG/EnFG/StromStG/KAV) sind hier gebündelt und verweisen operativ auf die EN-VO.
\end{quote}}



%% ----------------------------------------
\section[VO 17.13a: Verordnung zu Energiepreisen \& Netzentgelten (EN-VO)]{VO 17.13a: Verordnung zu Energiepreisen \& Netzentgelten (EN-VO)}
\noindent\textit{ID: 68 -- art-17-13-en-vo (meta)}

\begin{quote}
Konkretisiert Art. 17.13 (EN) schlank und parametrisierbar: ToU-Netzentgelte, §14a-Flex-Profile, dynamische Tarife, Mindest-Messanforderungen, Open-Data-Pflichten, Sozialmechanismen (nicht über Netzentgelte finanziert) sowie Übergang \& Evaluation. Querschnitt (IDs/eIDAS/Monitoring) richtet sich nach 17.1, 17.12, 17.14.
\end{quote}

- Geltung, Begriffe, Verweisrahmen
- Gilt für ÜNB/VNB, Lieferanten, Messstellenbetreiber und Letztverbraucher (Strom).

- Begriffe nach Art. 17.13; Identitäten/Signaturen/Zeitstempel nach 17.1/17.12; Daten-/Reporting-Schnittstellen nach 17.14.

- NIS2/IT-SiG, Datenschutz und Once-Only gelten querschnittlich (17.12/17.14) und werden hier nicht wiederholt.

- Zeitvariable Netzentgelte (ToU)
- Mindestens ein ToU-Modell (Peak/Off-Peak/Super-Off-Peak); optional Kapazitäts-/Leistungsteil (LP).

- AP-Multiplikatoren: Peak [VO 1.2--2.5×], Off-Peak [VO 0.6--0.9×], Super-Off-Peak [VO 0.3--0.6×].

- Fenster/Zonen DSO-spezifisch, jährlich fortzuschreiben; Veröffentlichung als Open Data (Schema: AR-VO).

- Messanforderung: mind. 15-min-Fähigkeit für ToU/LP; Übergänge/Fallbacks [VO] für Nicht-SMGW-Haushalte.

- Diskriminierungsverbot; Zonen-Mehr/Minderentgelte nur nach [VO-Methodik].

- §14a EnWG: Standardisierte Flex-Profile S/M/L
- Begünstigte Entgelte gegen netzdienliche Steuerbarkeit; Aggregatoren erhalten diskriminierungsfreien Zugang.

- 

ProfilReaktionszeitMin. DauerVerfügbarkeit/TagMax. Abschaltzeit

S≤ 15 min30 min≥ 6 h (Fenster)≤ 60 min/Tag; ≤ 10 h/Monat
M≤ 5 min15 min≥ 10 h≤ 120 min/Tag; ≤ 20 h/Monat
L≤ 60 s5 min24/7 (Ausnahmen [VO])≤ 240 min/Tag; ≤ 40 h/Monat

- Vergütung/Reduktion: Pay-as-Delivered oder rabattierte NNE-Klasse je Profil [VO \%].

- Schutzprofile: Mindesttemperatur, Medizingeräte; Härtefall-Bypass [VO].

- Nachweis: 15-min-Ereignislogs (QES/QSeal); Bilanzierung/Konfliktlösung gemäß AR-VO.

- Dynamische Lieferverträge \& Lieferantenpflichten
- Mindestens ein dynamischer Tarif (Spot-Index [EPEX DA/ID] oder ToU) für SMGW-fähige Kunden.

- Preisbestandteile T-1 bis 18:00 Uhr; Intraday-Updates [VO]; maschinenlesbar nach AR-VO-Schema.

- Standard-Preisblatt/Rechnung gemäß AR-VO; verpflichtender TCO-Hinweis (CapEx/OpEx/CO\$\_2\$).

- Fallback bei Kommunikationsausfall [VO]; keine Strafklauseln für Tarifwechsel.

- Messstellenbetrieb \& Sicherheit
- Mindestprofile: 15-min-Lastgang, Fernsteuer-Kanal, Ereignis-Signaturen; Latenz ≤ [VO] min.

- APIs/Identitäten/Signaturen: nach 17.12; Transport-/Schnittstellenprofile: AR-VO.

- Datenschutz: Minimalprinzip, Aufbewahrung [VO], Anonymisierung für Open-Data gemäß AR-VO.

- Open Data \& Netzstatus
- Pflichtendpunkte: /netstatus/zonen, /netstatus/fenster, /netstatus/engpaesse, /preise/nne, /§14a/events.

- Aktualität: Zonen/Fenster mind. jährlich; Engpässe mind. täglich; Preise T-1; Formate/Lizenzen nach AR-VO.

- Sozialmechanismen (außerhalb Netzentgelte)
- Zielgruppe [VO Kriterien]; Instrumente: Grundkontingent-Deckel [kWh/a] oder Zuschüsse [VO].

- Finanzierung ausschließlich haushalts-/sozialrechtlich (vgl. 17.2/17.7); keine Refinanzierung über Netzentgelte/Umlagen.

- Zielgenaue Entlastung Industrie
- Adressaten: energieintensive, transformierende Prozesse; MRV-gebunden, Effizienz-/Flex-Pflichten, No-Backsliding.

- Beihilfe-Gate (CEEAG/GBER) zwingend; Export-Rebates nur unions-/WTO-konform; Transparenz im AR-Dashboard.

- Übergang \& Evaluation
- Meilensteine: M0--M6 Open-Data/Preisblätter; M7--M12 §14a-Profile S/M/L; M17--M18 ToU flächig; M19--M24 Umlagen-Bündelung [VO].

- Pilotzonen [VO] mit Monitoring; KPIs/Evaluationsberichte nach AR-VO; ggf. Anpassungs-VO binnen 6 Monaten.


%% ----------------------------------------
\section[Artikel 17.14: Teil N -- Monitoring \& Ampel-Reporting (AR)]{Artikel 17.14: Teil N -- Monitoring \& Ampel-Reporting (AR)}
\noindent\textit{ID: 69 -- art-17-14-teil-n-ar (meta)}

\begin{quote}
Einheitliche, digitale Monitoring-/Reporting-Regeln für alle KlimaGG-Instrumente: schnell, rechtsstaatlich, transparent und revisionssicher. Daten fließen „once-only“; Ampel-Status triggert definierte Nachsteuerungen.
\end{quote}

- Geltung \& Verhältnis
- Gilt für das KlimaGG und harmonisierte Fachgesetze; Spezialbeschleuniger (Art. 8b) gehen vor.

- Querschnittsregeln zu IDs/Once-Only (17.1), Durchsetzung/Vollzug (17.4) und eIDAS/Signaturen (17.12) bleiben unberührt.

- Datenplattform \& Once-Only
- Elektronische Verfahrens-/Datenakte gemäß EGovG; offene Schnittstellen zum Registry (17.1), DPP/Verbraucher (17.5), eIDAS/Wallet (17.12), Außenwirtschaft/Zoll (17.3) und amtlicher Statistik.

- Verbot von Doppelmeldungen; qualifizierte Signaturen/Siegel für statusändernde Ereignisse; versionierte Änderungsjournale.

- Ampel-System \& Trigger
- Definition von Zielkorridoren für u. a. Emissionspfad/Budget, Permanenz-/Reversal-Risiken, Cap-Ausnutzung, Beihilfe-Status, Datenqualität.

- Grün: innerhalb des Korridors. Gelb: Abweichung ≤ [VO] Korridorpunkte. Rot: signifikante Zielverfehlung oder Risikoüberschreitung.

- Rot-Folgen: Maßnahmenplan binnen 30 Tagen; ggf. automatische Nachsteuerung gemäß 17.2 (Tilgung/Cap), Registry-Sperrvermerke und Clawbacks nach Sanktionsrahmen des Stammgesetzes.

- Zusätzliche Trigger: (a) aid-pending > 180 Tage oder > 25 \% der neuen Serien, (b) API-Uptime < 99,5 \%/Monat oder ≥ 2 kritische Incidents, (c) Audit-Fail-Rate ≥ 2 \% zwei Quartale in Folge.

- Berichtskaskade
- Quartal: Kennzahlen zu Krediten/Tilgungen/Rücklagen, Risiken/Einlösungen (CR-Floor), Cap-Ausnutzung, Beihilfe-Status, Netzdaten/§14a-Leistungen; Übermittlung an Stabilitätsrat/Statistik; Veröffentlichung im Dashboard.

- Jahr: Evaluationsbericht (Tempo, Zielerreichung, Rechtsschutzqualität, IT-Sicherheit, Nutzerfreundlichkeit) mit Maßnahmenplan.

- IT-Sicherheit \& BCM
- NIS2/IT-SiG-Pflichten: ISO/IEC 27001/BSI-Grundschutz, jährliche Pen-Tests, HSM-Schlüssel, Vier-Augen-Prinzip für kritische Operationen.

- Business Continuity: RTO ≤ 4 h, RPO ≤ 1 h, Georedundanz, halbjährliche Notfallübungen; Incidents binnen 24 h melden; betroffene Serien/Vorgänge vorsorglich sperren.

- Öffentlichkeit \& Barrierefreiheit
- Öffentliches Ampel-Dashboard (Status, Fristen, Änderungen) maschinenlesbar; barrierefrei, mehrsprachige Kurzfassungen.

- Zusätzlicher OPS-Feed mit Betriebsmetriken (Uptime, MTTR, Release-Log) gemäß ID-VO/AR-VO.

- Open-Data mit Pseudonymisierung; IFG/Transparenzrechte bleiben unberührt.

- Aufsicht \& Qualität
- Halbjährliche Fristen-/Qualitätsaudits; Organisationsverfügungen bei systemischen Verstößen.

- Verfahrensbezogene Verstöße unterfallen dem allgemeinen Sanktions-/Clawback-Rahmen des Stammgesetzes.

- Verordnungsermächtigung (AR-VO)
- Datenmodelle/APIs, Ampel-Schwellen/Korridore, Berichtsschemata, Open-Data-Profile, IT-Sicherheitsprofile und BCM-Vorgaben.

- Übergang
- 12-monatige Pilotphase für digitale Akte \& Schnittstellen; danach volle Anwendung. Altverfahren können optieren.

{\footnotesize\begin{quote}
Ampel schafft Verbindlichkeit ohne Rechtsschutz zu verkürzen; Schnittstellen verhindern Greenwashing und Datenbrüche.
\end{quote}}



%% ----------------------------------------
\section[VO 17.14a: Verordnung zu Monitoring \& Ampel-Reporting (AR-VO)]{VO 17.14a: Verordnung zu Monitoring \& Ampel-Reporting (AR-VO)}
\noindent\textit{ID: 70 -- art-17-14-ar-vo (meta)}

\begin{quote}
Konkretisiert Art. 17.14 (AR): Datenmodell, Schnittstellen, Fristen, Qualitätsstufen, Ampel-Logik, Veröffentlichung/Open-Data, Audit/QA, Governance, Übergang \& Evaluation. Querverweise: Registry/MRV (Hauptgesetz), Verbraucherrechte/Claims (17.5), Durchsetzung/Vollzug (17.4), Haushalt/TFK (17.2), Energie/EN-VO (17.13), IDs/DS (17.1), eIDAS/Signaturen (17.12).
\end{quote}

- Gegenstand, Begriffe \& Identifikatoren
- Einheitliche Datenerhebung/-veröffentlichung zu: Klimapfaden, Investitionskorridor, Rewards/Serien, Risiken/Rücklagen, Haushalts-/Tilgung, Vollzug/Fristen, Energie-KPIs, Sozial-/Arbeits-KPIs.

- Identifikatoren: Projekt-ID, Serien-ID, Anlagen-ID, LEI/EORI, Markt-/Messlokation, Vergabe-/Vertrags-ID, Geo-ID (EPSG/WGS84).

- Assurance-Level: AL-0 (vorläufig), AL-1 (prüfbereit), AL-2 (geprüft), AL-3 (behördlich bestätigt).

- Datenmodell \& Mindestinhalte
- Kernobjekte: Emissionsbudget (Soll/Ist/Varianz), Issuance/Transfer/Redeem, Permanenzklasse/Buffer, EU-Beihilfe-Status, TFK/KKR, Tilgung, Risiko-/Rücklagenstand, Vollzugsfristen, Energiepreis-/Netz-KPIs, TKUG/Weiterbildung (aggregiert).

- Zeitraster: mind. quartalsweise; Energiedaten 15-min/Tag aggregiert; Haushalt/Beihilfe quartalsweise.

\{
  "schema":"AR-VO/Core/1.0",
  "period":"2026-Q1",
  "national\_budget":\{"remaining\_mtCO2e":123.4,"path\_variance\_mtCO2e":-4.2\},
  "corridor":\{"cap\_pct\_gdp":"[VO]","utilization\_pct":61.3,"tilgung\_ok":true\},
  "rewards":\{"issued\_tCO2e":18500000,"redeemed\_tCO2e":6100000,
             "by\_line":\{"A":...,"B":...,"C":...\},"buffer\_tCO2e":420000\},
  "eu\_gate":\{"state\_aid\_cases":\{"pending":7,"approved":21,"rejected":1\}\},
  "risks":\{"guarantees\_outstanding\_eur":2340000000,"loss\_events":2\},
  "vollzug":\{"avg\_case\_days":54,"deadline\_breaches":12,"escalations":3\},
  "energy":\{"avg\_total\_price\_ct\_kWh":23.7,"nne\_tou\_spread\_pct":140,
            "sec14a\_events":12450,"load\_shift\_MW":580\},
  "social":\{"tkug\_participants":48200,"completion\_rate\_pct":72.4\},
  "assurance\_level":"AL-2",
  "qseal":"<registry-seal>"
\}

- Einlieferung, Schnittstellen \& Signaturen
- Abgabe per API (/v1/ar) durch Register-, Haushalts-, Energie-, Sozial- und Zollbehörden; Ereignisdaten werden once-only aus dem Registry gespiegelt (17.4/17.1).

- Signaturen: Datenpakete QSeal; contributor-seitig QES oder organisationsbezogene QSeal-Light [VO].

- Fristen: Quartal bis T+30, Jahresabschluss bis T+60; Korrekturen als Delta mit corr\_id.

- Pflichtfeed /open-data/ops mit Uptime, Incidents, MTTR und Release-Hinweisen (monatlich, QSeal-gesiegelt).

- Ampel-Logik (Trigger \& Maßnahmen)
- Budgetpfad (Art. 1): Grün ≤ +1 \% Abweichung; Gelb > +1 \% bis +3 \%; Rot > +3 \% ⇒ Maßnahmenplan binnen 30 Tagen.

- Korridor-Cap (17.2): Grün Ausnutzung ≤ [VO 70\%]; Gelb ≤ [VO 90\%]; Rot > [VO 90\%] ⇒ Priorisierung/Stop-Go-Prüfung.

- Tilgungspfad: Grün on track; Gelb −1 Quartal; Rot ≥ −2 Quartale ⇒ Nachsteuerung.

- Risiken/Rücklagen: Grün ≥ Mindestniveau; Gelb < Mindestniveau um ≤ [VO 10\%]; Rot < Mindestniveau um > [VO 10\%] ⇒ Nachdotierung/Publikation Loss-Events.

- Vollzug/Fristen (17.4): Grün ≤ [VO 10\%] Fristüberschreitungen; Gelb ≤ [VO 20\%]; Rot > [VO 20\%] ⇒ automatische Eskalation.

- Beihilfe-Trigger: eu\_gate.pending\_days > 180 oder eu\_gate.pending\_share > 0.25 ⇒ Volumenkappung neuer Ausgaben + Priorisierung nach €/tCO\$\_2\$ bis Genehmigung.

- OPS-Trigger: ops.uptime < 99.5\% oder ops.critical\_incidents ≥ 2 ⇒ Feature-Freeze, verpflichtender Fallback.

- Assurance-Trigger: audits.fail\_rate ≥ 2\% zwei Quartale ⇒ Stichprobenquote ↑ und Blacklisting-Prüfung.

- Energie-KPIs (17.13): Rot außerhalb ToU-Korridor ⇒ Anpassung Zonen/Fenster/§14a-Profile.

- Sozial/Qualifizierung (17.7): Grün Abschlussquote ≥ [VO 70\%]; Gelb ≥ [VO 60\%]; Rot < [VO 60\%] ⇒ Kurs-/Matching-Reform.

\{
  "schema":"AR-VO/Ampel/1.0",
  "signals":\{
    "budget\_path":"yellow",
    "corridor\_cap":"green",
    "debt\_service\_path":"green",
    "risk\_reserve":"yellow",
    "enforcement\_sla":"green",
    "energy\_kpis":"yellow",
    "social\_kpis":"green"
  \},
  "triggers":[
    \{"signal":"budget\_path","reason":"path\_variance\_mtCO2e = +2.1"\},
    \{"signal":"energy\_kpis","reason":"nne\_tou\_spread\_pct > upper\_korridor"\}
  ],
  "actions\_due\_in\_days":30
\}

- Veröffentlichung \& Open-Data
- Öffentliches Ampel-/Verfahrensdashboard mit Feeds: /open-data/ampel, /open-data/rewards, /open-data/energy, /open-data/vollzug.

- Lizenz DL-BY 2.0, Versionierung/Änderungsjournal; Personenbezug pseudonymisiert.

- Snapshot-Beweis: QSeal-gesiegelte Monats-/Quartalssnapshots mit Hash-Kette.

- Qualitätssicherung \& Audit
- Mehrstufige Validierung (Syntax/Plausibilität/Cross-Checks gg. ETS/CBAM/Zoll/Haushalt).

- Stichproben; Fehlerklassen: minor 15 Tage, major 30 Tage, critical sofort.

- Rechnungshof/Statistik erhalten Vollzugriff (minimal personenbezogen); jährlicher Prüfbericht.

- Streitbeilegung \& Ombud
- Digitale Ombudsstelle mit Akteneinsicht; Entscheidung binnen 30 Tagen (17.4-Bezug); veröffentlichter Tätigkeitsbericht.

- ADR-Schnittstelle für Markt-/Verbraucherstreitigkeiten (17.5-Bezug).

- Governance \& Zuständigkeiten
- Leitbehörde AR koordiniert; Stabilitätsrat-Spur: quartalsweise Übermittlung Budget/Korridor/Tilgung.

- EU-Beihilfe-Status via EU-Gate konsistent gespiegelt; Ampel-Trigger: „Beihilfe pending > [VO 180] Tage“.

- Übergang, Meilensteine \& Evaluation
- M0--M3: Kern-Open-Data (Rewards/Issuance/Vollzug); M4--M6: Ampel-Dashboard; M7--M12: Energie-/Sozial-KPIs integriert.

- Jährliche Evaluation: Datenlatenz, Fehlerquote, Ampel-Trefferqualität, Nutzungsmetriken, Maßnahmen-Erfolgsquote.


%% ----------------------------------------
\section[Artikel 17.S: Schlussknoten Harmonisierung]{Artikel 17.S: Schlussknoten Harmonisierung}
\noindent\textit{ID: 71 -- art-17-schlussknoten (meta)}

\begin{quote}
Sichert Auslegung, Pflege und Kohärenz dieses Kapitels -- ergänzend zu den Schlussbestimmungen (Art. 18), ohne neue Verfahren zu schaffen.
\end{quote}

- Kohärenz \& Vorrang: Bei Auslegungszweifeln in diesem Kapitel gelten Zielvorrang (Art. 1), Zuständigkeitsklarheit (Art. 8) und Daten-/Once-Only-Grundsätze (Art. 17.14) als Leitregeln. Spezialnormen dieses Kapitels gehen den allgemeinen Harmonisierungsklauseln vor.

- Dynamische EU-Bezüge: Verweisungen auf EU-Recht (z. B. ETS/CBAM/eIDAS/ESPR/CEEAG/GBER) erfolgen auf deren jeweils geltende Fassung; materiell-rechtliche Änderungen werden im Rahmen von Art. 18 redaktionell nachvollzogen.

- Redaktionelle Pflege: Die Bundesregierung veröffentlicht halbjährlich eine konsolidierte Lesefassung der in Art. 17 geänderten/verbundenen Vorschriften mit Änderungsjournal; rein redaktionelle Anpassungen (Nummern/Querverweise) erfolgen per VO nach Art. 18.

- Kollisionslösung: Kollisionen zwischen harmonisierten Fachregimen werden zieläquivalent gelöst (mildestes, wirksames Mittel); Doppelmeldungen und Doppelgutschriften bleiben unzulässig.

- Monitoring-Andockung: Kennzahlen, Fristen und Schnittstellen dieses Kapitels sind über die AR-VO (Art. 17.14) in das öffentliche Ampel-Dashboard einzuspeisen.

{\footnotesize\begin{quote}
Minimaler Abschlussanker für Art. 17: Auslegungshilfe, EU-Dynamik, Pflege -- die materiellen Zeitpunkte/Reviews stehen in Art. 18.
\end{quote}}



%% ----------------------------------------
\section[Artikel 18: Schlussbestimmungen]{Artikel 18: Schlussbestimmungen}
\noindent\textit{ID: 72 -- artikel18 (gesetz)}

\begin{quote}
Dieses Gesetz bündelt die zentralen Instrumente der Transformation -- Preis, Reward, Kredit, Vollzug -- in einem einheitlichen Ordnungsrahmen. Artikel 18 regelt Inkrafttreten, Übergänge, Rechtsbereinigung, Evaluation und digitale Grundsätze -- ohne neue Parallelverfahren.
\end{quote}

- Bezeichnung \& Rang: Dieses Artikelgesetz führt die Kurzbezeichnung Klima-Generationen-Gesetz (KlimaGG); öffentlich kann die Bezeichnung „Klima-Grundgesetz (KlimaGG)“ verwendet werden. Soweit dieses Gesetz speziellere Regelungen trifft, gehen sie abweichendem Bundesrecht vor; Unions- und Völkerrecht bleiben unberührt.

- Inkrafttreten (gestuft):
- Tag 0 (am Tag nach der Verkündung): Verfassungsänderungen (Art. 1), Institutionen (Agentur, Kommission, Gerichtsbarkeit), Serviceplattform \& Beschleunigung (VO/VwV) sowie der Bund-Länder-Pakt treten in Kraft.

- +3 Monate: Transformationskredit und Transformationsförderung gelten für Neuanträge; Standardkataloge/SLAs der Serviceplattform sind anzuwenden.

- +6 Monate: Carbon-Reward startet mit Pilotfenstern nach seinen Übergangsvorschriften; CO\$\_2\$-Steuer (Art. 7) beginnt mit dem bekanntgemachten Preisplan; die strompreis-bezogenen Leitplanken greifen gemäß den dortigen Übergangsregeln.

- +9--12 Monate: Klimadividende (Art. 8), CO\$\_2\$-Grenzausgleich (Art. 9), Kommunaler Klima- \& Resilienzfonds sowie Katastrophenschutz/Vorsorge greifen gemäß den Andock-/Notifizierungsakten nach Art. 17 ff.

- Übergangsrecht \& Bestandsschutz:
- Bestands-Vintages von Carbon-Rewards und bestehende Förder-/Vertragsverhältnisse behalten ihre Rechte; ein Opt-in in KlimaGG-Instrumente ist zulässig, sofern Doppelförderung ausgeschlossen ist.

- Bereits begonnene Verfahren in bestehenden Fachgesetzen (z. B. EEG/KWKG/EnFG/BEHG/TEHG) laufen weiter; Wechsel in KlimaGG-Instrumente ist zulässig nach Maßgabe der Harmonisierung (Art. 17 ff.).

- Laufende Vergaben bleiben im bisherigen Rahmen; neue Vergaben nutzen die LCC/CO\$\_2\$-Kriterien und Safe-Harbor-Muster der Serviceplattform.

- Rechtsbereinigung \& Konsolidierung: Abweichende Bundesnormen werden gemäß dem Harmonisierungskatalog (Art. 17 ff. und Anlagen) geordnet angepasst oder aufgehoben. Die Bundesregierung wird ermächtigt, binnen 6 Monaten eine Erste Rechtsbereinigungs-VO zu erlassen, die ausschließlich redaktionelle Anpassungen (Verweisungen, Nummerierungen, Aufhebungslisten, Konsolidierungen) umsetzt.

- Digital-First \& Once-Only: Verfahren nach diesem Gesetz werden digital geführt (barrierefrei). Bereits erfasste Daten dürfen nicht erneut angefordert werden (Once-Only); Papier-/Offline-Zugänge werden als Härtefallkanal bereitgehalten. Offene, dokumentierte APIs und die Open-Data-Regeln nach Art. 17.14 gelten verbindlich.

- Evaluation \& Review:
- Gesamtevaluation 24 Monate nach Inkrafttreten der Kerninstrumente (CO\$\_2\$-Steuer 7, Dividende 8, CBAM 9, Carbon-Reward), danach alle 48 Monate (Ziele Art. 2/1, Verteilung, Kosten/Nutzen, Vollzug).

- Verordnungsermächtigungen unterliegen einer Review-/Sunset-Pflicht von 36 Monaten, sofern im Einzelartikel nichts Abweichendes geregelt ist.

- Bei ausbleibender Review tritt die jeweilige VO auf eine definierte Baseline zurück (Default-Schwellen \& Methoden gem. Stammgesetz); laufende Einlösungen bleiben unberührt.

- Notifizierung \& Außenbezug: Erforderliche EU-Notifizierungen (Beihilfe/CBAM/Steuern) und völkerrechtliche Anzeigen werden durch die Bundesregierung betrieben; bis zur Genehmigung gelten die jeweiligen Übergangsregeln der Einzelartikel. Art. 17 regelt Verfahren/Fristen und Datenformate.

- Beihilfe-Kohärenz (Übergang): Für beihilferelevante Serien gilt der Redeem-Lock (17.1/17.2) bis zur Genehmigung; Parameter-Switch auf GBER-Standard ist durch Bekanntmachung zulässig.

- Sprach- \& Auslegungsregeln: Begriffe folgen den Definitionen in den Einzelartikeln (insb. Ziele/Methodik/Claims); Re-Basierungen erfolgen ambitionsneutral. Gleichstellungsform gilt.

- Teilnichtigkeit: Ist eine Bestimmung nichtig oder unanwendbar, bleibt das Gesetz im Übrigen wirksam; die nichtige Regel ist zieläquivalent fortzubilden.

- Bekanntmachung \& Wirksamkeit: Die Bundesregierung wird ermächtigt, eine amtliche konsolidierte Fassung dieses Gesetzes mit aktualisierten Verweisungen zu veröffentlichen. Dieses Gesetz tritt im Übrigen am 1. Januar des auf die Verkündung folgenden Jahres in Kraft; entgegenstehende/ersetzte Bundesnormen treten nach Maßgabe des Harmonisierungskatalogs zu den dort genannten Zeitpunkten außer Kraft.

{\footnotesize\begin{quote}
Schlanker Schlussstein: klare Starttreppe ohne Doppelregelungen, Bestandsschutz ohne Doppelförderung, Digital-/Open-Data-Pflicht, regelmäßige Reviews und eine rein redaktionelle erste Rechtsbereinigungs-VO zur sauberen Konsolidierung.
\end{quote}}



\end{document}
